\documentclass[11pt,a4paper]{article}
\usepackage[T1]{fontenc}
\usepackage{lmodern}
\usepackage[margin=27mm]{geometry}
\usepackage{amsmath,amssymb,amsthm,mathtools}
\usepackage{booktabs,array,longtable}
\usepackage{microtype}
\usepackage{xcolor}
\usepackage{hyperref}
\hypersetup{colorlinks=true,linkcolor=blue!45!black,citecolor=blue!45!black,urlcolor=blue!45!black,pdftitle={Spanning-tree leaves, signed moments, and finite-dimensional descent},pdfauthor={Research draft prepared with ChatGPT}}
\newtheorem{theorem}{Theorem}[section]
\newtheorem{lemma}[theorem]{Lemma}
\newtheorem{proposition}[theorem]{Proposition}
\newtheorem{corollary}[theorem]{Corollary}
\newtheorem{conjecture}[theorem]{Conjecture}
\theoremstyle{definition}
\newtheorem{definition}[theorem]{Definition}
\newtheorem{problem}[theorem]{Problem}
\theoremstyle{remark}
\newtheorem{remark}[theorem]{Remark}
\newcommand{\E}{\mathbb E}
\DeclareMathOperator{\Var}{Var}
\DeclareMathOperator{\Cov}{Cov}
\DeclareMathOperator{\Arg}{Arg}
\DeclareMathOperator{\supp}{supp}
\newcommand{\dd}{\,\mathrm d}
\newcommand{\R}{\mathbb R}
\newcommand{\N}{\mathbb N}
\newcommand{\ii}{\mathrm i}
\numberwithin{equation}{section}
\setlength{\emergencystretch}{2em}
\title{Spanning-tree leaves, signed moments,\protect\\and finite-dimensional descent\protect\\[4pt]\large Ten arguments on published mathematical questions}
\author{Research draft prepared with ChatGPT for Keio}
\date{Revised 3 November 2026\protect\\\small Complete arguments for review; historical priority not certified}
\begin{document}
\maketitle
\begin{abstract}
We give complete arguments for review on ten published questions.
For every connected graph with at least two vertices we prove
$L+b\ge 2\alpha+d_2$, the 2005 Graffiti.pc Conjecture~177. Here $L$
is the maximum number of spanning-tree leaves, $b$ the largest order
of an induced bipartite subgraph, and $d_2$ the second-smallest vertex
degree. The proof begins with a deletion lemma for maximal independent sets.
We give interior counterexamples to the global reverse modulus inequality
for Mittag--Leffler functions in Garrappa--Gerhold--Popolizio--Simon
Conjecture~25. For the cubic alternating example of
Bhatnagar--Rajkumar's telescoping continued fraction, negative partial
numerators occur infinitely often, with only $O(\log K)$ among the
first $K$. This refutes their eventual-sign Conjecture~6.1.

The other seven arguments concern elementary stability, coefficientwise
Hankel positivity, two DeTemple approximation problems, diagonal
Pad\'e bounds for the scaled gamma function, Wright-density
log-concavity, and infinite divisibility of negative powers.
The manuscript distinguishes the mathematical arguments from the
evidence for their historical priority. All ten are research-draft
results. Section~\ref{sec:status} identifies the remaining source
comparisons and the extent of independent checking.
\end{abstract}

\section{Scope, original targets, and status}

Each row below corresponds to one previously stated question and a complete
argument offered for mathematical review. A counterexample settles a
universal assertion when it meets the original hypotheses and violates
its conclusion. Two clauses proved by the same construction appear
in one row. The table covers ten targets. The additional requirements of
historical priority, a new principle for each solution, and a comparison
with the entire preceding corpus remain unverified.

\begin{center}
\small
\begin{tabular}{@{}>{\raggedright\arraybackslash}p{.31\textwidth}>{\raggedright\arraybackslash}p{.54\textwidth}r@{}}
\toprule
Published target & Result & Page\\
\midrule
Chen--Kauers, Problem 8 & Elementary stability implies $D$-finiteness and admits a finite adjoint criterion. & \pageref{stab:section}\\[4pt]
Sokal, Conjecture 6.9(b),(d) & Coefficientwise Hankel total positivity for functional digraphs and a positive measure for connected counts. & \pageref{sec:tp}\\[4pt]
Yang, Conjecture 1(i) & Strict monotonicity of the DeTemple quotient with both sharp endpoint limits. & \pageref{sec:yang}\\[4pt]
Yang--Chu--Zhang, Problem 1 & Coefficients, strict complete monotonicity, and the prescribed leading asymptotic at every order. & \pageref{sec:yang-allorders}\\[4pt]
Li--Chen, Conjecture 4.1 & Strict alternating bounds for every original diagonal Pad\'e approximant on $x>1/24$. & \pageref{sec:li-chen-pade}\\[4pt]
Ferreira--Simon, Section 6.4 & An interior rational counterexample to the Wright-density threshold. & \pageref{sec:wright-counterexample}\\[4pt]
Yano--Yano--Yor, Remark 2.9(iii) & A negative power of a half-$\alpha$-Cauchy variable fails infinite divisibility. & \pageref{sec:yyy}\\[4pt]
Graffiti.pc, WoWII 177 & $L+b\ge2\alpha+d_2$ for every connected graph of order at least two. & \pageref{sec:wow177}\\[4pt]
Garrappa et al., Conjecture 25 & Interior counterexamples to the reverse modulus inequality for each $1<\alpha<2$. & \pageref{ggps:section}\\[4pt]
Bhatnagar--Rajkumar (2026), Conjecture 6.1 & Infinitely many negative and positive partial numerators in Appendix A.3, with positive signs of density one. & \pageref{sec:br-sign}\\
\bottomrule
\end{tabular}
\end{center}

Three auxiliary questions clarify the contributions of this revision.
Can an independent set supply a vertex whose neighbors, except for one,
can be deleted while keeping the graph connected? What happens to a
modulus inequality when its leading asymptotic coefficient vanishes at
a parameter boundary? How fast can the negative index of a signed moment
form grow while its continued-fraction coefficients remain largely
positive? Sections~\ref{sec:wow177}, \ref{ggps:section}, and
\ref{sec:br-sign} answer these questions and apply the answers to the
three published targets in the last rows.

The other principal constructions are a differential-field descent,
a coefficientwise Bell--cumulant criterion, and a matrix marginal
representation for an orthogonal-polynomial transform. Here total
positivity follows the convention of the sources that every minor is
nonnegative. The arguments identify the classical inputs they use.

The \href{https://github.com/QEEDDD/Verifiable-Mathematical-Reasoning-Corpus}{Verifiable Mathematical Reasoning Corpus}
was inspected at
\href{https://github.com/QEEDDD/Verifiable-Mathematical-Reasoning-Corpus/commit/ecd87de3348dcb706a76bb12fde722d7a348a012}{\texttt{ecd87de3348dc}}.
Its earlier work on restricted WoWII cases, signed kernels, and boundary
limits informed the choice of questions. Sections below supply their own
mathematical arguments. The previously prepared Le Roy and high-dimensional
rook-recurrence results are separate from the ten entries.

\tableofcontents
\bigskip


\section{Elementary stability and a finite adjoint criterion}
\label{stab:section}

\paragraph{Status and scope.}
This section presents a proposed complete proof for review of a structural
classification answering Chen--Kauers Problem~8.  It is not a claim of
independently certified historical priority.  The problem asks for the class
of elementary functions whose every repeated indefinite integral is
elementary \cite[Problem~8]{stab:CK}.  We prove that such functions must be
$D$-finite, and then give a finite elementary-integrability criterion once
an annihilator is known.  The proof allows arbitrary finite algebraic
adjunctions interspersed with logarithmic and exponential adjunctions.

Throughout, $C$ is an algebraically closed field of characteristic zero,
$B=C(x)$, and $D x=1$.  All differential fields have constants $C$.
For ordinary complex elementary functions one takes $C=\mathbb C$ and
works with compatible germs in a common differential extension; choices
of integration constants do not affect stability.  A function is
\emph{elementarily stable} if it has an elementary $n$-th primitive for
every $n\geq 1$.

Chen established the moment formulation and several special families of
stable functions \cite{stab:Chen}.  Sahu, Seiss and Srinivasan gave a
generalized Liouville theorem for an elementary primitive of any fixed
depth \cite{stab:SSS}.  The complete reduction of Chen et al. treats
single-step integration in transcendental Liouvillian towers
\cite{stab:CDGHLL}.  The argument here uses the classical Liouville theorem
\cite{stab:Rosenlicht} and Rosenlicht's local function-field estimates
\cite{stab:Rosenlicht1969}, together with a uniform relative bound:
elementary primitives of integrands in a fixed finite-dimensional
$B$-space lie in one finite-dimensional $B$-space modulo a lower field.
A finite-dimensional differential-module argument then removes one
transcendental layer. The same bound gives a finite/cofinite dichotomy
for integrable polynomial weights. Chen's separate moment-set
conjecture remains unresolved in its $D$-finite case.

\subsection{Moments and a finite primitive-space lemma}

\begin{lemma}[Moment criterion]\label{stab:moments}
An elementary function $f$ is stable if and only if $x^n f$ has an
elementary primitive for every $n\geq0$.  The stable functions form a
$C[x][D]$-module.
\end{lemma}
\begin{proof}
Repeated integration by parts expresses a primitive of $x^n f$ in terms
of the first $n+1$ primitives of $f$.  Conversely, if $Dg_k=x^k f$ for
$0\leq k<n$, then
\[
 \frac1{(n-1)!}\sum_{k=0}^{n-1}(-1)^k
       \binom{n-1}{k}x^{n-1-k}g_k
\]
is an $n$-th primitive of $f$.  Addition and multiplication by $C[x]$
preserve the moment condition.  Differentiation preserves it because
$x^n Df=D(x^n f)-n x^{n-1}f$, with the second term absent for $n=0$.
\end{proof}

\begin{lemma}[Finite primitive space for independent weights]
\label{stab:finite-space}
Let $M$ be a differential $B$-module.  Suppose that
$f,g_0,\ldots,g_d\in M$, that the $g_i$ span a $B$-space of
dimension at most $d$, and that
\[
 Dg_i=p_i f\qquad(0\leq i\leq d),
\]
where $p_0,\ldots,p_d\in B$ are linearly independent over $C$.
Then $f$ belongs to $V=\operatorname{span}_B(g_0,\ldots,g_d)$,
$DV\subseteq V$, and $f$ is annihilated by a nonzero operator in
$B[D]$.
\end{lemma}
\begin{proof}
Suppose $f\notin V$.  Choose a nonzero relation
$\sum_{i\in J}a_i g_i=0$ with $J$ of minimal cardinality, and
normalize one coefficient to $1$.  Differentiating gives
\[
 \sum_{i\in J}a_i'g_i+
 \left(\sum_{i\in J}a_i p_i\right)f=0.
\]
Since $f\notin V$, both terms vanish separately.  The first relation
has smaller support unless every $a_i'=0$, so minimality forces all
$a_i\in C$.  The second relation then contradicts the $C$-linear
independence of the $p_i$.  Hence $f\in V$.  The equations
$Dg_i=p_if$ imply $DV\subseteq V$, and finite dimension gives a
linear dependence among successive derivatives of $f$.
\end{proof}


\subsection{One algebraic function-field layer}

We use strong induction on the number of transcendental generators in a
finite Liouvillian tower over $B$.  Arbitrary finite algebraic
adjunctions are allowed and are included in the induction hypothesis.
At a positive level write
\begin{equation}\label{stab:layer}
 E/F(t)/F/B,
\end{equation}
where $E/F(t)$ is finite algebraic, $F$ has one fewer transcendental
generator, $t$ is transcendental, and either
\[
 Dt=\beta\in F\quad\hbox{or}\quad Dt=\alpha t,\quad\alpha\in F.
\]
We call these the primitive and exponential cases.  Replacing $F$ by
its relative algebraic closure in $E$ is a finite algebraic enlargement
and does not increase the induction level.  We do so, and view $E$ as
the function field of a smooth projective geometrically integral curve
$X/F$, finite over the $t$-line.

\begin{lemma}[Properness after algebraic enlargement]\label{stab:proper}
In the primitive case, $\beta$ is not a derivative in any algebraic
extension of $F$.  In the exponential case, no nonzero rational multiple
of $\alpha$ is a logarithmic derivative of an element algebraic over $F$.
\end{lemma}
\begin{proof}
If $Du=\beta$ in a finite algebraic extension $F_1/F$, then normalized
trace gives a primitive in $F$.  Subtracting it from $t$ gives a new
constant, contrary to the hypotheses.  If $Du/u=q\alpha$ with
$q=a/b\in\mathbb Q\setminus\{0\}$ and $u$ algebraic over $F$, then
$D(u^b/t^a)=0$.  The field containing this quotient is algebraic over
$F(t)$, so its constants are still $C$.  This would make a nonzero power
of $t$ algebraic over $F$, a contradiction.  Here we use the standard
fact that constants in an algebraic differential extension are
algebraic over the original constant field; $C$ is algebraically closed.
\end{proof}

An irreducible $p\in F[t]$ of positive degree is called normal when
$p\nmid Dp$.  Every finite irreducible is normal in the primitive case.
Indeed, for monic $p$ of degree $d$, divisibility would imply $Dp=0$;
the coefficient of $t^{d-1}$ would give $Db=-d\beta$ for the next
coefficient $b$, contradicting Lemma~\ref{stab:proper}.  In the
exponential case the same argument shows that every irreducible other
than $t$ is normal: a monic special polynomial of degree $d$ with
nonzero constant term $b_0$ would satisfy $Db_0=d\alpha b_0$.

The local valuation arguments are classical. Rosenlicht proved continuity
of the derivation at function-field places and corresponding ramified-place
estimates for derivatives and logarithmic derivatives in primitive and
exponential extensions \cite[Lemmas~1--2]{stab:Rosenlicht1969}.
We record the bounds needed here, together with the fixed finite
coefficient extension used in the subsequent uniform argument.

\begin{lemma}[Local bounds and Puiseux coordinates]\label{stab:local}
Let $P$ lie over a finite normal divisor $p$, with ramification index
$e$, and normalize its valuation $v_P$ to have value group $\mathbb Z$.
Then
\begin{align}
 v_P(Dh)&=v_P(h)-e &&\text{if }v_P(h)<0,\label{stab:negative-local}\\
 v_P(Dh)&\geq1-e &&\text{if }h\text{ is a local unit},\label{stab:unit-local}\\
 v_P(Du/u)&\geq-e &&\text{if }u\in E^*.\label{stab:log-local}
\end{align}
At a selected place over infinity, a fixed finite algebraic extension
$F'/F$ gives an embedding of the base-changed function field into
$F'((t^{-1/e}))$ that commutes with $D$.  Analogous local expansions
hold at all the other places under consideration.
\end{lemma}
\begin{proof}
For the finite-place assertions, extend algebraic coefficients and
choose a root $c$ of $p$.  Since characteristic is zero,
$p_t(c)\ne0$, and differentiation of $p(c)=0$ gives
\[
 (Dp)(c)=p_t(c)\bigl((Dt)|_{t=c}-Dc\bigr).
\]
Thus $\gamma=(Dt)|_{t=c}-Dc$ is nonzero.  In a Puiseux coordinate
$z$ with $t-c=z^e$,
\[
 Dz=\frac{\gamma}{e}z^{1-e}+\text{terms of degree at least }1.
\]
For $h=\sum_{k\geq r}h_kz^k$, $h_r\ne0$, and $r<0$, the leading
term of $Dh$ is $(r\gamma/e)h_rz^{r-e}$.  Derivatives of coefficients
have exponent at least $r$, and derivatives of higher powers have
exponent at least $r+1-e$.  This proves
\eqref{stab:negative-local}.  For a unit, differentiating its constant
coefficient is regular and every positive power contributes exponent
at least $1-e$, proving \eqref{stab:unit-local}.  Factoring $u$ into a
power of $z$ times a unit gives \eqref{stab:log-local}.

For the finite coefficient extension, first extend $F$ by the finite
residue field of a selected point over infinity.  The corresponding
point after smooth base change is rational.  Its completed local ring
is $F_1[[z]]$, with the given $F_1$ as coefficient field.  Write
$t^{-1}=z^e u(z)$, where $u(0)\ne0$.  Adjoin an $e$-th root of
$u(0)$, a finite extension of $F_1$.  Hensel lifting, or the formal
binomial expansion in characteristic zero, supplies an $e$-th root
$U(z)$ of $u(z)$ over the resulting finite field $F'$.  Then
$w=zU(z)$ is a uniformizer with $w^e=t^{-1}$, and the completed field
is $F'((w))$.  The coefficientwise formal derivation having the
specified value on $t$ agrees on the algebraic function field with the
unique extension of $D$.  The finite-place constructions are the same
after replacing $t^{-1}$ by a local base coordinate.  For local bounds
one may work in an algebraic closure of $F$; only one selected place
needs the fixed finite coefficient extension for the global argument.
\end{proof}

For an expansion at infinity, write $\deg_P h=d$ when its leading
term is $h_dt^d$, with $d\in e^{-1}\mathbb Z$.  A rational point over
infinity also supplies a useful basis construction.  For any fixed
divisor $H$, the Riemann--Roch space $\mathcal L(H)$ is finite-dimensional
over $F'$.  Its valuation filtration at that point has graded pieces
of dimension at most one.  It therefore has a basis $w_1,\ldots,w_s$
with distinct leading degrees $d_1>\cdots>d_s$, each leading
coefficient normalized to $1$.  In the primitive case
$\deg_P Dw_j<d_j$; in the exponential case the degree-$d_j$ coefficient
of $Dw_j$ is $d_j\alpha$.  If $1\in\mathcal L(H)$, the basis can
include $1$ as its unique degree-zero element: subtract the leading
constant from every other element initially having degree zero.

\subsection{Uniform first-order preimages}

\begin{lemma}[Uniform preimage space]\label{stab:uniform}
Let $K/B$ be a finite Liouvillian tower with no new constants, allowing
arbitrary finite algebraic adjunctions.  For $a\in K$ and a
finite-dimensional $B$-space $V\subseteq K$, the set
\[
 \{g\in K: Dg+ag\in V\}
\]
is contained in a finite-dimensional $B$-subspace of $K$.
\end{lemma}
\begin{proof}
Use strong induction on the transcendental length.  At length zero
$K/B$ is finite algebraic, and the assertion is immediate.  At positive
length use \eqref{stab:layer}, with $K=E$.

We first bound the pole divisor of every solution $g$.  At a finite
normal place $P$, let $M$ bound the pole orders of generators of $V$,
and suppose $q=v_P(g)<0$.  If $v_P(a)<-e$, multiplication by $a$
dominates differentiation, so $q+v_P(a)\geq-M$.  If $v_P(a)>-e$,
differentiation dominates, so $q-e\geq-M$.  If $v_P(a)=-e$, the local
leading coefficient is
\[
 h_q\bigl(q\gamma/e+a_{-e}\bigr).
\]
It can vanish for at most one integer $q$.  This gives a finite bound
in every case.  Outside the finite set consisting of poles of $a$,
poles of generators of $V$, and places over branch points, the bound
forces $g$ to be regular.

Next consider every place above infinity.  Set $A=\deg_P a$, and let
$N$ bound the degrees of generators of $V$.  When $A>0$, multiplication
dominates because differentiation does not increase Puiseux degree.
Consequently $d=\deg_P g\leq N-A$.

In the exponential case with $A\leq0$, a solution with $d>N$ must
have leading coefficient satisfying
\[
 \frac{Dh_d}{h_d}+d\alpha+a_0=0.
\]
There is at most one rational $d$ for which this is possible with
$h_d$ algebraic over $F$: two choices contradict
Lemma~\ref{stab:proper} after dividing their leading coefficients.
Thus $d$ is bounded.  At each place over zero the same argument uses
the lowest Puiseux exponent.  A pole of $a$ gives a direct lower bound
on that exponent; otherwise there is at most one exceptional exponent
below those occurring in $V$.

In the primitive case with $A\leq0$, suppose $d>N+1$.  Divide the
equation by $g$ and write, purely formally,
\[
 g=h_dt^d\exp(\Lambda),\qquad
 \Lambda=\log\bigl(1+\text{strictly lower Puiseux powers}\bigr).
\]
The right side of $Dg/g+a=v/g$ has degree below $-1$.  If
$\Lambda_1$ is the coefficient of $t^{-1}$ in $\Lambda$, comparison
at $t^{-1}$ gives
\begin{equation}\label{stab:primitive-resonance}
 d\beta+D\Lambda_1+a_{-1}=0.
\end{equation}
Indeed, $\beta\partial_t\Lambda$ has degree strictly below $-1$,
even when fractional exponents occur.  Two distinct rational values
of $d$ satisfying \eqref{stab:primitive-resonance} would give a
primitive of $\beta$ algebraic over $F$, excluded by
Lemma~\ref{stab:proper}.  There is at most one exceptional degree.
This proves a bound at every infinity place, not merely at the place
chosen later for coefficient extraction.

The local bounds yield a fixed divisor $H$ containing every pole
divisor of $g$.  Make one infinity place rational over a fixed finite
extension $F'/F$ and choose the normalized echelon basis of
$\mathcal L(H)$ described above.  Write $g=\sum h_jw_j$ with
$h_j\in F'$.  If $a$ has positive degree at that place, the highest
coefficient of $ag$ determines $h_1$ algebraically from a coefficient
of $v\in V$, so $h_1$ lies in a fixed finite $B$-space.  Otherwise
the highest coefficient equation has one of the forms
\[
 Dh_1+a_0h_1\in V_1,
 \qquad
 Dh_1+(a_0+d_1\alpha)h_1\in V_1,
\]
with $V_1\subset F'$ a fixed finite $B$-space.  The induction hypothesis
applies because the finite algebraic enlargement $F'$ has smaller
transcendental length.

Subtract $h_1w_1$ and descend through the finite basis.  Each new
right side remains in a fixed finite $B$-space: the derivative of a
finite $B$-space is contained in a finite $B$-space, since $DB\subseteq
B$.  Induction bounds every $h_j$, and hence $g$, in a finite $B$-space
inside $EF'$.  Intersecting that space with $E$ proves the assertion.
The bounds used here are existential; their effectivity is not needed.
\end{proof}

\subsection{Moving logarithms and descent modulo the lower field}

\begin{lemma}[Relative elementary primitive space]\label{stab:relative}
In \eqref{stab:layer}, let $V\subset E$ be a finite-dimensional
$B$-space.  There exist a fixed finite algebraic extension $F'/F$, an
elementary extension $A$ of $EF'$, and a fixed finite-dimensional
$B$-space $W\subset A$ such that every elementarily integrable $v\in V$
admits
\[
 v=Dw+q,\qquad w\in W,\quad q\in F'.
\]
The fields and the space are independent of $v$.
\end{lemma}
\begin{proof}
Liouville's theorem gives
\begin{equation}\label{stab:liouville}
 v=Dg+\sum_j c_j\frac{Du_j}{u_j},
 \qquad g\in E,\quad u_j\in E^*,\quad c_j\in C.
\end{equation}
The version with all terms in $E$ follows from the usual algebraic
version by trace for the rational part and norm for logarithmic
derivatives.

Let $S$ contain the poles of generators of $V$, all places over branch
points, all infinity places, and all places above zero in the
exponential case.  At $P\notin S$, the derivation is regular and
transverse.  A pole of $g$ would give a pole of $Dg$ of order at least
two, whereas logarithmic derivatives have pole order at most one and
$v$ is regular.  Hence $g$ is regular there.  The simple-pole
coefficient then forces
\[
 \sum_j c_j\operatorname{ord}_P(u_j)=0.
\]
Thus the formal divisor combination $\sum_jc_j\operatorname{div}(u_j)$
is supported on the fixed finite set $S$.

These combinations form a $C$-subspace of $C^S$.  Choose a basis of
its image and one fixed logarithmic representative for each basis
element, writing
\[
 \ell_i=\sum_k c_{ik}\log u_{ik},\qquad \delta_i=D\ell_i\in E.
\]
If a $C$-linear logarithmic combination has zero divisor combination,
choose a $\mathbb Q$-basis of the finite $\mathbb Q$-span of its
coefficients and clear denominators.  Each resulting integer product
of the $u_j$ has divisor zero.  Since $X$ is projective and $F$ is
relatively algebraically closed in $E$, that product belongs to
$F^*$.  Its logarithmic derivative belongs to $F$.  Therefore
\eqref{stab:liouville} can be written
\begin{equation}\label{stab:relative-equation}
 v=Dg+\sum_i\lambda_i\delta_i+q_0,
 \qquad \lambda_i\in C,\quad q_0\in F.
\end{equation}
This controls the moving logarithms modulo the lower field, rather
than assuming their arguments were fixed initially.

We next bound the pole divisor of the original rational part $g$.
At a finite normal place of ramification $e$, a pole of $g$ of order
$m$ produces a pole of $Dg$ of order $m+e$, whereas all logarithmic
derivatives have pole order at most $e$.  The fixed pole bound for
$v$ bounds $m$.  At exponential zero and infinity, logarithmic
derivatives have respectively nonnegative lowest exponent and
nonpositive degree.  A nonzero leading exponent of $g$ cannot
disappear on differentiation, by Lemma~\ref{stab:proper}.  This bounds
its poles there.  At primitive infinity, logarithmic derivatives
have degree at most zero.  If $N$ is the maximum of zero and the
degrees of generators of $V$, then $Dg$ has degree at most $N$.
A degree $d>N+1$ would force $d\beta+D\Lambda_1=0$ by the
$t^{-1}$ coefficient of $Dg/g$, contradicting
Lemma~\ref{stab:proper}.  These estimates apply at every special
place.  Together with regularity outside $S$, they put every $g$ in
one fixed Riemann--Roch space.

Make a selected infinity place rational over fixed finite $F'/F$.
Enlarge the bounded-divisor space to include $1$, and choose an
echelon basis with normalized leading coefficients, distinct degrees,
and $1$ as its unique degree-zero element.  Descend through this
basis in \eqref{stab:relative-equation}.  At every nonzero degree,
$q_0\in F'$ contributes no coefficient.  The leading coefficient
$h_j$ of $g$ satisfies a first-order equation over $F'$ with coefficient
$0$ in the primitive case or $d_j\alpha$ in the exponential case,
and forcing in a fixed finite $B$-space.  Lemma~\ref{stab:uniform}
at the smaller transcendental length confines $h_j$ to a finite
$B$-space.  At degree zero subtract its unrestricted coefficient
$h_0$ and absorb $Dh_0$ into $q_0$.  Continue with all negative
degrees; the lower-field term still contributes no coefficient there.
Consequently
\[
 g=w_*+h,\qquad h\in F',
\]
where $w_*$ ranges in a fixed finite $B$-space.  Adjoin the finitely
many logarithms occurring in the fixed $\ell_i$ and set
$w=w_*+\sum_i\lambda_i\ell_i$.  Then $v=Dw+q$ for some $q\in F'$,
and all such $w$ lie in one fixed finite-dimensional $B$-space.
\end{proof}

\begin{theorem}[D-finiteness of elementary-stable functions]
\label{stab:dfinite}
Every elementary-stable function over $C(x)$ is $D$-finite over $C(x)$.
\end{theorem}
\begin{proof}
Induct on the transcendental length of an elementary field containing
$f$.  At length zero the function is algebraic and hence $D$-finite.
At positive length use \eqref{stab:layer} and apply
Lemma~\ref{stab:relative} to $V=Bf$.  By Lemma~\ref{stab:moments},
every $x^nf$ is elementarily integrable.  There are $w_n$ in one fixed
finite $B$-space such that
\[
 Dw_n=x^nf-q_n,\qquad q_n\in F'.
\]
The quotient $A/F'$ is a differential $B$-module, since $F'$ contains
$B$ and is stable under $D$.  Lemma~\ref{stab:finite-space}, applied
in this quotient to finitely many of the $w_n$, gives a nonzero
$L\in B[D]$ with $Lf\in F'$.  Clear denominators and take
$L\in C[x][D]$.  This step is necessary because stability is not
being asserted to be preserved by arbitrary rational multiplication.

Lemma~\ref{stab:moments} makes $Lf$ elementary-stable.  The field
$F'$ has smaller transcendental length, with only a finite algebraic
enlargement, so induction makes $Lf$ $D$-finite.  A nonzero polynomial-
coefficient annihilator $M$ of $Lf$ gives the nonzero annihilator
$ML$ of $f$.
\end{proof}

\subsection{A finite criterion for the original classification problem}

Let $L=\sum_{j=0}^r a_j(x)D^j\ne0$ annihilate $f$, with
$a_j\in C[x]$.  Expand its formal adjoint with coefficients on the
left:
\[
 P=L^*=\sum_{j=0}^r(-D)^ja_j(x)=\sum_{i=0}^r b_i(x)D^i.
\]
Over nonzero $b_i$, define
\begin{equation}\label{stab:cutoff-data}
 s=\max_i(\deg b_i-i),\qquad
 \rho(z)=\sum_{\deg b_i-i=s}\operatorname{lc}(b_i)(z)_i,
\end{equation}
where $(z)_i=z(z-1)\cdots(z-i+1)$.  The polynomial $\rho$ is nonzero,
because the represented falling factorials have distinct degrees.
Choose $N\geq0$ with $N+s\geq0$ and $\rho(n)\ne0$ for every integer
$n\geq N$, and put $m=N+s$.

\begin{proposition}[Finite adjoint-moment criterion]\label{stab:adjoint}
For an elementary $f$ annihilated by $L$, stability is equivalent to
elementary integrability of
\[
 f,\;xf,\;\ldots,\;x^{m-1}f.
\]
When $m=0$, the test list is empty.
\end{proposition}
\begin{proof}
For $n\geq N$, the leading term of $P(x^n)$ is
$\rho(n)x^{n+s}$.  Triangular elimination by degree gives
\begin{equation}\label{stab:adjoint-span}
 C[x]=P(C[x])+\operatorname{span}_C\{1,x,\ldots,x^{m-1}\}.
\end{equation}
Only a vector-space sum is used; $P(C[x])$ is not asserted to be an
ideal.  The Lagrange identity is
\[
 qL(f)-fL^*(q)=D J_L(f,q),
\]
with the explicit concomitant
\[
 J_L(f,q)=\sum_{j=1}^r\sum_{k=0}^{j-1}
       (-1)^kD^k(a_jq)D^{j-1-k}f.
\]
Since $Lf=0$, $fP(q)$ is the derivative of the elementary function
$-J_L(f,q)$.  Equation~\eqref{stab:adjoint-span} therefore reduces
elementary integrability of every polynomial multiple of $f$ to the
displayed finite list.  Lemma~\ref{stab:moments} completes the proof.
\end{proof}

\begin{theorem}[Classification by an annihilator and finite tests]
\label{stab:classification}
An elementary function $f$ is stable if and only if it is $D$-finite
and, for one nonzero polynomial-coefficient annihilator, the finite
tests of Proposition~\ref{stab:adjoint} pass.
\end{theorem}
\begin{proof}
Necessity of $D$-finiteness is Theorem~\ref{stab:dfinite}; the remaining
equivalence is Proposition~\ref{stab:adjoint}.
\end{proof}

\begin{remark}[One elementary primitive is not enough]\label{stab:example}
Take $f=2x\exp(x^2)$.  It is $D$-finite and has the elementary primitive
$\exp(x^2)$.  An annihilator and its adjoint are
\[
 L=xD-(2x^2+1),\qquad L^*=-xD-2x^2-2.
\]
Thus $s=2$, $\rho=-2$, $N=0$, and the criterion tests both $f$ and
$xf$.  The first passes, while the second fails because
\[
 \int 2x^2e^{x^2}\,dx=xe^{x^2}-\int e^{x^2}\,dx,
\]
and the last integral is not elementary.  The finite criterion therefore
distinguishes stability from $D$-finiteness plus one elementary integral.
\end{remark}

\begin{remark}[Effective decision and constant-field assumptions]
\label{stab:algorithm}
Assume a proper elementary-field presentation with effective arithmetic,
zero tests, the requisite algebraic representations, and an elementary-
integrability decision procedure.  Also assume that a bound beyond all
nonnegative integer zeros of the polynomial $\rho$ can be computed;
this is automatic for effectively represented algebraic complex
coefficients, and follows from any certified complex-root bound.
These are effective
input assumptions of symbolic integration; no algorithmic claim is
made for arbitrary unrepresented complex constants.

Dovetail two searches.  The first tests $f,xf,x^2f,\ldots$ and returns
``not stable'' at the first nonintegrable moment.  The second computes
$f,Df,D^2f,\ldots$ and tests $B$-linear dependence.  Once an
annihilator is found, apply the finite tests above and return their
answer.  A nonstable input has a failed moment, whereas a stable input
is $D$-finite by Theorem~\ref{stab:dfinite}; hence one branch terminates.

Finite $B$-linear dependence is decidable from the proper presentation:
the transcendental generators are algebraically independent over $B$,
and the full field is finite algebraic over their rational function
field.  Expand in an algebraic basis, clear denominators in the
transcendental generators, and equate the finitely many monomial
coefficients.  This gives finite linear algebra over $B$.  Thus an
independent total test for $D$-finiteness is not assumed.  Nor does the
algorithm require effective versions of the existential bounds in
Lemma~\ref{stab:uniform}; those bounds guarantee its termination.
The elementary-integrability component is the usual finite-terms
problem \cite{stab:Risch}; it is kept explicit as an input assumption
rather than conflated with the new stability decision.
\end{remark}

% Insert after Remark stab:algorithm and before the Audit boundary paragraph.
% Requires the generalized replacement of Lemma stab:finite-space supplied
% in finite_primitive_replacement.tex.  Reuses the existing Section 2 results.
% This strengthening is not an additional resolution of a named problem.

\subsection{Polynomial weights and the remaining moment conjecture}
\label{stab:weights-subsection}

The relative primitive-space bound also controls sparse sets of
integrable moments.  For an elementary function $f$ over $C(x)$, put
\[
 I_f=\{p\in C[x]:pf\text{ has an elementary primitive}\}.
\]
This is a $C$-vector space; it need not be an ideal of $C[x]$.

\begin{theorem}[Polynomial-weight dichotomy]
\label{stab:weight-dichotomy}
For every elementary function $f$ over $C(x)$, exactly one of the
following alternatives holds:
\begin{enumerate}
\item $f$ is $D$-finite and $I_f$ has finite codimension in $C[x]$;
\item $f$ is not $D$-finite and $I_f$ is finite-dimensional over $C$.
\end{enumerate}
In particular, infinitely many elementarily integrable monomial moments
$x^n f$ force $f$ to be $D$-finite.
\end{theorem}
\begin{proof}
The case $f=0$ is immediate, with $I_f=C[x]$.
The degree argument giving \eqref{stab:adjoint-span} applies to every
nonzero operator $P\in C[x][D]$: its image $P(C[x])$ has finite
codimension in $C[x]$.  If $Lf=0$ for some nonzero
$L\in C[x][D]$, the Lagrange identity gives
$L^*(C[x])\subseteq I_f$.  Thus $D$-finiteness implies the first
alternative.

It remains to prove that infinite dimension of $I_f$ forces
$D$-finiteness.  Induct on the transcendental length of an elementary
tower containing $f$.  At length zero, $f$ is algebraic.  At a positive
layer \eqref{stab:layer}, apply Lemma~\ref{stab:relative} to $Bf$.
Its conclusion supplies a fixed finite algebraic extension $F'/F$,
an elementary extension $A$ of $EF'$, and a fixed finite-dimensional
$B$-space $W\subset A$ such that, for every $p\in I_f$,
\[
 pf=Dw_p+q_p,\qquad w_p\in W,\quad q_p\in F'.
\]
Let $d$ be the dimension of the image of $W$ in the differential
$B$-module $A/F'$.  Choose $d+1$ polynomials in $I_f$ that are
linearly independent over $C$.  Lemma~\ref{stab:finite-space}, applied
to their corresponding primitive classes, gives a nonzero
$L\in B[D]$ with $Lf\in F'$; if $f\in F'$ already, one can simply
take $L=1$.  Multiplying $L$ on the left by a common denominator
preserves this conclusion and permits $L\in C[x][D]$.
Put $h=Lf$ and $P=L^*$.

For every $q\in C[x]$, the Lagrange identity expresses $qh-fP(q)$
as the derivative of an elementary function.  Therefore
\begin{equation}\label{stab:weight-descent}
 I_h=P^{-1}(I_f).
\end{equation}
The intersection $I_f\cap P(C[x])$ is infinite-dimensional, because
$P(C[x])$ has finite codimension.  Its inverse image under $P$ is
infinite-dimensional as well.  Equation~\eqref{stab:weight-descent}
therefore makes $I_h$ infinite-dimensional.  The field $F'$ has smaller
transcendental length, so induction makes $h$ $D$-finite.  Composing a
nonzero annihilator of $h$ with $L$ gives a nonzero annihilator of $f$.

This proves the dichotomy.  Finally, distinct monomials are linearly
independent over $C$, so an infinite set of integrable monomial moments
makes $I_f$ infinite-dimensional.
\end{proof}

Chen's final conjecture in \cite[Section~5]{stab:Chen} asks whether,
for every elementary $f$, the unchanged set
\[
 S_f=\{n\geq0:x^nf\text{ is elementary integrable over }C(x)\}
\]
is a Skolem set.  We use the usual convention: a finite union of
infinite arithmetic progressions together with a finite set.  Theorem~
\ref{stab:weight-dichotomy} proves the assertion for non-$D$-finite
$f$, since $S_f$ is then finite.  The other case has the following
precise reduction.

\begin{proposition}[Holonomic obstruction sequences]
\label{stab:moment-obstructions}
If $f$ is elementary and $D$-finite, then $S_f$ is the simultaneous
zero set of finitely many $P$-recursive sequences over $C$.
Consequently, a Skolem theorem for all $P$-recursive sequences over $C$
would imply Chen's full elementary-moment conjecture.
\end{proposition}
\begin{proof}
The quotient $Q_f=C[x]/I_f$ is finite-dimensional by
Theorem~\ref{stab:weight-dichotomy}.  Write $e_n=[x^n]\in Q_f$,
choose a nonzero annihilator $L\in C[x][D]$ of $f$, and put
$P=L^*=\sum_i b_i(x)D^i$.  Since $P(C[x])\subseteq I_f$,
expansion of $P(x^n)$ gives, for all sufficiently large $n$,
\begin{equation}\label{stab:obstruction-recurrence}
 \sum_j c_j(n)e_{n+j}=0,
 \qquad
 c_j(z)=\sum_i [x^{i+j}]b_i(x)\,(z)_i.
\end{equation}
Only finitely many $j$ occur; coefficients with negative polynomial
indices are understood to be zero.  The polynomials $c_j$ are not
all zero, since the falling factorials $(z)_i$ are linearly independent
and $P\ne0$.  Thus each coordinate sequence of $e_n$ in a fixed
$C$-basis of $Q_f$ is $P$-recursive.  These coordinates vanish
simultaneously exactly when $x^n\in I_f$.  If $Q_f=0$, take the
single identically zero sequence.

Finite intersections of Skolem sets are Skolem sets.  The conditional
conclusion follows, with the non-$D$-finite case supplied by
Theorem~\ref{stab:weight-dichotomy}.
\end{proof}

The remaining zero-set step is substantial.  The finite-dimensional
quotient $Q_f$ does not in general carry multiplication by $x$, so
\eqref{stab:obstruction-recurrence} need not have constant coefficients.
For example, take $f=e^{x^2}$ and $P=-D-2x$.  Then $e_0\ne0$,
$e_1=0$, and
\[
 e_{n+1}=-\frac n2 e_{n-1}\quad(n\geq1).
\]
In particular, $x\in I_f$ but $x^2\notin I_f$.  The
Skolem--Mahler--Lech theorem for constant-coefficient recurrences
therefore cannot be invoked solely from finite dimension of $Q_f$.
Rubel's holonomic zero-set question and the special positive results
of Bell, Burris and Yeats \cite{stab:BBY} concern precisely the kind of
additional recurrence information needed here.  The reduction above
leaves Chen's full moment-set conjecture unresolved; it is a
strengthening of the present structural argument, not a second named
problem resolved by this section.

\paragraph{Audit boundary.}
The claimed mathematical conclusion is the complete characterization
in Theorem~\ref{stab:classification}, not merely a theorem for rational,
logarithmic, or unramified inputs.  The decisive additional steps are
the all-place bounds before Riemann--Roch, the finite logarithmic-divisor
space modulo the lower field, and the quotient-module argument.
Independent reviews within this research session found no fatal gap in
these steps.  This is not external peer review, formal verification,
or proof that no identical argument exists in the literature.


\section{A Bell--cumulant criterion and functional digraphs}
\label{sec:tp}

This section proves a criterion for coefficientwise Hankel total positivity
of polynomial sequences defined by an exponential formula.  Its application
establishes the mathematical assertions in Sokal's
Conjecture~6.9(b) and~(d) \cite{Sokal2022}, with their original indexing and
normalization.  The substantive step is a passage from positive semidefinite
coefficient kernels to \emph{all} Hankel minors, using auxiliary polynomial
weights and Schur functions.  The proof does not assume that a canonical
production matrix is totally positive.

We use ``totally positive'' in the non-strict convention: every finite minor
is nonnegative.  A real sequence $(a_n)_{n\geq0}$ is a Stieltjes moment
sequence if $a_n=\int_{[0,\infty)}x^n\,d\rho(x)$ for a positive measure
with all moments finite.  The classical Stieltjes criterion identifies this
property with total positivity of $[a_{i+j}]_{i,j\geq0}$; equivalently,
both $[a_{i+j}]$ and $[a_{i+j+1}]$ are positive semidefinite on every
finite principal block.  For polynomials in $t$, coefficientwise Hankel
total positivity means that every minor of the Hankel matrix belongs to
$\mathbb R_{\geq0}[t]$.

\subsection{The exact criterion}

Let $(\kappa_n)_{n\geq1}$ be an arbitrary real sequence, and define formal
series and polynomials by
\begin{equation}
 K(z)=\sum_{n\geq1}\kappa_n\frac{z^n}{n!},
 \qquad
 M_t(z)=\exp(tK(z))=\sum_{n\geq0}m_n(t)\frac{z^n}{n!}.
 \label{eq:tp-bell-definition}
\end{equation}
Thus $m_0(t)=1$ and $m_n(t)$ is the complete exponential Bell polynomial
evaluated at $t\kappa_1,\ldots,t\kappa_n$.  No convergence of $K$ is
assumed.

\begin{theorem}[Bell--cumulant criterion]
\label{thm:tp-bell}
For the sequence in \eqref{eq:tp-bell-definition}, the following statements
are equivalent:
\begin{enumerate}
 \item The shifted cumulant sequence $(\kappa_{n+1})_{n\geq0}$ is a
 Stieltjes moment sequence.
 \item The sequence $(m_n(t))_{n\geq0}$ is coefficientwise
 Hankel-totally positive in $t$.
 \item For every real $t>0$, the numerical sequence
 $(m_n(t))_{n\geq0}$ is a Stieltjes moment sequence.
\end{enumerate}
\end{theorem}

The necessity direction uses the standard lowest-coefficient argument;
compare the related extraction in Zhu's Theorem~3.12(iii)
\cite{Zhu2021}. The substantive direction proved here is sufficiency.

We first give two elementary positivity lemmas.  Call a formal bivariate
series $A(u,v)=\sum_{i,j\geq0}a_{ij}u^iv^j$ a
\emph{positive semidefinite coefficient kernel} if every finite principal
block of its coefficient matrix $[a_{ij}]$ is real positive semidefinite.
This definition concerns coefficient matrices, rather than pointwise
evaluation of $A$.

\begin{lemma}
\label{lem:tp-kernel}
Positive semidefinite coefficient kernels are closed under addition,
multiplication, and multiplication by nonnegative real scalars.
\end{lemma}
\begin{proof}
Only multiplication requires explanation.  Fix a truncation degree $N$.
Write the required coefficient blocks of two kernels as Gram matrices,
$a_{ij}=\langle v_i,v_j\rangle$ and
$b_{ij}=\langle w_i,w_j\rangle$, for $0\leq i,j\leq N$.
For $0\leq n\leq N$, set
\[
 q_n=\sum_{i=0}^n v_i\otimes w_{n-i}.
\]
The coefficient of $u^nv^m$ in their product is
$\langle q_n,q_m\rangle$.  Its block through degree $N$ is therefore
positive semidefinite.  Since $N$ is arbitrary, this proves the assertion.
All sums used here are finite.
\end{proof}

\begin{lemma}
\label{lem:tp-mixed-determinant}
If $G_1,\ldots,G_d$ are real positive semidefinite $r\times r$ matrices,
then
\[
 \det\!\left(\sum_{j=1}^d z_jG_j\right)
 \in\mathbb R_{\geq0}[z_1,\ldots,z_d].
\]
\end{lemma}
\begin{proof}
Factor $G_j=V_jV_j^{\mathsf T}$ and concatenate the columns of the
$V_j$ into a matrix $V$.  Then $\sum_jz_jG_j=VDV^{\mathsf T}$, where
$D$ is diagonal and each column belonging to $V_j$ receives the diagonal
entry $z_j$.  Cauchy--Binet expands the determinant as a sum over
$r$-element column sets $S$ of
\[
 (\det V_S)^2\prod_{c\in S}D_{cc}.
\]
Every coefficient is a sum of nonnegative real squares.
\end{proof}

\begin{proof}[Proof of Theorem~\ref{thm:tp-bell}]
We prove the substantive implication (1)$\Rightarrow$(2) first.  Choose a
positive measure $\rho$ on $[0,\infty)$ such that
\begin{equation}
 \kappa_{n+1}=\int x^n\,d\rho(x),\qquad n\geq0.
 \label{eq:tp-cumulant-measure}
\end{equation}

\smallskip\noindent
\emph{Step 1: positive semidefinite coefficient kernels.}
For every integer $s\geq1$, the coefficient matrix of $K^{(s)}(u+v)$ is
\[
 \left[\frac{\kappa_{i+j+s}}{i!j!}\right]_{i,j\geq0}.
\]
Its finite quadratic forms are integrals of squares multiplied by
$x^{s-1}$, so it is positive semidefinite.  The same is true of
\begin{equation}
 C(u,v)=K(u+v)-K(u)-K(v)
       =\sum_{i,j\geq1}\frac{\kappa_{i+j}}{i!j!}u^iv^j.
 \label{eq:tp-centered-kernel}
\end{equation}
Indeed, its nonzero coefficient block has quadratic form
\[
 \sum_{i,j\geq1}c_ic_j\frac{\kappa_{i+j}}{i!j!}
 =\int x\left(\sum_{i\geq1}\frac{c_i x^{i-1}}{i!}\right)^2d\rho(x)
 \geq0,
\]
where the $c_i$ have finite support.  This expression remains valid when
$\rho$ has an atom at zero.

For $s\geq0$, put
\begin{equation}
 Q_s(u,v;t)=\frac{M_t^{(s)}(u+v)}{M_t(u)M_t(v)}.
 \label{eq:tp-Q}
\end{equation}
The quotients are well-defined formally because $M_t(0)=1$.
For $s\geq1$, Bell differentiation gives
\begin{equation}
 Q_s=e^{tC(u,v)}s!
 \sum_{\substack{\ell_1,\ldots,\ell_s\geq0\\
                   \ell_1+2\ell_2+\cdots+s\ell_s=s}}
 \prod_{j=1}^s
 \frac{t^{\ell_j}\bigl(K^{(j)}(u+v)\bigr)^{\ell_j}}
      {\ell_j!\,(j!)^{\ell_j}},
 \label{eq:tp-bell-Q}
\end{equation}
and $Q_0=e^{tC}$.  Lemma~\ref{lem:tp-kernel} shows that every coefficient
of $t$ in each $Q_s$ is a positive semidefinite coefficient kernel.
Moreover, every finite coefficient block of $Q_s$ is polynomial in $t$:
$C$ has positive degree in both $u$ and $v$, while the Bell polynomial in
\eqref{eq:tp-bell-Q} has finite degree in $t$.

\smallskip\noindent
\emph{Step 2: weighted leading determinants are Schur-positive.}
Fix $r\geq1$ and $p\geq1$.  Introduce auxiliary variables
$a=(a_1,\ldots,a_p)$ and the polynomial
\[
 P_a(x)=\prod_{\ell=1}^p(x+a_\ell)
       =\sum_{s=0}^p e_{p-s}(a)x^s,
\]
where $e_j$ denotes the elementary symmetric polynomial.  Let $L_t$ be
the $\mathbb R[t]$-linear functional on $\mathbb R[t][x]$ determined by
$L_t(x^n)=m_n(t)$.  Consider
\begin{equation}
 W_{r,p}(a;t)=
 \det\bigl[L_t(x^{i+j}P_a(x))\bigr]_{0\leq i,j<r}.
 \label{eq:tp-weighted-det}
\end{equation}
Its factorial-normalized matrix has generating kernel
\begin{align}
 \sum_{s=0}^p e_{p-s}(a)M_t^{(s)}(u+v)
 &=M_t(u)M_t(v)\sum_{s=0}^p e_{p-s}(a)Q_s(u,v;t).
 \label{eq:tp-Appell-kernel}
\end{align}
Let $G_s(t)$ be the $r\times r$ coefficient block of $Q_s$, and let
$A_r(t)$ be the unit lower-triangular matrix
\[
 (A_r)_{ij}=[z^{i-j}]M_t(z)\quad(i\geq j),
 \qquad (A_r)_{ij}=0\quad(i<j).
\]
Coefficient extraction in \eqref{eq:tp-Appell-kernel}, followed by
finite truncation, gives the exact congruence
\[
 \bigl[L_t(x^{i+j}P_a(x))\bigr]_{0\leq i,j<r}
 =D_rA_r\left(\sum_{s=0}^p e_{p-s}(a)G_s(t)\right)A_r^{\mathsf T}D_r,
 \qquad D_r=\operatorname{diag}(0!,1!,\ldots,(r-1)!).
\]
No terms outside the displayed finite blocks contribute: the factors
$A_r$ are lower triangular.  Since $\det A_r=1$, it follows that
\begin{equation}
 W_{r,p}(a;t)=\left(\prod_{i=0}^{r-1}(i!)^2\right)
 \det\left(\sum_{s=0}^p e_{p-s}(a)G_s(t)\right).
 \label{eq:tp-Appell-det}
\end{equation}

Write $G_s(t)=\sum_k t^kG_{s,k}$.  Each $G_{s,k}$ is a real positive
semidefinite matrix, by Step~1, and the sum is finite.  Apply
Lemma~\ref{lem:tp-mixed-determinant} with independent weights for these
matrices, and then substitute $t^ke_{p-s}(a)$ for the corresponding
weight.  Equation~\eqref{eq:tp-Appell-det} expresses $W_{r,p}$ as a
nonnegative real linear combination of terms
\[
 t^N e_{j_1}(a)\cdots e_{j_r}(a).
\]
The Pieri rule, or the Littlewood--Richardson rule, implies that products
of elementary symmetric polynomials are Schur-positive
\cite[Chapter~I]{Macdonald1995}.  Consequently every coefficient of $t$
in $W_{r,p}(a;t)$ has nonnegative coefficients when expanded in the Schur
basis in the $p$ auxiliary variables.

\smallskip\noindent
\emph{Step 3: extract every minor with consecutive initial row indices.}
For a partition $\lambda$ of length at most $r$, padded with zeros, define
\[
 J_\lambda=(\lambda_r,\lambda_{r-1}+1,\ldots,\lambda_1+r-1),
 \qquad \delta=(0,1,\ldots,r-1),
\]
and write $D_{I,J}(t)=\det[m_{i_a+j_b}(t)]_{a,b=1}^r$ for increasing
nonnegative index tuples $I,J$.  With
$\Delta(x)=\prod_{1\leq a<b\leq r}(x_b-x_a)$, the alternant formula
for Schur functions and the formal Andr\'eief identity give
\begin{equation}
 D_{\delta,J_\lambda}(t)
 =\frac1{r!}L_t^{\otimes r}\bigl(\Delta(x)^2s_\lambda(x)\bigr).
 \label{eq:tp-one-sided-minor}
\end{equation}
Here $L_t^{\otimes r}$ acts separately in the variables $x_1,\ldots,x_r$.
The identity is purely algebraic: expand both alternants, apply the
functional term by term, and group the $r!$ equal contributions obtained
by simultaneous permutation of the variables.  In particular, no
representing measure for $L_t$ is being assumed.

The rectangular form of the dual Cauchy identity is
\begin{equation}
 \prod_{i=1}^r\prod_{\ell=1}^p(x_i+a_\ell)
 =\sum_{\lambda\subseteq(p^r)}s_\lambda(x)s_{\widehat\lambda}(a),
 \qquad
 \widehat\lambda=(r-\lambda'_p,\ldots,r-\lambda'_1),
 \label{eq:tp-rectangle-Cauchy}
\end{equation}
where $\lambda'$ is the conjugate partition
\cite[Chapter~I]{Macdonald1995}.  For example, it follows from
the ordinary dual Cauchy identity after replacing $a_\ell$ by
$a_\ell^{-1}$ and multiplying by $\prod_\ell a_\ell^r$; the resulting
identity is polynomial.  Applying the formal Andr\'eief identity to
\eqref{eq:tp-weighted-det}, then using
\eqref{eq:tp-rectangle-Cauchy} and \eqref{eq:tp-one-sided-minor}, yields
\begin{equation}
 W_{r,p}(a;t)=\sum_{\lambda\subseteq(p^r)}
 D_{\delta,J_\lambda}(t)s_{\widehat\lambda}(a).
 \label{eq:tp-extract}
\end{equation}
The map $\lambda\mapsto\widehat\lambda$ is injective, and its images
have length at most $p$.  Their Schur polynomials are therefore linearly
independent in $p$ variables.  Step~2 and \eqref{eq:tp-extract} imply
\[
 D_{\delta,J_\lambda}(t)\in\mathbb R_{\geq0}[t].
\]
Every partition of length at most $r$ occurs by taking
$p\geq\max(1,\lambda_1)$.

\smallskip\noindent
\emph{Step 4: pass to arbitrary row and column index sets.}
Every increasing nonnegative tuple of length $r$ is uniquely $J_\lambda$
for some partition $\lambda$ of length at most $r$.  Thus for arbitrary
$I=J_\lambda$ and $J=J_\mu$, the same alternant calculation gives
\[
 D_{I,J}(t)=\frac1{r!}L_t^{\otimes r}
   \bigl(\Delta(x)^2s_\lambda(x)s_\mu(x)\bigr).
\]
The Littlewood--Richardson identity
\[
 s_\lambda s_\mu=\sum_\nu c_{\lambda\mu}^{\nu}s_\nu,
 \qquad c_{\lambda\mu}^{\nu}\in\mathbb Z_{\geq0},
\]
holds after specialization to $r$ variables, with terms of length
greater than $r$ vanishing \cite[Chapter~I]{Macdonald1995}.
It follows that
\begin{equation}
 D_{I,J}(t)=\sum_{\ell(\nu)\leq r}
             c_{\lambda\mu}^{\nu}D_{\delta,J_\nu}(t)
           \in\mathbb R_{\geq0}[t].
 \label{eq:tp-LR-minors}
\end{equation}
This proves (2) for all minors, not just principal or solid minors.

\smallskip\noindent
To finish the equivalence, (2)$\Rightarrow$(3) follows by evaluation at
$t>0$ and the classical Stieltjes criterion.  Suppose (3) holds.  For
arbitrary increasing nonnegative tuples $I=(i_a)$ and $J=(j_b)$ of length
$r$, the shifted minor
\[
 \det[m_{i_a+j_b+1}(t)]_{a,b=1}^r
\]
is nonnegative for every $t>0$, since it is a minor of the moment Hankel
matrix, using rows $I$ and columns $J+1$.  For $n\geq1$,
$m_n(t)=t\kappa_n+O(t^2)$ as a polynomial at zero.  Dividing this minor
by $t^r$ and taking $t\downarrow0$ proves
\[
 \det[\kappa_{i_a+j_b+1}]_{a,b=1}^r\geq0.
\]
The shifted cumulant Hankel matrix is totally positive.  The Stieltjes
criterion now gives (1).
\end{proof}

\begin{remark}
The measure in the hypothesis of Theorem~\ref{thm:tp-bell} is numerical.
The proof does not assert that replacing the $\kappa_n$ by arbitrary
polynomials with coefficientwise Hankel positivity preserves the same
conclusion jointly in all parameters.  Its determinant lemma uses real
positive semidefinite coefficient matrices.  This distinction matters
for refinements by additional combinatorial statistics.
\end{remark}

\subsection{An explicit measure for connected functional digraphs}

Let $T(z)$ be the rooted-tree series, characterized by $T(0)=0$ and
$T(z)=ze^{T(z)}$, and put
\begin{equation}
 K(z)=-\log(1-T(z))=\sum_{n\geq1}\kappa_n\frac{z^n}{n!}.
 \label{eq:tp-functional-K}
\end{equation}
The following proposition establishes the exact shifted-cumulant input
needed above, as well as its factorial-normalized version.

\begin{proposition}
\label{prop:tp-Lambert-measure}
Let $W$ be the principal Lambert function.  The function
\begin{equation}
 w(x)=\frac1\pi\operatorname{Arg}\bigl(1+W(-1/x+i0)\bigr),
 \qquad 0<x<e,
 \label{eq:tp-density}
\end{equation}
is a probability density on $(0,e)$ and satisfies $0<w(x)<1/2$.
For the coefficients in \eqref{eq:tp-functional-K},
\begin{equation}
 \frac{\kappa_{n+1}}{(n+1)!}
   =\int_0^e x^n w(x)\,dx,\qquad n\geq0.
 \label{eq:tp-normalized-moments}
\end{equation}
Consequently both $(\kappa_{n+1}/(n+1)!)_{n\geq0}$ and
$(\kappa_{n+1})_{n\geq0}$ are Stieltjes moment sequences.
\end{proposition}

\begin{proof}
The identity $T(-s)=-W(s)$ gives
\[
 h(s)=-K(-s)=\log(1+W(s)),\qquad F(s)=\frac{h(s)}s,
\]
where the singularity of $F$ at zero is removable and $F(0)=1$.
We derive its Stieltjes representation directly.  We use the standard
principal-branch properties of $W$, including its cut
$(-\infty,-1/e]$ and its boundary parametrization
\cite{Lambert2012}.

For $t>1/e$, there is a unique $v\in(0,\pi)$ such that
\begin{equation}
 t=\frac{v}{\sin v}e^{-v\cot v},
 \qquad W(-t+i0)=-v\cot v+iv.
 \label{eq:tp-W-boundary}
\end{equation}
The parametrization follows by imposing that $We^W$ be negative real.
For completeness, its $t$-coordinate increases strictly from $1/e$ to
infinity: the derivative of its logarithm is
$1/v-2\cot v+v\csc^2v$, whose product with $v\sin^2v$ is
$(\sin v-v\cos v)^2+v^2\sin^2v>0$.
Furthermore $1-v\cot v>0$ because
$\sin v-v\cos v$ has derivative $v\sin v>0$ and limit zero at the
origin.  Hence
\[
 \theta(t):=\operatorname{Arg}(1+W(-t+i0))\in(0,\pi/2),
\]
with $\theta(1/e+)=\pi/2$ and $\theta(t)\to0$ as $t\to\infty$.

The principal logarithm in $h$ introduces no additional cut: $W$ maps
the upper half-plane into itself, while $1+W(s)>0$ for real $s>-1/e$.
Thus $F$ is holomorphic on $\mathbb C\setminus(-\infty,-1/e]$, after
removing its singularity at zero, and is real symmetric.  On the upper
side of its cut,
\begin{equation}
 \operatorname{Im}F(-t+i0)=-\frac{\theta(t)}t.
 \label{eq:tp-boundary-imaginary}
\end{equation}
At the branch endpoint, the square-root expansion of $W+1$ shows that
$F$ has only a logarithmic singularity.  At infinity, the principal-branch
asymptotics $W(s)=\log s-\log\log s+o(1)$ give
$F(s)=O(\log\log|s|/|s|)$ along a large keyhole contour.
Accordingly the small endpoint-circle and large-circle terms in Cauchy's
formula both vanish.  Combining the two sides of the cut using
\eqref{eq:tp-boundary-imaginary} gives, for $s>0$,
\begin{equation}
 F(s)=\frac1\pi\int_{1/e}^{\infty}
             \frac{\theta(t)}{t(t+s)}\,dt
     =\int_0^e\frac{w(x)}{1+sx}\,dx,
 \label{eq:tp-Stieltjes-F}
\end{equation}
where the second equality is the substitution $x=1/t$.
The bounds on $\theta$ show that $w$ is positive and bounded by $1/2$.
Letting $s\downarrow0$ in \eqref{eq:tp-Stieltjes-F} gives
$\int_0^e w(x)\,dx=F(0)=1$.

Both sides of \eqref{eq:tp-Stieltjes-F} are analytic for $|s|<1/e$.
The geometric series for $(1+sx)^{-1}$ converges uniformly on every
smaller closed disk, uniformly for $x\in[0,e]$.  Since
\[
 F(s)=\sum_{n\geq0}(-1)^n\frac{\kappa_{n+1}}{(n+1)!}s^n,
\]
comparison of Taylor coefficients proves
\eqref{eq:tp-normalized-moments}.

Finally let $X$ have density $w$ and, independently, let $G$ have the
gamma density $g e^{-g}$ on $(0,\infty)$.  Since
$\E[G^n]=\Gamma(n+2)=(n+1)!$, we obtain
\[
 \E[(XG)^n]=(n+1)!\int_0^e x^n w(x)\,dx=\kappa_{n+1}.
\]
Thus the product $XG$ supplies a positive representing measure for the
unnormalized shifted cumulants as well.
\end{proof}

\begin{remark}
The existence assertion in Proposition~\ref{prop:tp-Lambert-measure}
also follows from established complete-Bernstein theory.  The principal
function $W$ is complete Bernstein \cite{Lambert2012}; so is
$\log(1+s)$, and complete Bernstein functions are closed under composition.
It follows that $h(s)=\log(1+W(s))$ is complete Bernstein and that $h(s)/s$
is Stieltjes \cite{Schilling2012}.  The contour calculation above makes
the representing density and normalization explicit.  The analytic
closure principle itself is an existing theorem, not a new ingredient
claimed here.
\end{remark}

\subsection{The original functional-digraph assertions}

A functional digraph on $[n]$ has exactly one outgoing edge at each
vertex, with loops permitted.  Let $\psi_{n,k}$ count such digraphs with
$k$ weakly connected components, and define
\[
 \Psi_n(y)=\sum_{k=0}^n\psi_{n,k}y^k,
 \qquad \Psi_0(y)=1.
\]
Each connected component is a directed cycle of rooted labelled trees.
The labelled exponential formula therefore gives
\begin{equation}
 \sum_{n\geq0}\Psi_n(y)\frac{z^n}{n!}
  =\exp\!\left(y\sum_{\ell\geq1}\frac{T(z)^\ell}{\ell}\right)
  =(1-T(z))^{-y}=e^{yK(z)},
 \qquad \kappa_n=\psi_{n,1}.
 \label{eq:tp-functional-EGF}
\end{equation}
These are exactly the definitions used in Sokal's Conjecture~6.9
\cite{Sokal2022}.

\begin{corollary}[Sokal's Conjecture~6.9(b) and~(d)]
\label{cor:tp-Sokal}
Every finite minor of $[\Psi_{i+j}(y)]_{i,j\geq0}$ has nonnegative
coefficients in $y$.  In addition,
\[
 \left(\frac{\psi_{n+1,1}}{(n+1)!}\right)_{n\geq0}
\]
is a Stieltjes moment sequence, with the probability density
\eqref{eq:tp-density} as a representing measure.
\end{corollary}
\begin{proof}
Proposition~\ref{prop:tp-Lambert-measure} gives the numerical Stieltjes
property of $(\kappa_{n+1})_{n\geq0}$.  Apply
Theorem~\ref{thm:tp-bell} with $t=y$ and use
\eqref{eq:tp-functional-EGF}.  The normalized assertion is precisely
\eqref{eq:tp-normalized-moments} with $\kappa_n=\psi_{n,1}$.
\end{proof}

Zhu's Proposition~6.18 \cite{Zhu2024} already establishes the
coefficient-triangle assertion~6.9(a) and the numerical shifted-moment
assertion~6.9(c), and proves pointwise Stieltjes positivity of
$\Psi_n(y)$.  The argument here proves the full coefficientwise
assertion~(b), and supplies a direct measure for~(d).  It does not rely
on the indexing of the cumulant summation displayed in that earlier
proof.  The four mathematical clauses of the original Conjecture~6.9
are thus accounted for, with the earlier results explicitly credited.

\begin{remark}
No joint coefficientwise assertion is made for the additional
cyclic-vertex parameter in Sokal's Conjecture~6.10, or for the
unnormalized ordered-forest sequence in Conjecture~6.7.  Neither is
obtained by treating the numerical cumulant hypothesis as a
coefficientwise hypothesis in a second indeterminate.
\end{remark}

\subsection{An obstruction to the canonical production-matrix criterion}

The following exact computation explains why the usual sufficient
criterion based on the canonical exponential-Riordan production matrix
does not directly prove Corollary~\ref{cor:tp-Sokal}.

\begin{proposition}
\label{prop:tp-production-obstruction}
The canonical production matrix of the coefficient triangle
$[\psi_{n,k}]_{n,k\geq0}$ has a negative entry:
\[
 P_{17,1}=-1120228.
\]
\end{proposition}
\begin{proof}
The column exponential generating functions are $K(z)^k/k!$.
Differentiating them shows that the production entries are
\[
 P_{n,k}=\frac{n!}{(k-1)!}[w^{n-k+1}]A(w)\quad(k\geq1),
 \qquad P_{n,0}=0,
 \qquad A(w)=K'(K^{-1}(w)).
\]
Indeed, expanding $A(K(z))K(z)^{k-1}/(k-1)!$ in the basis
$K(z)^n/n!$ gives exactly these coefficients.
For the present $K$, write $T=1-e^{-w}$ and $z=Te^{-T}$.
Since $K'(z)=e^T/(1-T)^2$, this gives
\begin{equation}
 A(w)=\exp(2w+1-e^{-w}).
 \label{eq:tp-production-A}
\end{equation}
Let $a_n=A^{(n)}(0)$, and set $\ell_1=3$ and
$\ell_j=(-1)^{j+1}$ for $j\geq2$.  Differentiating
$A'=A(2+e^{-w})$ repeatedly gives the integer recurrence
\[
 a_0=1,\qquad
 a_{n+1}=\sum_{j=0}^n\binom nj\ell_{j+1}a_{n-j}.
\]
It yields, by exact integer arithmetic,
\begin{align*}
 (a_0,\ldots,a_8)
   &=(1,3,8,19,41,84,171,347,690),\\
 (a_9,\ldots,a_{17})
   &=(1385,2825,5438,11077,24535,33720,102623,350605,-1120228).
\end{align*}
Finally $P_{17,1}=17![w^{17}]A(w)=a_{17}$.
\end{proof}

Thus this particular production matrix is not even entrywise
nonnegative.  This observation excludes direct application of a
sufficient criterion requiring its total positivity; it does not
exclude other production matrices or alternative path representations.


\section{The DeTemple quotient in Yang's monotonicity conjecture}
\label{sec:yang}

Let $\psi=\Gamma'/\Gamma$ denote the digamma function and put
\begin{equation}\label{eq:yang-R}
 R(x)=\psi\!\left(x+\tfrac12\right)-\log x,\qquad x>0.
\end{equation}
Yang~\cite[Conjecture~1(i)]{Yang2014} conjectured that the quotient
\begin{equation}\label{eq:yang-original}
 H(x)=
 \frac{-24\left(x^4-\frac{31}{336}\right)R(x)
                  +x^2-\frac7{40}}
      {24\left(x^2+\frac7{40}\right)R(x)-1}
\end{equation}
is increasing on $(0,\infty)$. The same assertion appears as
Conjecture~1(i) in the published article of Yang, Chu, and
Zhang~\cite{YangChuZhang2015}. We prove it for the stated quotient on
the entire stated domain.

\begin{theorem}\label{thm:yang-quotient}
The denominator in \eqref{eq:yang-original} is strictly positive, and
$H'(x)>0$ for every $x>0$. Moreover,
\begin{equation}\label{eq:yang-sharp-range}
 \frac{155}{294}<H(x)<\frac{11165}{8284},\qquad x>0,
\end{equation}
and both bounds are sharp, being the limits at zero and infinity,
respectively.
\end{theorem}

The proof identifies the derivative of a quotient of Stieltjes transforms
with the deficit in the covariance form of the Cauchy--Schwarz inequality.
We first establish this identity for an arbitrary positive measure.

\begin{proposition}[A variance identity for a Stieltjes quotient]
\label{prop:yang-variance}
Let $T$ be a nonnegative random variable with finite second moment whose
support contains at least three points. Define
\[
 m_1=\E[T],\qquad m_2=\E[T^2],\qquad
 F(s)=\E\!\left[\frac1{s+T}\right],\qquad s>0,
\]
and
\[
 D(s)=(s+m_1)F(s)-1.
\]
Then $D(s)>0$, and the function
\begin{equation}\label{eq:yang-J}
 J(s)=\frac{-(s^2-m_2)F(s)+s-m_1}{D(s)}
\end{equation}
satisfies
\begin{equation}\label{eq:yang-variance-identity}
 J'(s)=
 \frac{\Var(T)\Var((s+T)^{-1})
          -\Cov(T,(s+T)^{-1})^2}{D(s)^2}>0.
\end{equation}
\end{proposition}

\begin{proof}
Fix $s>0$ and write $U=(s+T)^{-1}$. Since
$TU=1-sU$, we have
\begin{equation}\label{eq:yang-D-covariance}
 D(s)=-\Cov(T,U).
\end{equation}
If $T'$ is an independent copy of $T$, with $U'=(s+T')^{-1}$, then
\[
 2\Cov(T,U)=\E[(T-T')(U-U')]<0.
\]
Indeed, $t\mapsto(s+t)^{-1}$ is strictly decreasing, and
$\mathbb P(T\ne T')>0$. All expectations are finite. This proves
$D(s)>0$.

Put $\sigma^2=m_2-m_1^2$. Algebraic rearrangement of
\eqref{eq:yang-J} gives
\begin{equation}\label{eq:yang-J-rearranged}
 J(s)=\frac{\sigma^2F(s)}{D(s)}-s+m_1.
\end{equation}
Differentiation under the expectation is justified locally in $s>0$
by the bound $(s+T)^{-2}\le s^{-2}$. Thus
\[
 F'(s)=-\E[U^2],\qquad
 D'(s)=F(s)+(s+m_1)F'(s).
\]
It follows that
\[
 F'(s)D(s)-F(s)D'(s)
 =-F'(s)-F(s)^2
 =\Var(U).
\]
Differentiating \eqref{eq:yang-J-rearranged} and using
\eqref{eq:yang-D-covariance} proves the identity in
\eqref{eq:yang-variance-identity}.

Cauchy--Schwarz makes its numerator nonnegative. Equality would imply
$U=aT+b$ almost surely for some real $a,b$. Consequently the support of
$T$ would lie among the zeros of
\[
 (s+t)(at+b)-1.
\]
This is a nonzero polynomial of degree at most two: it cannot vanish
identically, as comparison of its coefficients immediately shows. Its
zero set contains at most two points, contrary to the support hypothesis.
The inequality is therefore strict.
\end{proof}

\begin{lemma}[The positive measure for the DeTemple remainder]
\label{lem:yang-measure}
For $x>0$,
\begin{equation}\label{eq:yang-stieltjes}
 R(x)=\int_0^\infty
 \frac{du}{(x^2+u)(e^{2\pi\sqrt{u}}+1)}.
\end{equation}
The measure
\begin{equation}\label{eq:yang-probability}
 d\mu(u)=\frac{24\,du}{e^{2\pi\sqrt{u}}+1},\qquad u>0,
\end{equation}
is a probability measure. Its first three moments are
\begin{equation}\label{eq:yang-moments}
 m_1=\frac7{40},\qquad m_2=\frac{31}{336},\qquad
 m_3=\frac{127}{1280}.
\end{equation}
\end{lemma}

\begin{proof}
Binet's integral for the digamma function reads
\[
 \psi(x)=\log x-\frac1{2x}
 -2\int_0^\infty
       \frac{t\,dt}{(x^2+t^2)(e^{2\pi t}-1)}.
\]
Combining it with the duplication identity
\[
 \psi\!\left(x+\tfrac12\right)
       =2\psi(2x)-\psi(x)-2\log2
\]
and changing the variable $t$ to $2t$ in the integral arising from
$\psi(2x)$ gives
\begin{align*}
 R(x)
 &=2\int_0^\infty\frac{t}{x^2+t^2}
       \left(\frac1{e^{2\pi t}-1}
                    -\frac2{e^{4\pi t}-1}\right)dt\\
 &=2\int_0^\infty
       \frac{t\,dt}{(x^2+t^2)(e^{2\pi t}+1)}.
\end{align*}
The substitution $u=t^2$ proves \eqref{eq:yang-stieltjes}.

For every integer $j\ge0$, direct integration yields
\begin{equation}\label{eq:yang-all-moments}
 \int_0^\infty u^j\,d\mu(u)
 =\frac{48\,\Gamma(2j+2)}{(2\pi)^{2j+2}}
       \bigl(1-2^{-2j-1}\bigr)\zeta(2j+2).
\end{equation}
For example, the identity follows by setting $v=2\pi\sqrt u$ and
expanding $(e^v+1)^{-1}=\sum_{\ell\ge1}(-1)^{\ell-1}e^{-\ell v}$.
The integrated series is absolutely summable for the powers
$v^{2j+1}$ occurring here, so termwise integration is justified.
At $j=0$ the expression in \eqref{eq:yang-all-moments} is $1$.
At $j=1,2,3$ the usual evaluations of the even zeta values give
\eqref{eq:yang-moments}.
\end{proof}

\begin{proof}[Proof of Theorem~\ref{thm:yang-quotient}]
Let $T$ have distribution \eqref{eq:yang-probability}. Its density is
strictly positive throughout $(0,\infty)$, and it has moments of every
order. Lemma~\ref{lem:yang-measure} identifies its Stieltjes transform as
\[
 F(s)=24R(\sqrt s).
\]
With the moment values in \eqref{eq:yang-moments}, the denominator in
\eqref{eq:yang-original} equals $D(x^2)$, and its numerator equals the
numerator in \eqref{eq:yang-J} evaluated at $s=x^2$. Thus
\[
 H(x)=J(x^2),\qquad H'(x)=2xJ'(x^2)>0
\]
by Proposition~\ref{prop:yang-variance}. The same proposition proves that
the denominator is positive everywhere on the original domain.

It remains to evaluate the endpoints. Since $\psi$ is continuous at
$1/2$,
\[
 R(x)=-\log x+\psi(1/2)+o(1)\qquad(x\downarrow0).
\]
Consequently, dividing the numerator and denominator of $J(s)$ by
$F(s)$ shows that
\[
 \lim_{x\downarrow0}H(x)=\frac{m_2}{m_1}=\frac{155}{294}.
\]
At infinity, the finite third moment justifies the Stieltjes expansion
\[
 F(s)=\frac1s-\frac{m_1}{s^2}+\frac{m_2}{s^3}
                         -\frac{m_3}{s^4}+o(s^{-4}).
\]
For clarity, this follows from finite geometric division of
$(s+T)^{-1}$ and dominated convergence; convergence of an infinite
asymptotic series is not being assumed. Substitution gives
\[
 D(s)=\frac{m_2-m_1^2}{s^2}+o(s^{-2}),\qquad
 -(s^2-m_2)F(s)+s-m_1
       =\frac{m_3-m_1m_2}{s^2}+o(s^{-2}).
\]
Therefore
\[
 \lim_{x\to\infty}H(x)
 =\frac{m_3-m_1m_2}{m_2-m_1^2}
 =\frac{11165}{8284}.
\]
Strict monotonicity now proves \eqref{eq:yang-sharp-range}, including
sharpness of both constants.
\end{proof}

\begin{remark}\label{rem:yang-scope}
The result proves the separately stated monotonicity assertion in
Conjecture~1(i). The all-orders coefficient problem is treated separately
in Section~\ref{sec:yang-allorders}. The companion Conjecture~1(ii) is
false: Kang's earlier preprint \cite{Kang2026} supplies a strictly interior
counterexample. It is not counted as a new result here. The integral representation and Cauchy--Schwarz
inequality used here are classical; the quotient identity
\eqref{eq:yang-variance-identity} supplies their application to the
specific conjecture.
\end{remark}


\section{All orders in Yang's complete-monotonicity problem}
\label{sec:yang-allorders}

\paragraph{Problem and result.}
Let
\[
 R(x)=\psi(x+\tfrac12)-\log x,\qquad x>0,
\]
where $\psi=\Gamma'/\Gamma$.
Problem~1 of Yang's 2014 preprint \cite{Yang2014}, repeated in the 2015 article of
Yang, Chu, and Zhang \cite{YangChuZhang2015}, asks for the coefficients for which
\[
 \left(\sum_{j=0}^{n+1}a_jx^{2j}\right)R(x)
       -\sum_{j=0}^{n}b_jx^{2j}
\]
is completely monotone and has a finite nonzero leading coefficient at
order $x^{-2n-4}$. There is an unavoidable positive scaling freedom.
The theorem below determines every solution and proves strict complete
monotonicity at every order. The argument is presented as an internally audited complete-proof candidate.
The bounded priority review is recorded in the status section; historical
novelty is not certified.

Define the probability measure
\begin{equation}\label{ya:mu}
 d\mu(u)=\frac{24\,du}{e^{2\pi\sqrt u}+1},\qquad u>0,
\end{equation}
and its moments
\begin{equation}\label{ya:moments}
 m_j=\int u^j\,d\mu(u)
 =\frac{24(1-2^{-2j-1})|B_{2j+2}|}{2j+2},\qquad j\ge0.
\end{equation}
Here $B_j$ are the Bernoulli numbers; in particular $m_0=1$.
Let $\pi_N$ be the monic orthogonal polynomial of degree $N$ for $\mu$,
and write
\[
 H_N=\int\pi_N(u)^2\,d\mu(u)>0,
 \qquad P_N(s)=(-1)^N\pi_N(-s)=\sum_{j=0}^Na_js^j.
\]
Thus $a_N=1$. Put
\begin{equation}\label{ya:Q}
 Q_{N-1}(s)=\int
       \frac{P_N(s)-P_N(-u)}{s+u}\,d\mu(u).
\end{equation}
The quotient in the integrand is a polynomial, so $Q_{N-1}$ has degree
at most $N-1$.

\begin{theorem}\label{ya:main}
For every $N\ge1$, the function
\begin{equation}\label{ya:E}
 E_N(x)=P_N(x^2)R(x)-\frac1{24}Q_{N-1}(x^2)
\end{equation}
is strictly completely monotone on $(0,\infty)$: for every integer
$k\ge0$ and every $x>0$,
\[
 (-1)^kE_N^{(k)}(x)>0.
\]
Moreover,
\begin{equation}\label{ya:asymptotic}
 \lim_{x\to\infty}x^{2N+2}E_N(x)=\frac{H_N}{24}.
\end{equation}
Every solution of the stated problem with $n=N-1$ is a positive
constant multiple of \eqref{ya:E}.

In explicit finite-dimensional terms, if
\[
 \Delta_N=\det(m_{i+j})_{i,j=0}^{N-1},\qquad \Delta_0=1,
\]
then
\begin{equation}\label{ya:determinant}
 \pi_N(u)=\frac1{\Delta_N}
 \det\begin{pmatrix}
 m_0&m_1&\cdots&m_N\\
 m_1&m_2&\cdots&m_{N+1}\\
 \vdots&\vdots&&\vdots\\
 m_{N-1}&m_N&\cdots&m_{2N-1}\\
 1&u&\cdots&u^N
 \end{pmatrix},\qquad
 H_N=\frac{\Delta_{N+1}}{\Delta_N}.
\end{equation}
The normalized coefficients in the question are the coefficients
$a_j$ of $P_N$, together with
\begin{equation}\label{ya:bcoefficients}
 b_k=\frac1{24}\sum_{j=k+1}^N
       a_j(-1)^{j-k-1}m_{j-k-1},\qquad 0\le k<N.
\end{equation}
\end{theorem}

The structural question behind the proof is when a Dirichlet average of
an orthogonal-polynomial ensemble is a mixture of gamma laws of one
fixed shape. Lemma~\ref{ya:matrix-marginal} supplies a sufficient
spectral condition, and the Yang weight satisfies it at every matrix size.
This new auxiliary question is not counted as a pre-existing conjecture.
The decisive positivity statement is obtained from a one-dimensional
marginal of a positive Hermitian matrix. We give the matrix argument
and its transform identity explicitly; no convergence of an infinite
moment expansion is assumed.

\subsection{A matrix marginal and gamma mixtures}

For a positive function $f$ and a positive-definite Hermitian matrix
$X$, the notation $\det f(X)$ means the product of $f$ evaluated at the
eigenvalues of $X$.

\begin{lemma}[A spectral product class]\label{ya:spectral-class}
Suppose
\begin{equation}\label{ya:product-class}
 f(u)=C u^{-\beta}e^{-c\sqrt u}
       \prod_{j\ge1}(1+d_ju)^{-\gamma_j},\qquad u>0,
\end{equation}
where $C>0$, $\beta,c,\gamma_j\ge0$, $d_j>0$, and
$\sum_j\gamma_jd_j<\infty$; finite products are allowed.
If $A\ge0$ is Hermitian, $e\ne0$, and $A+see^*>0$ for every $s>0$,
then
\[
 s\longmapsto\det f(A+see^*)
\]
is completely monotone on $(0,\infty)$.
\end{lemma}
\begin{proof}
The matrix determinant lemma shows that
$\det(I+d_j(A+see^*))$ is a positive constant times
$1+s d_je^*(I+d_jA)^{-1}e$. Its negative power is completely
monotone. Also $\det(A+see^*)$ is affine in $s$, with nonnegative
coefficients and positive values for $s>0$; its negative power has the
same property. This assertion follows directly from the determinant
lemma when $A>0$, and by a limit from $A+\varepsilon I$ otherwise.

For the square-root factor, the scalar identity
\[
 \sqrt{x}=\frac1{2\pi}\int_0^\infty
                t^{-1/2}\log(1+x/t)\,dt,\qquad x\ge0,
\]
first follows by differentiation for $x>0$ and then by continuity at
zero. Functional calculus and the determinant lemma give
\begin{equation}\label{ya:root-log}
 \operatorname{Tr}\sqrt{A+see^*}-\operatorname{Tr}\sqrt A
 =\frac1{2\pi}\int_0^\infty
        t^{-1/2}\log\bigl(1+s b(t)\bigr)\,dt,
 \qquad b(t)=e^*(A+tI)^{-1}e\ge0.
\end{equation}
The integral converges, including for singular $A$, since
$b(t)\le\|e\|^2/t$: the resulting majorant is integrable both at zero
and at infinity for each fixed $s>0$.
Approximate this positive integral by finite positive quadrature sums
on compact subintervals of $(0,\infty)$ and then let those intervals
increase. The corresponding exponentials are positive products of
functions $(1+s b)^{-a}$ with $a\ge0$, hence are completely monotone.
The class of completely monotone functions is closed under
pointwise limits that are finite on $(0,\infty)$. Therefore
$\exp(-c\operatorname{Tr}\sqrt{A+see^*})$ is completely monotone.
Finally, the infinite product in \eqref{ya:product-class} converges to
a positive finite value on every positive-definite matrix, by
$\sum_j\gamma_jd_j<\infty$. Apply the same pointwise-limit closure to
its finite products and combine the factors.
\end{proof}

\begin{lemma}[Gamma-mixture matrix marginal]\label{ya:matrix-marginal}
Let $\alpha>0$, let $r\ge1$, and suppose a probability measure $\nu$
on $(0,\infty)$ has density
\[
 w(u)=u^{\alpha-1}f(u)>0
\]
and moments of every order. Assume that
$s\mapsto\det f(A+see^*)$ is completely monotone whenever
$A\ge0$ and $A+see^*>0$ for $s>0$.
Let $(U_1,\ldots,U_r)$ have density proportional to
\[
 \prod_{i<j}(U_i-U_j)^2\prod_{i=1}^r w(U_i)
\]
on $(0,\infty)^r$. Independently let
$(V_1,\ldots,V_r)$ have the Dirichlet$(1,\ldots,1)$ distribution,
and put $S=\sum_iV_iU_i$.
Then $S$ is a positive mixture of gamma laws with common shape
$p=\alpha+r-1$ and varying positive rates.
\end{lemma}
\begin{proof}
Consider the probability density proportional to
\[
 \det(X)^{\alpha-1}\det f(X)
\]
on the cone of positive-definite complex Hermitian $r\times r$
matrices, with respect to Lebesgue measure. Its integral is finite
because the spectral Jacobian is a constant times
$\prod_{i<j}(U_i-U_j)^2$ and all the required moments are finite.
The spectral change of variables also shows that its eigenvalues have
the indicated density and that its eigenvectors are independently
Haar unitary. The squared absolute values of the first row of a Haar
unitary matrix have the Dirichlet$(1,\ldots,1)$ distribution. For
example, this follows by normalizing independent standard complex
Gaussians, whose squared absolute values are independent exponentials.
Consequently $X_{11}$ has the law of $S$.

For $r>1$, use the Schur-complement parametrization
\begin{equation}\label{ya:schur}
 X=\begin{pmatrix}
 s&\sqrt{s}\,z^*\\
 \sqrt{s}\,z&B+zz^*
 \end{pmatrix},\qquad
 s>0,\quad B>0,\quad z\in\mathbb C^{r-1}.
\end{equation}
It is bijective, $\det X=s\det B$, and its real Jacobian is
$s^{r-1}$. Indeed the $2r-2$ real coordinates of the off-diagonal
vector are scaled by $\sqrt s$, while translation of the lower block
has Jacobian one. Furthermore $X=LL^*$, where
\[
 L=\begin{pmatrix}\sqrt s&0\\z&B^{1/2}\end{pmatrix}.
\]
Thus $X$ has the same eigenvalues as
\[
 L^*L=A+s e_1e_1^*,\qquad
 A=\begin{pmatrix}
 z^*z&z^*B^{1/2}\\B^{1/2}z&B
 \end{pmatrix}\ge0.
\]
The density of $S$ consequently has the form
\begin{equation}\label{ya:marginal-density}
 p_S(s)=s^{p-1}G(s),\qquad
 G(s)=C_r\int_{B>0}\int_{\mathbb C^{r-1}}
       (\det B)^{\alpha-1}\det f(A+s e_1e_1^*)\,dz\,dB.
\end{equation}
Each integrand is completely monotone in $s$. Positive integration
preserves complete monotonicity whenever the result is finite. Here
finiteness holds for every $s>0$: the integral is nonincreasing, so an
infinite value at any $s_0>0$ would contradict
$\int_0^\infty s^{p-1}G(s)\,ds=1$.
This argument can be formalized by exhausting the parameter domain
with compact sets and applying pointwise-limit closure to the finite
integrals. The case $r=1$ is the same calculation without $B,z$.

By Bernstein's theorem there is a positive measure $\rho$ on
$[0,\infty)$ such that
$G(s)=\int e^{-\lambda s}\rho(d\lambda)$.
Tonelli's theorem and normalization yield
\[
 1=\Gamma(p)\int_{(0,\infty)}\lambda^{-p}\rho(d\lambda),
 \qquad \rho(\{0\})=0.
\]
Hence $\tau(d\lambda)=\Gamma(p)\lambda^{-p}\rho(d\lambda)$ is a
probability measure and
\[
 p_S(s)=\int_{(0,\infty)}
       \frac{\lambda^p}{\Gamma(p)}s^{p-1}e^{-\lambda s}
       \,\tau(d\lambda),
\]
which is the claimed gamma mixture.
\end{proof}

\subsection{An exact orthogonal-polynomial transform}

\begin{lemma}[Divided-difference identity]\label{ya:divided-difference}
Let $\nu$ be a probability measure on $(0,\infty)$ with all moments
and infinite support. Let $\pi_N$ be its monic orthogonal polynomial
and $H_N=\int\pi_N^2\,d\nu$. Take $r=N+1$ and form $U,V,S$ from
$\nu$ as in Lemma~\ref{ya:matrix-marginal}.
For a $C^N$ function $F$ for which the following integrals converge,
\begin{equation}\label{ya:divided-difference-identity}
 \int\pi_N(u)F(u)\,d\nu(u)
       =\frac{H_N}{N!}\,\mathbb E[F^{(N)}(S)].
\end{equation}
In particular, this identity applies to
$F(u)=\sin(t\sqrt u)/\sqrt u$ for every fixed $t>0$.
\end{lemma}
\begin{proof}
Write $Z_r=\int\prod_{i<j}(u_i-u_j)^2\prod_i d\nu(u_i)$.
Expanding the Vandermonde determinants and using orthogonality gives
$Z_r=r!\prod_{j=0}^{r-1}H_j$. Heine's formula for $\pi_N$, followed
by symmetrization in the $N+1$ integration variables, then gives
\[
 \int\pi_N(u)F(u)\,d\nu(u)
       =H_N\,\mathbb E\bigl[[U_1,\ldots,U_{N+1}]F\bigr].
\]
Explicitly, the symmetrized numerator is the full Vandermonde square
multiplied by the divided difference
\[
 \sum_i\frac{F(U_i)}{\prod_{j\ne i}(U_i-U_j)}.
\] The normalizing factor is
$Z_{N+1}/((N+1)N!\prod_{j<N}H_j)=H_N$.
The Hermite--Genocchi formula gives
\[
 [U_1,\ldots,U_{N+1}]F
 =\frac1{N!}\mathbb E\left[
 F^{(N)}\left(\sum_iV_iU_i\right)\,\middle|\,U\right],
\]
proving \eqref{ya:divided-difference-identity}.
For the stated sine function,
\begin{equation}\label{ya:bessel-derivative}
 F^{(N)}(u)
 =\frac{(-1)^Nt^{2N+1}}{4^N(3/2)_N}
       {}_0F_1\left(N+\tfrac32;-\frac{t^2u}{4}\right).
\end{equation}
The derivative formula follows from the entire power series of the
single-variable function $F$, not from any moment expansion of
$\nu$. The bound
$|{}_0F_1(p;-v)|\le1$ for $p\ge\tfrac12$, $v\ge0$, follows from
its cosine integral against a beta probability measure (with the
endpoint $p=\tfrac12$ giving a cosine). Thus $F$ and $F^{(N)}$ are
bounded here. The Heine integrals converge absolutely by the finite
moments, so all symmetrizations and integrations above are justified.
\end{proof}

\begin{lemma}[Yang kernel positivity]\label{ya:kernel-positive}
Let $q(t)=1/t-1/(2\sinh(t/2))$ and
$h_N(t)=P_N(D_t^2)q(t)$, with $P_N$ defined from \eqref{ya:mu}.
There is a probability measure $\tau_N$ on $(0,\infty)$ such that
\begin{equation}\label{ya:positive-kernel}
 h_N(t)=\frac{H_N}{24(2N+1)!}\,t^{2N+1}
        \int_{(0,\infty)}e^{-t^2/(4\lambda)}
                          \,\tau_N(d\lambda),\qquad t>0.
\end{equation}
In particular $h_N(t)>0$, and
$h_N(\sqrt z)/z^{N+1/2}$ is completely monotone in $z>0$.
\end{lemma}
\begin{proof}
Binet's digamma integral and the duplication formula give
\begin{equation}\label{ya:Stieltjes}
 R(x)=2\int_0^\infty
       \frac{t\,dt}{(x^2+t^2)(e^{2\pi t}+1)}
     =\frac1{24}\int_0^\infty\frac{d\mu(u)}{x^2+u}.
\end{equation}
Indeed, subtracting Binet's formula at $x$ from twice its value at
$2x$, and then using
$\psi(x+\tfrac12)=2\psi(2x)-\psi(x)-2\log2$, replaces
$(e^{2\pi t}-1)^{-1}$ by
$(e^{2\pi t}-1)^{-1}-2(e^{4\pi t}-1)^{-1}
=(e^{2\pi t}+1)^{-1}$.
The ordinary digamma integral also gives
$R(x)=\int_0^\infty e^{-xt}q(t)\,dt$.
Since the Laplace transform of
$\sin(t\sqrt u)/\sqrt u$ is $(x^2+u)^{-1}$,
uniqueness of Laplace transforms yields
\[
 q(t)=\frac1{24}\int
       \frac{\sin(t\sqrt u)}{\sqrt u}\,d\mu(u).
\]
Fubini is valid, for example by
$|\sin(t\sqrt u)/\sqrt u|\le t$; both sides are continuous.
Differentiation under the integral is justified to every finite order
by the moments of $\mu$ and gives
\[
 h_N(t)=\frac{(-1)^N}{24}\int\pi_N(u)
                   \frac{\sin(t\sqrt u)}{\sqrt u}\,d\mu(u).
\]

Apply Lemma~\ref{ya:matrix-marginal} with $\alpha=3/2$ and
\[
 f(u)=\frac{24u^{-1/2}}{e^{2\pi\sqrt u}+1}
 =12u^{-1/2}e^{-\pi\sqrt u}
        \prod_{j\ge1}\left(1+\frac{4u}{(2j-1)^2}\right)^{-1}.
\]
The last equality is Euler's product for $\cosh(\pi\sqrt u)$.
The reciprocal-square coefficients are summable, so this function
satisfies Lemma~\ref{ya:spectral-class}.
Thus, for $r=N+1$, the random variable $S$ in
Lemma~\ref{ya:divided-difference} is a mixture of gamma laws with
shape $p=N+3/2$ and positive rate distribution $\tau_N$.
Equations \eqref{ya:divided-difference-identity} and
\eqref{ya:bessel-derivative} imply
\[
 \frac{h_N(t)}{t^{2N+1}}
 =\frac{H_N}{24\,4^NN!(3/2)_N}
       \mathbb E\left[{}_0F_1\left(p;-\frac{t^2S}{4}\right)\right].
\]
For each fixed $\lambda>0$, integration against the gamma density
of shape $p$ and rate $\lambda$ gives
\[
 \int_0^\infty {}_0F_1(p;-zs/4)
       \frac{\lambda^p}{\Gamma(p)}s^{p-1}e^{-\lambda s}\,ds
       =e^{-z/(4\lambda)},\qquad z\ge0.
\]
This identity follows by termwise integration of the hypergeometric
series: the series of integrals of absolute values is
$\exp(z/(4\lambda))<\infty$ for the fixed $\lambda$.
The subsequent integration over $\tau_N$ is justified by the bound
$|{}_0F_1(p;-v)|\le1$; no exponential moment of $\tau_N$ is used.
Finally $4^NN!(3/2)_N=(2N+1)!$, proving
\eqref{ya:positive-kernel}. Strict positivity follows because
$\tau_N$ is a probability measure on positive finite rates.
\end{proof}

\subsection{Coefficients, complete monotonicity, and uniqueness}

\begin{proof}[Proof of Theorem~\ref{ya:main}]
The measure \eqref{ya:mu} is positive on $(0,\infty)$ and has all
moments. Substituting $v=2\pi\sqrt u$ and expanding
$(e^v+1)^{-1}$ gives
\[
 m_j=\frac{48\Gamma(2j+2)}{(2\pi)^{2j+2}}
          (1-2^{-2j-1})\zeta(2j+2),
\]
which is \eqref{ya:moments}, including $m_0=1$.
All Hankel determinants $\Delta_N$ are positive by positive
definiteness of the moment form. The standard Gram determinant
construction gives \eqref{ya:determinant}; it may also be obtained
directly by solving the $N$ orthogonality equations.
Expanding the polynomial quotient in \eqref{ya:Q} proves
\eqref{ya:bcoefficients}.

The sine representation of $q$ gives
\[
 q^{(2j)}(0)=0,\qquad
 q^{(2j+1)}(0)=\frac{(-1)^jm_j}{24}.
\]
Its derivatives of every fixed order are bounded, by the finite
moments and $\int u^{-1/2}\,d\mu(u)<\infty$.
Repeated integration by parts therefore shows
\[
 \int_0^\infty e^{-xt}q^{(2j)}(t)\,dt
 =x^{2j}R(x)-\frac1{24}\sum_{\ell=0}^{j-1}
                (-1)^\ell m_\ell x^{2j-2-2\ell}.
\]
Combining these identities with the coefficients of $P_N$ yields
\begin{equation}\label{ya:laplace-E}
 E_N(x)=\int_0^\infty e^{-xt}h_N(t)\,dt.
\end{equation}
By Lemma~\ref{ya:kernel-positive},
\[
 0<h_N(t)\le\frac{H_N}{24(2N+1)!}t^{2N+1},\qquad t>0.
\]
Consequently all derivatives of \eqref{ya:laplace-E} can be taken
under the integral, and they have the asserted strict signs.
Also, dominated convergence in \eqref{ya:positive-kernel} gives
\[
 \lim_{t\downarrow0}\frac{h_N(t)}{t^{2N+1}}
        =\frac{H_N}{24(2N+1)!}.
\]
Changing variables $t=y/x$ in \eqref{ya:laplace-E} and applying the
same polynomial majorant proves \eqref{ya:asymptotic}.

For uniqueness, let $A$ and $B$ be real polynomials of degrees at
most $N$ and $N-1$, respectively, such that
$A(x^2)R(x)-B(x^2)$ has a finite nonzero limit after multiplication
by $x^{2N+2}$. Set $s=x^2$ and
\[
 Q_A(s)=\int\frac{A(s)-A(-u)}{s+u}\,d\mu(u).
\]
Polynomial division in \eqref{ya:Stieltjes} shows that the
polynomial part must vanish, so $B=Q_A/24$, and the remainder is
\[
 \frac1{24}\int\frac{A(-u)}{s+u}\,d\mu(u).
\]
Finite geometric division, with a remainder controlled by the
available moments, now forces
\[
 \int u^jA(-u)\,d\mu(u)=0,\qquad 0\le j<N.
\]
The space of polynomials of degree at most $N$ orthogonal to all
lower-degree polynomials is one-dimensional. Therefore
$A(s)=cP_N(s)$ and $B(s)=cQ_{N-1}(s)/24$ for a nonzero real $c$.
Complete monotonicity implies nonnegativity of the function; its
nonzero asymptotic coefficient is $cH_N/24$, so necessarily $c>0$.
Conversely, every such positive multiple satisfies all requirements
by the result just proved.
\end{proof}


\section{The diagonal Pad\'e conjectures of Li and Chen}
\label{sec:li-chen-pade}

\paragraph{Original problem.}
Write
\[
 F(x)=\frac{\Gamma(x+1)}{\sqrt{2\pi x}(x/e)^x},\qquad x>0.
\]
Li and Chen \cite{LiChen2017} define the diagonal approximant
\[
 R_k(x)=\frac{x^k+a_1x^{k-1}+\cdots+a_k}
                   {x^k+b_1x^{k-1}+\cdots+b_k}
 \quad\hbox{by}\quad F(x)-R_k(x)=O(x^{-2k-1}).
\]
Their equation (3.16) is the coefficient-matching system for this
condition. Their Conjecture 4.1 asserts $F(x)<R_k(x)$ for odd $k$
and $F(x)>R_k(x)$ for even $k$, throughout $x\ge1$.
Their Conjecture 3.1 asserts $a_j=(-1)^jb_j$.
The theorem below proves both original statements. The coefficient
symmetry is classical: McCabe's discussion of reciprocal series
\cite[pp.~67--68]{McCabe1989} explicitly includes the scaled gamma
series and credits N\'emeth--Zim\'anyi \cite{NemethZimanyi1982}.
It is included for completeness, not as an additional new closure.

\begin{theorem}\label{lc:main}
For every integer $k\ge1$, the normalized coefficient system (3.16)
of Li and Chen has a unique solution. Its numerator and denominator
are relatively prime, have degree exactly $k$, and satisfy
$a_j=(-1)^jb_j$ for $1\le j\le k$.
The denominator is positive for every $x>1/24$, and on that interval
\[
 (-1)^{k+1}\bigl(R_k(x)-F(x)\bigr)>0.
\]
In particular, Conjecture 4.1 holds with its original domain $x\ge1$.
\end{theorem}

The proof uses Binet's integral, the Herglotz representation theorem,
and positive Gaussian and Gauss--Radau quadrature. These are classical
ingredients. The crucial function to approximate is a M\"obius
transform of $F$, rather than $F$ itself.

\begin{lemma}[Classical Cauer--Herglotz representation, with moment extraction]
\label{lc:positive-real}
Suppose $h$ is holomorphic in $\operatorname{Re}z>0$,
$\operatorname{Re}h(z)\ge0$, and
$h(\overline z)=\overline{h(z)}$.
If $\lim_{x\to\infty}xh(x)=M<\infty$, then a finite positive measure
$\mu$ on $[0,\infty)$ of mass $M$ satisfies
\[
 h(z)=\int_{[0,\infty)}\frac{z}{z^2+u}\,d\mu(u),
 \qquad \operatorname{Re}z>0.
\]
If, in addition, $h(x)$ has an asymptotic expansion to every order in
odd inverse powers of $x$, then $\mu$ has moments of every order and
\[
 h(x)\sim\sum_{j\ge0}\frac{(-1)^jm_j}{x^{2j+1}},
 \qquad m_j=\int u^j\,d\mu(u).
\]
\end{lemma}
\begin{proof}
The representation step below is classical; Cauer's positive-real
form appears, for example, as equation (71) in the arXiv version
1103.5640 of Kalugin, Jeffrey, and Corless \cite{Kalugin2011}.
The representation is due to Cauer \cite{Cauer1932}. We spell out its reduction from Herglotz's
theorem to fix the symmetry and normalization conventions.
The function $K(w)=ih(-iw)$ is a Herglotz function on the upper
half-plane and obeys $K(-\overline w)=-\overline{K(w)}$.
In its unique Herglotz representation
\[
 K(w)=a+bw+\int_{\mathbb R}
 \left(\frac1{t-w}-\frac{t}{1+t^2}\right)d\sigma(t),
 \quad a\in\mathbb R,\ b\ge0,\quad
 \int\frac{d\sigma(t)}{1+t^2}<\infty,
\]
this symmetry forces $a=0$ and $\sigma$ to be even. Pair the terms
at $t$ and $-t$, and push forward $\sigma$ by $t\mapsto t^2$.
The result is
\[
 h(z)=bz+\int_{[0,\infty)}\frac{z}{z^2+u}\,d\mu(u),
 \qquad \int\frac{d\mu(u)}{1+u}<\infty.
\]
For real $x>0$ all terms are nonnegative. Consequently
\[
 xh(x)=bx^2+\int\frac{x^2}{x^2+u}\,d\mu(u)
\]
and monotone convergence show that $b=0$ and $\mu([0,\infty))=M$.

Set $S(s)=h(\sqrt s)/\sqrt s$. Its all-order expansion has the form
$S(s)\sim\sum_{j\ge0}d_js^{-j-1}$. The mass identity gives $m_0=d_0$.
Inductively, if $m_0,\ldots,m_{n-1}$ are finite and
$d_j=(-1)^jm_j$ for $j<n$, finite geometric division gives
\[
 (-1)^ns^{n+1}\left(S(s)-
       \sum_{j=0}^{n-1}\frac{(-1)^jm_j}{s^{j+1}}\right)
   =\int\frac{s u^n}{s+u}\,d\mu(u).
\]
The asymptotic expansion gives the finite limit $(-1)^nd_n$ on the
left; monotone convergence gives $m_n$ on the right. This proves
the moment assertion inductively, without assuming convergence of
any infinite asymptotic series.
\end{proof}

\begin{lemma}[The measure for the scaled gamma function]
\label{lc:measure}
There is a positive measure $\mu$ on $(0,\infty)$, of mass $1/24$,
with infinite support and moments of every order, such that
\[
 H(x):=\frac{F(x)-1}{F(x)+1}
      =\int\frac{x}{x^2+u}\,d\mu(u),\qquad x>0.
\]
\end{lemma}
\begin{proof}
Use the holomorphic logarithm of $\Gamma$ and the principal logarithm
of $z$ in the right half-plane, and put
\[
 g(z)=\log\Gamma(z+1)-(z+\tfrac12)\log z+z-\tfrac12\log(2\pi).
\]
Binet's second formula, followed by integration by parts, gives
\[
 g(z)=-\frac1\pi\int_0^\infty
          \log(1-e^{-2\pi t})\frac{z}{z^2+t^2}\,dt,
 \qquad \operatorname{Re}z>0.
\]
The identity holds first on the positive real axis and then by
holomorphic continuation; the displayed integral converges locally
uniformly in the right half-plane. Since
\[
 \frac{z}{z^2+t^2}=\frac12\left(\frac1{z-it}+\frac1{z+it}\right),
\]
its real part is strictly positive there. Thus $\operatorname{Re}g>0$.
The function $h(z)=\tanh(g(z)/2)$ is holomorphic in that half-plane:
the poles of $\tanh$ have zero real part, whereas
$\operatorname{Re}(g/2)>0$. Moreover,
\[
 \operatorname{Re}h(z)=
 \frac{\sinh(\operatorname{Re}g(z))}
      {\cosh(\operatorname{Re}g(z))+\cos(\operatorname{Im}g(z))}>0.
\]
It is real-symmetric, and on the positive real axis it equals $H$.
Stirling's all-order expansion
\[
 g(x)\sim\sum_{m\ge1}
       \frac{B_{2m}}{2m(2m-1)x^{2m-1}}
\]
and the odd analytic Taylor expansion of $\tanh$ show that $h$ has
an all-order expansion in odd inverse powers, beginning with
$1/(24x)$. Lemma~\ref{lc:positive-real} gives $\mu$, its mass, and
all its moments.

Finally $F(x)\sim(2\pi x)^{-1/2}$ as $x\downarrow0$, so
$H(x)\to1$. The measure has no atom at zero, since such an atom
would force $H(x)\ge\mu(\{0\})/x$. If its support were finite and
contained only positive points, its representation would instead
give $H(x)\to0$. Hence the support is infinite.
\end{proof}

\begin{proof}[Proof of Theorem~\ref{lc:main}]
Let $\mu$ be the measure in Lemma~\ref{lc:measure}, and write
$S(s)=\int(s+u)^{-1}\,d\mu(u)$. All quadratures below exist because
the moment form is strictly positive on every nonzero polynomial.

For $k=2n$, $n\ge1$, take the $n$-point Gaussian quadrature for
$\mu$, with nodes $u_i>0$ and weights $w_i>0$. It is exact through
degree $2n-1$, and $\sum_iw_i=1/24$. Put
\[
 S^-_n(s)=\sum_{i=1}^n\frac{w_i}{s+u_i},
 \qquad h_{2n}(x)=xS^-_n(x^2).
\]
Hermite interpolation of $u\mapsto(s+u)^{-1}$ at the Gaussian
nodes, with both values and first derivatives prescribed, gives
\[
 S^-_n(s)<S(s),\qquad s>0.
\]
Indeed the pointwise interpolation remainder has the positive sign
of $\prod_i(u-u_i)^2$, since its derivative of order $2n$ is
$(2n)!/(s+u)^{2n+1}>0$. The inequality is strict after integration
because the support of $\mu$ is infinite. Exactness and finite
geometric expansion also give
\[
 S(s)-S^-_n(s)=O(s^{-2n-1}),\qquad
 H(x)-h_{2n}(x)=O(x^{-4n-1}).
\]

For $k=2n+1$, $n\ge0$, take Gauss--Radau quadrature with the
prescribed node $u_0=0$ and $n$ positive nodes. Its positive weights
$w_0,\ldots,w_n$ sum to $1/24$, and it is exact through degree
$2n$. The positive nodes are the zeros of the degree-$n$ monic
orthogonal polynomial for $u\,d\mu(u)$. Set
\[
 S^+_n(s)=\frac{w_0}{s}+\sum_{i=1}^n\frac{w_i}{s+u_i},
 \qquad h_{2n+1}(x)=xS^+_n(x^2).
\]
For completeness, positivity of the weights follows by applying
exactness to the squares of the cardinal polynomials of degree
$n$; their integrals are strictly positive. Existence and exactness
follow by dividing a polynomial of degree at most $2n$ by
$u\pi_n(u)$, where $\pi_n$ is orthogonal for $u\,d\mu$.
Hermite interpolation now prescribes the value at zero and the
values and derivatives at the other nodes. Its remainder has the
negative sign of
$u\prod_{i=1}^n(u-u_i)^2$, since the derivative of order $2n+1$
of $(s+u)^{-1}$ is negative. Thus
\[
 S(s)<S^+_n(s),\qquad
 S(s)-S^+_n(s)=O(s^{-2n-2}),\qquad
 H(x)-h_{2n+1}(x)=O(x^{-4n-3}).
\]
These statements include $n=0$, when $S^+_0(s)=1/(24s)$.

Every discrete transform obeys
\[
 0<h_k(x)\le\frac1{24x}.
\]
Therefore $1-h_k(x)>0$ for $x>1/24$. On this interval define
\[
 \widetilde R_k(x)=\frac{1+h_k(x)}{1-h_k(x)}.
\]
The function $y\mapsto(1+y)/(1-y)$ is strictly increasing for
$y<1$, and $0<H(x)<1$ for positive $x$. The preceding strict
quadrature inequalities imply the claimed parity bounds for
$\widetilde R_k$.

It remains to verify exactly that these are the approximants in
the original question. When $k=2n$, the reduced expression for
$h_k$ has the form $A(x)/B(x)$ with $B$ monic and even of degree
$2n$, and $A$ odd of degree $2n-1$. When $k=2n+1$, it has $B$
monic and odd of degree $2n+1$, and $A$ even of degree $2n$.
There are no cancellations, because every pole of the discrete
Stieltjes transform has a strictly positive weight, including the
pole at zero in the Radau case. Hence
\[
 \widetilde R_k(x)=\frac{B(x)+A(x)}{B(x)-A(x)}
\]
is reduced and has monic numerator and denominator of degree $k$.
The displayed error estimates imply
$F(x)-\widetilde R_k(x)=O(x^{-2k-1})$.
Thus its coefficients satisfy precisely (3.16).

Two rational functions with numerator and denominator degrees at
most $k$, normalized to tend to one at infinity, and satisfying
this approximation condition must agree: their cross-multiplied
numerator has degree at most $2k$ but is $O(x^{-1})$ at infinity.
Since the constructed numerator and denominator are coprime and
both have degree $k$, their monic coefficient pair is unique.
Consequently $\widetilde R_k=R_k$ exactly.

The parities of $A$ and $B$ show
$[B+A](-x)=(-1)^k[B-A](x)$, proving
$a_j=(-1)^jb_j$. Finally $B(x)>0$ for $x>0$, since it is a product
of factors $x^2+u_i$, with one factor $x$ in the odd case.
Thus $B(x)-A(x)=B(x)(1-h_k(x))>0$ for $x>1/24$, as claimed.
\end{proof}

\paragraph{Priority boundary.}
The source and its original conjectures have been verified. A
bounded search did not locate an earlier solution of Conjecture 4.1
or this specific Stieltjes representation of $(F-1)/(F+1)$.
Luke's 1970 paper \cite{Luke1970}, N\'emeth--Zim\'anyi
\cite{NemethZimanyi1982}, and the unpreviewed pages of McCabe
\cite{McCabe1989} remain relevant unresolved prior-art leads.
Their complete texts could not be inspected, so an equivalent earlier
inequality theorem cannot be excluded. The inspected McCabe preview
establishes the classical coefficient symmetry, not the alternating
inequalities. This is not a certification of historical priority. The classical
Herglotz and quadrature theorems are not claimed as new.

\section{An interior counterexample to the Ferreira--Simon Wright conjecture}
\label{sec:wright-counterexample}

For $\rho>-1$, define the Wright function by
\[
 \phi(\rho,\beta,z)
 =\sum_{n=0}^{\infty}\frac{z^n}{n!\,\Gamma(\beta+\rho n)}.
\]
The reciprocal Gamma function is understood at its zeros, so that this
definition also covers nonpositive integral values of $\beta+\rho n$.
For $0<\alpha<1$ and $\beta\geq0$, the function
\[
 \varphi_{\alpha,\beta}(x)
 =\Gamma(\alpha+\beta)\phi(-\alpha,\beta,-x),\qquad x>0,
\]
is the probability density considered by Ferreira and Simon
\cite{FerreiraSimon2024}.

For $0<\beta<1$, let $\alpha_*(\beta)$ denote the unique solution in
$(0,1)$ of
\begin{equation}
 \frac{1}{\Gamma(\beta-\alpha)^2}
 =\frac{1}{\Gamma(\beta)\Gamma(\beta-2\alpha)}.
 \label{eq:wright-threshold-definition}
\end{equation}
Set $\alpha_*(0)=1$, as in the original paper. Section~6.4 of
\cite{FerreiraSimon2024} conjectures the equivalence between
$\alpha\leq\alpha_*(\beta)$, log-concavity of this density, and convexity
on $\mathbb R_+$ of the reciprocal Mittag--Leffler function
$1/E_{\alpha,\alpha+\beta}(-t)$. In particular, the conjecture asserts
\begin{equation}
 \alpha\leq\alpha_*(\beta)
 \quad\Longrightarrow\quad
 \varphi_{\alpha,\beta}\text{ is log-concave}.
 \label{eq:wright-original-implication}
\end{equation}
We disprove this implication within its original parameter domain.

The threshold test established in the same source will be used only in
the region $0<\beta<2\alpha-1<1$, where
\begin{equation}
 \alpha>\alpha_*(\beta)
 \quad\Longleftrightarrow\quad
 R(\alpha,\beta):=
 \frac{\Gamma(\beta)\Gamma(\beta-2\alpha)}
      {\Gamma(\beta-\alpha)^2}<1.
 \label{eq:wright-threshold-test}
\end{equation}
Consequently $R(\alpha,\beta)>1$ in this region implies
$\alpha<\alpha_*(\beta)$: the strict reverse side is excluded by
\eqref{eq:wright-threshold-test}, and equality is excluded by
\eqref{eq:wright-threshold-definition}.

\subsection{A limiting profile at the boundary}

\begin{theorem}
\label{thm:wright-boundary}
For all sufficiently small $\varepsilon>0$, the density
$\varphi_{1-\varepsilon,0}$ is strictly logarithmically convex at
$x=1/2$. More generally, if $K$ is a compact subinterval of
$(\sqrt2-1,1)$, then
\[
 \bigl(\log\varphi_{1-\varepsilon,0}\bigr)''(x)>0
 \qquad(x\in K)
\]
for every sufficiently small $\varepsilon>0$. Thus
\eqref{eq:wright-original-implication} is false.
\end{theorem}

\begin{proof}
Write $f_\varepsilon(x)=\phi(-1+\varepsilon,0,-x)$. Euler reflection,
including its continuous interpretation at Gamma poles, yields
\begin{equation}
 f_\varepsilon(x)
 =\frac1\pi\sum_{n=1}^{\infty}
   \frac{\Gamma(1+(1-\varepsilon)n)}{n!}
   \sin(\pi\varepsilon n)x^n.
 \label{eq:wright-reflected-series}
\end{equation}
Indeed,
$1/\Gamma(-u)=-\Gamma(1+u)\sin(\pi u)/\pi$ and
$\sin(\pi(1-\varepsilon)n)=(-1)^{n+1}\sin(\pi\varepsilon n)$.
The $n=0$ coefficient vanishes because $1/\Gamma(0)=0$.

Log-convexity of Gamma on $[1,n+1]$ gives
\[
 0<\Gamma(1+(1-\varepsilon)n)\leq n!.
\]
The coefficient of $x^n$ in $f_\varepsilon/\varepsilon$ therefore has
absolute value at most $n$, and it converges to $n$ as
$\varepsilon\downarrow0$. On every interval $|x|\leq r<1$, this bound
also gives a summable majorant for any fixed number of derivative
series. Thus, for every fixed nonnegative integer $k$,
\begin{equation}
 \frac{f_\varepsilon(x)}{\varepsilon}
 \longrightarrow h(x):=\sum_{n=1}^{\infty}nx^n
 =\frac{x}{(1-x)^2}
 \quad\text{in }C^k\text{ on compact subintervals of }(-1,1).
 \label{eq:wright-profile-limit}
\end{equation}
In particular,
\begin{equation}
 \frac{(f_\varepsilon')^2-f_\varepsilon f_\varepsilon''}
      {\varepsilon^2}
 \longrightarrow
 (h')^2-hh''
 =\frac{1-2x-x^2}{(1-x)^6}
 \label{eq:wright-profile-curvature}
\end{equation}
locally uniformly. The last expression equals $-16$ at $x=1/2$ and
is strictly negative for $\sqrt2-1<x<1$. The limiting profile is
strictly positive there, so \eqref{eq:wright-profile-limit} also
ensures positivity of $f_\varepsilon$ on the compact intervals under
consideration. Since
\[
 (\log f_\varepsilon)''
 =-\frac{(f_\varepsilon')^2-f_\varepsilon f_\varepsilon''}
          {f_\varepsilon^2},
\]
the claimed strict logarithmic convexity follows. Multiplication by
$\Gamma(1-\varepsilon)>0$ does not change logarithmic curvature.
Finally, $1-\varepsilon<1=\alpha_*(0)$, proving the contradiction to
\eqref{eq:wright-original-implication}.
\end{proof}

\subsection{A completely explicit interior witness}

The preceding limiting argument already disproves the conjecture.
The next theorem supplies rational parameters with $\beta>0$, using
only inequalities and exact rational arithmetic.

\begin{theorem}
\label{thm:wright-interior}
Let
\[
 \varepsilon=\frac1{10000},\qquad
 \alpha=\frac{9999}{10000},\qquad
 \beta=\frac1{100000000}=\varepsilon^2.
\]
Then $\alpha<\alpha_*(\beta)$, but
\[
 \bigl(\log\varphi_{\alpha,\beta}\bigr)''(1/2)>0.
\]
In particular, these parameters give a counterexample strictly inside
the original admissible region.
\end{theorem}

\begin{proof}
We first record the elementary Gamma bounds used below. Let
$\psi=\Gamma'/\Gamma$ and let $\gamma$ be Euler's constant.
Monotonicity of $\psi$, together with
$\psi(1)=-\gamma$, $0<\gamma<1$, and
$\psi(n+1)=H_n-\gamma<n$, gives
\begin{equation}
 -1<\psi(t)<n\qquad(1\leq t\leq n+1,\ n\geq1).
 \label{eq:wright-digamma-bound}
\end{equation}
Here $H_n=\sum_{j=1}^n1/j$. Also, for $0\leq u\leq1$,
\begin{equation}
 \Gamma(1+u)\leq1,
 \qquad
 \Gamma(1+u)\geq e^{-u}\geq1-u.
 \label{eq:wright-gamma-unit-bounds}
\end{equation}
The upper bound follows from log-convexity and
$\Gamma(1)=\Gamma(2)=1$; the lower bound follows by integrating
$\psi\geq-1$.

Put
\[
 f_{\alpha,b}(x)=\phi(-\alpha,b,-x),\qquad
 g_0(x)=\frac{f_{\alpha,0}(x)}{\varepsilon},\qquad
 g(x)=\frac{f_{\alpha,\beta}(x)}{\varepsilon}.
\]
By \eqref{eq:wright-reflected-series},
\[
 g_0(x)=\sum_{n=1}^{\infty}A_nx^n,\qquad
 A_n=R_n\frac{\sin(\pi\varepsilon n)}{\pi\varepsilon},\qquad
 R_n=\frac{\Gamma(n+1-\varepsilon n)}{n!}.
\]
Log-convexity gives $0<R_n\leq1$. Furthermore,
\[
 \log R_n
 =-\int_{n+1-\varepsilon n}^{n+1}\psi(t)\,dt
 \geq-\varepsilon n^2
\]
by \eqref{eq:wright-digamma-bound}, so
$1-R_n\leq\varepsilon n^2$. The inequality
$|\sin z-z|\leq|z|^3/6$ holds for every real $z$. Hence
\begin{align}
 |A_n-n|
 &\leq n(1-R_n)
   +nR_n\left|\frac{\sin(\pi\varepsilon n)}
                         {\pi\varepsilon n}-1\right| \notag\\
 &\leq\varepsilon n^3+\frac{\pi^2}{6}\varepsilon^2n^3
 \leq3\varepsilon n^3.
 \label{eq:wright-coefficient-error}
\end{align}
The last inequality uses $\pi^2/6<2$ and $0<\varepsilon<1$.

At $x=1/2$, summing \eqref{eq:wright-coefficient-error} and its first
two derivative majorants gives
\begin{equation}
 |g_0-h|\leq78\varepsilon,\qquad
 |g_0'-h'|\leq900\varepsilon,\qquad
 |g_0''-h''|\leq11184\varepsilon.
 \label{eq:wright-boundary-error}
\end{equation}
For completeness, the geometric-power sums used here are
\[
 \sum_{n\geq1}\frac{n^3}{2^n}=26,\qquad
 \sum_{n\geq1}\frac{n^4}{2^n}=150,\qquad
 \sum_{n\geq1}\frac{n^5}{2^n}=1082.
\]
They follow by repeatedly applying $x\,d/dx$ to $x/(1-x)$.
The second-derivative constant in
\eqref{eq:wright-boundary-error} is
$12(1082-150)=11184$.

We next control the change in $b$. Write
\[
 c_n(b)=\frac{(-1)^n}{n!\,\Gamma(b-\alpha n)}.
\]
For $n\geq1$ and $0\leq b\leq\alpha/2$, reflection gives
\[
 c_n(b)=\frac{(-1)^n\Gamma(1+\alpha n-b)
                     \sin(\pi(b-\alpha n))}{\pi n!}.
\]
The Gamma argument lies in $[1,n+1]$, where its ratio to $n!$ is at
most one and $|\psi|\leq n$. Differentiation with respect to $b$
therefore yields
\begin{equation}
 |c_n'(b)|\leq n/\pi+1\leq n+1\qquad(n\geq1).
 \label{eq:wright-beta-coefficients}
\end{equation}
For $n=0$, recurrence gives
\[
 c_0(b)=\frac{b}{\Gamma(1+b)},\qquad
 c_0'(b)=\frac{1-b\psi(1+b)}{\Gamma(1+b)}.
\]
When $0\leq b\leq1/2$, \eqref{eq:wright-digamma-bound} with $n=1$
shows that the numerator is positive and at most $1+b$.
By \eqref{eq:wright-gamma-unit-bounds},
$|c_0'(b)|\leq(1+b)/(1-b)\leq3$.

Summing these coefficient bounds at $x=1/2$ gives
\begin{align}
 |f_{\alpha,\beta}-f_{\alpha,0}|&\leq6\beta,\notag\\
 |f_{\alpha,\beta}'-f_{\alpha,0}'|&\leq16\beta,\notag\\
 |f_{\alpha,\beta}''-f_{\alpha,0}''|&\leq96\beta.
 \label{eq:wright-beta-error}
\end{align}
Indeed,
\[
 3+\sum_{n\geq1}\frac{n+1}{2^n}=6,\qquad
 2\sum_{n\geq1}\frac{n(n+1)}{2^n}=16,\qquad
 4\sum_{n\geq1}\frac{n^3-n}{2^n}=96.
\]
The convergent majorants justify all differentiations and summations.

Combining \eqref{eq:wright-boundary-error} and
\eqref{eq:wright-beta-error}, using $\beta/\varepsilon=\varepsilon$,
we obtain at $x=1/2$
\[
 |g-h|\leq E_0:=\frac{21}{2500},\qquad
 |g'-h'|\leq E_1:=\frac{229}{2500},\qquad
 |g''-h''|\leq E_2:=\frac{141}{125}.
\]
Since $(h,h',h'')=(2,12,80)$ at this point, direct expansion of
$D_g=(g')^2-gg''$ gives
\begin{align*}
 |D_g+16|
 &\leq24E_1+E_1^2+2E_2+80E_0+E_0E_2\\
 &=\frac{32151661}{6250000}<6.
\end{align*}
Thus $D_g<-10$, while $g\geq2-E_0>0$. This proves
$(\log g)''(1/2)>0$ and hence the same strict inequality for
$\varphi_{\alpha,\beta}$.

It remains to verify the side of the conjectured threshold. Put
$a=\varepsilon+\beta$ and $b=2\varepsilon+\beta$. Gamma recurrence
gives the exact identity
\begin{equation}
 R(\alpha,\beta)
 =\frac{\Gamma(1+\beta)\Gamma(1+b)}{\Gamma(1+a)^2}
   \frac{a^2(1-a)^2}{\beta b(1-b)(2-b)}.
 \label{eq:wright-positive-ratio}
\end{equation}
Here $a,b<1/2$. The first factor is at least $1/4$ by
\eqref{eq:wright-gamma-unit-bounds}. Also,
$a\geq\varepsilon$, $1-a>1/2$, $b\leq3\varepsilon$,
$1-b<1$, and $2-b<2$, with every denominator factor positive.
Consequently
\[
 R(\alpha,\beta)\geq\frac{\varepsilon}{96\beta}
 =\frac{625}{6}>100>1.
\]
The parameters satisfy $0<\beta<2\alpha-1<1$, so
\eqref{eq:wright-threshold-test} and
\eqref{eq:wright-threshold-definition} imply
$\alpha<\alpha_*(\beta)$, as required.
\end{proof}

Both inequalities in Theorem~\ref{thm:wright-interior} are strict.
The derivative-series majorants above also establish joint continuity
in $(\alpha,\beta)$ near this pair. The Gamma ratio is continuous
there as well. Hence the counterexamples include a nonempty open
subset of the two-dimensional parameter domain, rather than just the
boundary family of Theorem~\ref{thm:wright-boundary}.

\subsection{Scope and priority boundary}

The result is a negative answer to the unchanged universal implication
\eqref{eq:wright-original-implication}, and therefore to the proposed
three-way equivalence. It does not classify the full log-concavity
region and does not separately decide equivalence between the
log-concavity and reciprocal-convexity properties. The boundary
family, the explicit interior pair, and its open neighborhood address
one conjecture and are counted as one result.

The argument uses classical reflection, Gamma inequalities,
coefficient estimates, and a curvature determinant. No historically
new general principle is claimed. A bounded literature check through
3 November 2026 found no identical resolution; this included the
original paper, author publication lists, and the later preprint
\emph{Mittag--Leffler functions and convex ordering}
(arXiv:2512.04940v1). The absence of a result in that search does not
prove priority. The explicit proof above is intended to make the
mathematical assertion independently checkable while keeping that
historical limitation separate.


\section{A negative answer to the Yano--Yano--Yor negative-power question}
\label{sec:yyy}

For $\alpha>1$, let $C_\alpha$ have the symmetric $\alpha$-Cauchy density
\begin{equation}
 f_{C_\alpha}(x)=\frac{\alpha\sin(\pi/\alpha)}{2\pi}
                  \frac{1}{1+|x|^\alpha},\qquad x\in\mathbb R.
 \label{eq:yyy-density}
\end{equation}
Yano, Yano and Yor~\cite[Remark~2.9(iii)]{YanoYanoYor2009}
asked whether $|C_\alpha|^{-p}$ is self-decomposable, or at least
infinitely divisible, for $p>0$. We answer this universal question
negatively. The counterexample lies in $1<\alpha<2$, so it also applies
if that smaller range motivated the question.

\begin{theorem}
\label{thm:yyy-explicit}
The random variable
\[
 Y=|C_{4/3}|^{-1/12}
\]
is not infinitely divisible, and hence is not self-decomposable.
Its density and an independent-gamma representation are
\begin{equation}
 f_Y(y)=\frac{8\sqrt2}{\pi}\frac{y^3}{1+y^{16}},\qquad y>0,
 \qquad
 Y\ \stackrel{d}{=}\
 \left(\frac{G_{1/4}}{G_{3/4}}\right)^{1/16},
 \label{eq:yyy-gb2}
\end{equation}
where $G_{1/4}$ and $G_{3/4}$ are independent gamma random variables
with the indicated shapes and unit scale.
\end{theorem}

We give a proof using only three moments. In particular, no numerical
approximation is needed to certify the counterexample.

\begin{lemma}
\label{lem:yyy-third-cumulant}
If a nonnegative infinitely divisible random variable $V$ has finite
third moment, then
\[
 \kappa_3(V):=\E[V^3]-3\E[V]\E[V^2]+2\E[V]^3\ge0.
\]
\end{lemma}

\begin{proof}
For each positive integer $n$, write $V$ in distribution as a sum of
$n$ independent copies of a convolution root $Z_n$. The essential lower
endpoint of the support of such a sum is $n$ times the essential lower
endpoint for a summand. An unbounded-below root would give an
unbounded-below sum. Since $V$ is nonnegative, every $Z_n$ is therefore
nonnegative. Its third moment is finite, because each nonnegative
summand is bounded by their sum.

Put $\mu=\E[V]$ and $v=\Var(V)$. Independence gives
\[
 \E[Z_n]=\frac{\mu}{n},\qquad
 \E[Z_n^2]=\frac{v}{n}+\frac{\mu^2}{n^2}.
\]
Additivity of the third cumulant and $\E[Z_n^3]\ge0$ imply
\[
 \kappa_3(V)=n\kappa_3(Z_n)
 \ge -\frac{3\mu v}{n}-\frac{\mu^3}{n^2}.
\]
Let $n$ tend to infinity.
\end{proof}

\begin{lemma}
\label{lem:yyy-factorization}
Let $1<\alpha<2$ and $0<p<1/3$, and put
\[
 \theta=\frac{\pi}{\alpha},\qquad
 h=\frac{\pi p}{\alpha},\qquad
 a=\cot\theta,\qquad t=\tan h.
\]
Define
\begin{align*}
 B_1&=1-at,\\
 B_2&=1-t^2-2at,\\
 B_3&=1-3t^2-at(3-t^2).
\end{align*}
These three numbers are strictly positive, and
\begin{equation}
 \kappa_3(|C_\alpha|^{-p})
 =\frac{2t^3(1+a^2)(1+t^2)^{3/2}
       \bigl[a+(a^2+3)t-at^2\bigr]}
       {B_1^3B_2B_3}.
 \label{eq:yyy-cumulant-factor}
\end{equation}
\end{lemma}

\begin{proof}
For $j=1,2,3$, the condition $jp<1$ and the beta integral give
\begin{align}
 m_j:=\E[|C_\alpha|^{-jp}]
 &=\frac{\alpha\sin\theta}{\pi}
     \int_0^\infty\frac{x^{-jp}}{1+x^\alpha}\,dx
 =\frac{\sin\theta}{\sin(\theta-jh)}.
 \label{eq:yyy-moments}
\end{align}
Indeed, after substituting $u=x^\alpha$, the beta-integral parameter is
$(1-jp)/\alpha\in(0,1)$. Also
$0<h<\theta/3<\pi/3$, so $\cos h>0$. The angle-subtraction identities
show that
\[
 B_j=\frac{\sin(\theta-jh)}{\sin\theta\cos^j h}>0,
 \qquad
 m_j=\frac{(1+t^2)^{j/2}}{B_j}
 \quad(j=1,2,3).
\]
Consequently
\[
 m_3-3m_1m_2+2m_1^3
 =\frac{(1+t^2)^{3/2}}{B_1^3B_2B_3}
       \bigl(B_1^3B_2-3B_1^2B_3+2B_2B_3\bigr).
\]
For completeness, the numerator identity needed here is
\begin{align*}
 B_1^3B_2-3B_1^2B_3+2B_2B_3
 &=2a(1+a^2)t^3
   +2(a^4+4a^2+3)t^4-2a(1+a^2)t^5\\
 &=2t^3(1+a^2)\bigl[a+(a^2+3)t-at^2\bigr].
\end{align*}
Substitution proves~\eqref{eq:yyy-cumulant-factor}.
\end{proof}

\begin{proof}[Proof of Theorem~\ref{thm:yyy-explicit}]
Take $\alpha=4/3$ and $p=1/12$. Then
$a=-1$, $h=\pi/16$, and the third moment exists because $3p=1/4<1$.
We first verify the exact bound
\[
 t=\tan(\pi/16)<\frac15.
\]
Since $2<289/144$, we have $\sqrt2<17/12$, and hence
\[
 \tan(\pi/8)=\sqrt2-1<\frac5{12}
       =\tan\bigl(2\arctan(1/5)\bigr).
\]
Both arguments on the right and left lie in $(0,\pi/2)$.
Strict monotonicity of tangent proves the asserted bound. It follows that
\[
 a+(a^2+3)t-at^2=-1+4t+t^2
 <-1+\frac45+\frac1{25}=-\frac4{25}<0.
\]
All other factors in~\eqref{eq:yyy-cumulant-factor} are positive.
Thus $\kappa_3(Y)<0$, contradicting the necessary condition in
Lemma~\ref{lem:yyy-third-cumulant} for infinite divisibility.
Self-decomposability implies infinite divisibility, so it fails as well.

The change of variables $x=y^{-12}$ in the density of $|C_{4/3}|$
gives the density in~\eqref{eq:yyy-gb2}. The ratio of independent gamma
variables with shapes $1/4$ and $3/4$ has beta-prime density
$z^{-3/4}/\bigl(B(1/4,3/4)(1+z)\bigr)$. Transforming by $y=z^{1/16}$,
and using $B(1/4,3/4)=\pi\sqrt2$, gives the same density.
\end{proof}

The factorization gives more than a single witness.

\begin{theorem}
\label{thm:yyy-region}
For every $1<\alpha<2$, define
\[
 c=-\cot(\pi/\alpha)>0,\qquad
 t_0(c)=\frac{2c}
 {c^2+3+\sqrt{(c^2+1)(c^2+9)}}.
\]
Then $|C_\alpha|^{-p}$ is not infinitely divisible whenever
\begin{equation}
 0<p<\frac{\alpha}{\pi}\arctan t_0(c).
 \label{eq:yyy-region}
\end{equation}
\end{theorem}

\begin{proof}
With $a=-c$, the sign-controlling bracket in
\eqref{eq:yyy-cumulant-factor} is
\[
 ct^2+(c^2+3)t-c.
\]
Its positive root is $t_0(c)$, since
$(c^2+3)^2+4c^2=(c^2+1)(c^2+9)$; its other root is negative.
It is therefore strictly negative for $0<t<t_0(c)$.
Moreover,
\[
 t_0(c)<\frac{c}{c^2+3}
 \le\frac1{2\sqrt3}<\tan(\pi/6).
\]
Condition~\eqref{eq:yyy-region} implies
$h<\pi/6<\theta/3$, because $\theta=\pi/\alpha>\pi/2$.
In particular $p<1/3$, so Lemma~\ref{lem:yyy-factorization} applies.
Its third cumulant is strictly negative, and
Lemma~\ref{lem:yyy-third-cumulant} completes the proof.
\end{proof}

\paragraph{Relation to later work and the scope of the conclusion.}
Wang~\cite[Section~1.3]{Wang2024} repeats the negative-power question
and proves sufficient conditions for infinite divisibility. For negative
powers with $1<\alpha\le2$, the sufficient condition there is
$p\ge\alpha/2$, which does not include the counterexamples above.
Wang's later classification~\cite{Wang2026} concerns $C_\alpha$ itself:
it is infinitely divisible exactly for $1<\alpha\le2$. Its introduction
leaves the other Yano--Yano--Yor questions for further research; that
classification does not settle the negative powers considered here.

Bosch and Simon~\cite[final paragraph]{BoschSimon2015} also ask about
infinite divisibility of the generalized beta-prime family
$(G_b/G_a)^s$ for $s<1$, after giving a sufficient subregion.
The representation~\eqref{eq:yyy-gb2} places our example at
$(a,b,s)=(3/4,1/4,1/16)$. This is a closely related application of the
same witness, and is not counted as a second independent target.

Theorem~\ref{thm:yyy-explicit} is a complete negative answer to the
literal universal question in Remark~2.9(iii); it is not a classification
of all $(\alpha,p)$. The proof uses the classical obstruction supplied
by a negative cumulant. A bounded search of the original questions,
related beta-prime power literature, and the cited subsequent papers
found no identical resolution. This records the scope of the prior-art
check; it does not establish exhaustive coverage or priority.


\section{A rooted neighborhood lemma and Graffiti.pc Conjecture 177}
\label{sec:wow177}

For a finite simple connected graph $G$, let $L(G)$ be the maximum
number of leaves in a spanning tree, $b(G)$ the maximum order of an
induced bipartite subgraph, and $\alpha(G)$ the independence number.
Write $d_2(G)$ for the second entry of the list of all vertex degrees
in nondecreasing order, with repetitions. In particular, $d_2(G)$
equals the minimum degree when two vertices attain the minimum.
Graffiti.pc Conjecture~177 (2005) asserts that, whenever $|V(G)|\ge2$,
\begin{equation}\label{wow:target}
 L(G)+b(G)\ \ge\ 2\alpha(G)+d_2(G).
\end{equation}
The question appears in West's record of the Graffiti conjectures
\cite{wow:West}. DeLaVi\~na and Waller proved the lower bound
$2\alpha(G)+1$ and explicitly distinguished it from the open
second-degree bound \cite[Theorem~3 and the paragraph following it]{wow:DW}.

\paragraph{A rooted deletion problem.}
Suppose $I$ is a maximal independent set and $r\in I$ is prescribed.
Can one find a different vertex $u\in I$ and keep just one of its
neighbors so that deleting all its other neighbors leaves a connected
graph? The freedom to prescribe $r$ matters for \eqref{wow:target}.
It allows the construction to avoid a vertex of exceptionally small
degree. Write $N(u)$ for the open neighborhood and
$N[u]=N(u)\cup\{u\}$ for the closed neighborhood.

\begin{lemma}[Rooted neighborhood deletion]\label{wow:star}
Let $G$ be connected and let $I$ be a maximal independent set with
$|I|\ge2$. For each $r\in I$ there are $u\in I\setminus\{r\}$ and
$x\in N(u)$ such that $G-(N(u)\setminus\{x\})$ is connected.
\end{lemma}
\begin{proof}
For $v\in I\setminus\{r\}$, let $A(v)$ be the component of $G-N[v]$
containing $r$. This component exists because $I$ is independent.
Choose $u$ maximizing $|A(u)|$ and put $A=A(u)$. Some vertex
$x\in N(u)$ has a neighbor in $A$. Indeed, $G$ is connected, and every
neighbor of $A$ outside $A$ belongs to $N(u)$.

Let $B$ be any other component of $G-N[u]$. It contains a vertex
$v\in I$. To see this, take any $w\in B\setminus I$ if $B$ contains
such a vertex. Maximality of $I$ supplies an $I$-neighbor of $w$.
That neighbor cannot be $u$ because $w\notin N[u]$, and all other
vertices of $I$ are outside $N[u]$. It therefore lies in $B$. If
$B\subseteq I$ the assertion is immediate.

There are no edges from $v$ to $A$. If $v$ were also nonadjacent to
$x$, the connected set $A\cup\{x\}$ would lie in $G-N[v]$. It
contains $r$, so $|A(v)|\ge |A|+1$, contrary to the choice of $u$.
Thus $x$ has a neighbor in $B$. Vertex $x$ now joins every component
of $G-N[u]$ and is adjacent to $u$. These components together with
$x$ and $u$ comprise exactly the asserted connected graph.
\end{proof}

\begin{theorem}[Leaf bound relative to a maximal independent set]
\label{wow:relative}
Let $G$ be connected and let $I$ be a maximal independent set with
$|I|\ge2$. Put $C=V(G)\setminus I$ and $t=\alpha(G[C])$.
Let $d_2(I)$ be the second-smallest degree in $G$ among the vertices
of $I$, counted with multiplicity. Then
\begin{equation}\label{wow:relative-eq}
 L(G)\ \ge\ |I|+d_2(I)-t.
\end{equation}
\end{theorem}
\begin{proof}
Choose $r\in I$ of minimum degree among the vertices of $I$. Apply
Lemma~\ref{wow:star} to find $u\ne r$ and $x\in N(u)$. Write
$d=\deg_G(u)\ge d_2(I)$, $R=N(u)\setminus\{x\}$, and $C'=C\setminus R$.
Then $|R|=d-1$, $x\in C'$, and $H=G-R$ is connected. Since $I$ is
independent, every vertex of $I$ has a neighbor in $C'$ in $H$.

Let $m$ be the number of components of $G[C']$. Choose one vertex
from each component to see that $m\le t$. At most $m-1$ vertices
from $I$ suffice to connect $C'$ in $H$. Indeed, contract the
components of $G[C']$ to obtain a connected bipartite incidence graph
between the components and $I$. Starting from one component, select
an $I$-vertex incident with a reached and an unreached component,
and reach all components incident with it. While an unreached
component remains, connectedness supplies such a choice. Each step
adds at least one component. Let $J\subseteq I$ be the selected
vertices and $D_0=C'\cup J$. This set is connected, dominates $I$,
and satisfies
\[
 |D_0|\le |C'|+m-1=|C|-d+m.
\]

If every vertex of $R$ has a neighbor in $C'$, then $D_0$ is a
connected dominating set of $G$ of size at most $|C|-d+t$. Otherwise
take $y\in R$ with no neighbor in $C'$. Together with one vertex
from each component of $G[C']$, vertex $y$ forms an independent set
in $G[C]$. Thus $m+1\le t$. The set $D=D_0\cup\{u\}$ is connected
because $ux$ is an edge with $x\in C'$. It dominates every vertex
of $R$ because $R\subseteq N(u)$, and has size at most
$|C|-d+m+1\le |C|-d+t$.

In both cases there is a connected dominating set of size at most
$|C|-d+t$. Take a spanning tree of this set and attach each remaining
vertex by one edge to a neighbor in it. The resulting spanning tree
has at least $|I|+d-t$ leaves. Since $d\ge d_2(I)$, this gives
\eqref{wow:relative-eq}.
\end{proof}

\begin{corollary}[Graffiti.pc Conjecture~177]\label{wow:main}
Every finite simple connected graph with at least two vertices
satisfies \eqref{wow:target}.
\end{corollary}
\begin{proof}
For a graph with $\alpha(G)\ge2$ take a maximum independent set $I$.
Adjoining an independent set of $t=\alpha(G[C])$ vertices in its
complement gives an induced bipartite graph with parts of sizes
$|I|$ and $t$. Hence $b(G)\ge |I|+t$. Adding this to
\eqref{wow:relative-eq} gives
\[
 L(G)+b(G)\ge 2|I|+d_2(I)\ge 2\alpha(G)+d_2(G).
\]
The last inequality holds because taking a subset of the degree list
cannot decrease its second order statistic. If $\alpha(G)=1$, then
$G=K_n$. For $n\ge3$ a spanning star gives $L=n-1$, and
$b=2$, $d_2=n-1$, so equality holds. For $K_2$, both vertices are
leaves of its spanning tree, giving $L+b=4\ge3$.
\end{proof}

\paragraph{Sharpness.}
Let $G$ be the join of an independent set of $a\ge2$ vertices and a
clique of $c\ge1$ vertices. The independent vertices have degree $c$,
and the clique vertices are universal. Consequently
$\alpha(G)=a$, $d_2(G)=c$, and $L(G)=a+c-1$. An induced bipartite
subgraph can include all independent vertices and one clique vertex.
Including two clique vertices excludes every independent vertex by a
triangle, and three clique vertices already form a triangle. Thus
$b(G)=a+1$ and both sides of \eqref{wow:target} equal $2a+c$.


\section{Interior counterexamples to the reverse Mittag--Leffler inequality}
\label{ggps:section}

For $\alpha>0$ and $\beta>0$ put
\[
 E_{\alpha,\beta}(z)=\sum_{n=0}^{\infty}
                 \frac{z^n}{\Gamma(\alpha n+\beta)}.
\]
Garrappa, Gerhold, Popolizio and Simon \cite{GGPS2025}, in Conjecture~25, propose that for $1<\alpha<2$ the following three
conditions are equivalent: the global inequality
\begin{equation}\label{ggps:original}
 |E_{\alpha,\beta}(z)|\ge E_{\alpha,\beta}(\Re z)\qquad(z\in\mathbb C),
\end{equation}
the membership $\beta\in[\alpha-1,\alpha]$, and complete monotonicity
of $1/E_{\alpha,\beta}$ on $(0,\infty)$. Their Proposition~3 proves
complete monotonicity throughout this parameter interval. Their
Proposition~14 excludes the regions $0\le\beta<\alpha-1$ and
$\beta>\alpha$ from the global modulus inequality. Their Corollary~7
allows a compact exceptional set of arguments when
$\alpha-1<\beta\le\alpha$.

\begin{theorem}\label{ggps:main}
For every $1<\alpha<2$, set $C_\alpha=-1/\Gamma(-1-\alpha)>0$.
Fix $h>0$ satisfying $\arctan h<\pi(1-\alpha/2)$.
For all sufficiently large real $s$, uniformly for
\[
  \alpha-1\le\beta\le\alpha-1+\frac{C_\alpha}{2s},
\]
one has
\[
  |E_{\alpha,\beta}(-s+ihs)|<E_{\alpha,\beta}(-s).
\]
Consequently, for every $\alpha\in(1,2)$ there is a nonempty open
interval of $\beta$ strictly between $\alpha-1$ and $\alpha$ on which
\eqref{ggps:original} fails. This refutes Conjecture~25 at strictly
interior parameters where its reciprocal complete-monotonicity condition
holds.
\end{theorem}

The point of using oblique rays is that a positive leading term of
inverse-power order $m$ on the negative axis acquires the smaller
modulus factor $(1+h^2)^{-m/2}$. At $\beta=\alpha-1$ the usual
term of order $s^{-1}$ vanishes. Allowing $\beta-(\alpha-1)$ to vary
on the scale $s^{-1}$ retains both of the first two coefficients and
provides an interior range of counterexamples.

\begin{lemma}\label{ggps:expansion}
Fix $1<\alpha<2$. In any closed angular sector
$\phi\le |\arg z|\le\pi$ with $\phi>\pi\alpha/2$, one has
\[
 E_{\alpha,\beta}(z)=-\frac{z^{-1}}{\Gamma(\beta-\alpha)}
                      -\frac{z^{-2}}{\Gamma(\beta-2\alpha)}
                      +O(|z|^{-3})
\]
as $|z|\to\infty$. The bound is uniform for real $\beta$ in a fixed
compact interval.
\end{lemma}
\begin{proof}
We include the uniform remainder estimate for this classical sectorial
expansion. Choose $\pi/2<\theta<\phi/\alpha$ and a Hankel contour
$H_\theta$ running inwards on the ray $\arg u=-\theta$ to the unit
circle, counterclockwise through $1$ to the ray $\arg u=\theta$, then
outwards. Use the principal branch for powers of $u$. Hankel's
reciprocal gamma formula gives
\[
 E_{\alpha,\beta}(z)=\frac{1}{2\pi i}\int_{H_\theta}
        \frac{e^u u^{\alpha-\beta}}{u^\alpha-z}\,du.
\]
Indeed expansion for $|z|<1$ gives precisely the defining coefficients.
Analytic continuation along each ray of the stated sector is justified
by the separation between $\arg z$ and the angles of $u^\alpha$.
In particular $|u^\alpha-z|\ge c(|u|^\alpha+|z|)$ with $c>0$
independent of $u$ and $z$ in these sets. Use the exact decomposition
\[
 \frac{1}{u^\alpha-z}=-\frac1z-\frac{u^\alpha}{z^2}
                +\frac{u^{2\alpha}}{z^2(u^\alpha-z)}.
\]
The first two integrals give the displayed reciprocal gamma values.
The last is bounded in modulus by $|z|^{-3}$ times a constant
multiple of $\int_{H_\theta}|e^u u^{3\alpha-\beta}|\,|du|$.
This integral is uniformly bounded on the indicated $\beta$ intervals
because $|u|\ge1$ and the exponential decays along the rays.
\end{proof}

\begin{proof}[Proof of Theorem~\ref{ggps:main}]
Write $C=C_\alpha$, $d=s(\beta-\alpha+1)$ and $w=1-ih$.
Here $0\le d\le C/2$ and $-sw$ belongs to a closed sector from
Lemma~\ref{ggps:expansion}. At the simple zero of the reciprocal
gamma function,
\[
 \frac1{\Gamma(-1+t)}=\frac{t(t-1)}{\Gamma(1+t)}=-t+O(t^2).
\]
Also $-1/\Gamma(-1-\alpha+t)=C+O(t)$. The constant $C$ is positive
because $-1-\alpha\in(-3,-2)$. Substitution in the lemma gives
\begin{align*}
 s^2 E_{\alpha,\beta}(-s)&=C-d+O(s^{-1}),\\
 s^2 E_{\alpha,\beta}(-sw)&=\frac{C-dw}{w^2}+O(s^{-1}),
\end{align*}
uniformly on $0\le d\le C/2$. The first quantity is positive for
large $s$. For the leading terms the squared ratio of moduli satisfies
\[
 \frac{|C-dw|^2}{|w|^4(C-d)^2}
 =\frac{1+h^2\{d/(C-d)\}^2}{(1+h^2)^2}
 \le\frac1{1+h^2}<1.
\]
Since $C-d\ge C/2$, the uniform error terms preserve this strict
inequality for large $s$. Fixing any one such $s$ large enough that
$C/(2s)<1$ gives the asserted open interval of interior parameters.
\end{proof}

\paragraph{One finite interior witness.}
The same strict failure occurs for
\[
  \alpha=\tfrac32,\qquad \beta=\tfrac{501}{1000},\qquad
                z=-400+100i.
\]
With $E=E_{3/2,501/1000}$, the certified enclosures are
\[
 410213<10^{11}E(-400)<410214,\qquad
 390048<10^{11}|E(-400+100i)|<390049.
\]
Here is a reproducible finite check using real and complex ball
arithmetic. With precision 320 bits, evaluate the first 300 terms of
the defining series by repeated integer-power multiplication and the
reciprocal gamma function. For $r=|z|$ and
$T_k=r^k/\Gamma(3k/2+501/1000)$ one has
\[
 T_{k+2}/T_k=\frac{r^2}{g_k(g_k+1)(g_k+2)}<\tfrac1{500}
    \quad(k\ge300),\qquad g_k=3k/2+501/1000,
\]
since $r^2\le170000$ and $g_k>450$. Bound the sum of both omitted
parities by $500(T_{300}+T_{301})/499<10^{-210}$. The finite
calculation supports the explicit witness. The all-$\alpha$ analytic
argument above proves the refutation independently of finite arithmetic.

\paragraph{Scope and priority.}
The theorem determines failure in an open region adjacent to the
lower edge for every $1<\alpha<2$. The full set of parameters on
which \eqref{ggps:original} holds is still undetermined. The proof
uses a classical Hankel representation and a boundary scaling
argument. A bounded literature search found no earlier refutation
of Conjecture~25 in this interior region. The separate Kummer
question in Remark~6 of the same source has public counterexamples
and contributes no claim here.


\section{Sign oscillation in the telescoping continued fraction for the cubic beta tail}
\label{sec:br-sign}

In their September 2026 preprint \emph{The method of telescoping continued
fractions} \cite{BhatnagarRajkumar2026}, Bhatnagar and Rajkumar ask in
Conjecture~6.1 whether the partial numerators $a_j$ of their Algorithm TCF
eventually have a constant sign in each of their examples. We treat the
example $m=3$, $\epsilon=-1$ in their Appendix~A, equation (A.3).
We use $f(x)=(x+1/2)^3$ as in the body of their paper. The appendix uses
$f(x)=(2x+1)^3$, which divides the tail and its first partial numerator
by eight. All the resulting signs are the same, as verified below.

\begin{theorem}\label{brsign:main}
For the cubic alternating tail, Algorithm TCF produces infinitely many
strictly positive and infinitely many strictly negative partial numerators.
Among its first $k$ partial numerators, at most $O(\log(k+1))$ are
negative. In particular the positive signs have natural density one.
Consequently Conjecture~6.1 of Bhatnagar--Rajkumar is false.
\end{theorem}

We use a finite-dimensional question about a signed measure. Let $w(t)=
\operatorname{sech}(\pi t)$, and restrict the form
\[
  (P,Q)\longmapsto L(PQ),\qquad L(P)=\int_{\mathbb R}P(t)w''(t)\,dt,
\]
to real polynomials of degree less than $n$. How many independent negative
directions does it have? The answer needed here is that their number tends
to infinity, but its proportion tends to zero. Comparing this fact with
the block decomposition forced by the original algorithm determines the
oscillation of its signs.

\begin{lemma}\label{brsign:inertia}
Let $\nu_n$ be the negative index of this form on polynomials of degree
less than $n$. Then $\nu_n\to\infty$ and $\nu_n=O(\log(n+1))$.
\end{lemma}
\begin{proof}
We have
\[
   w''(t)=\pi^2w(t)\bigl(1-2\operatorname{sech}^2(\pi t)\bigr).
\]
Its sign is negative for $|t|<c$, positive for $|t|>c$, where
$c=\pi^{-1}\operatorname{arcosh}\sqrt2$. Also $|w''|\le\pi^2w$.
Polynomials are dense in $L^2(w\,dt)$. Indeed, if $h$ is orthogonal to
all polynomials, then $\int h(t)w(t)e^{zt}\,dt$ is holomorphic for
$|\operatorname{Re}z|<\pi/2$ by Cauchy--Schwarz. All derivatives at zero
vanish, and uniqueness of the Fourier transform gives $h=0$.
Approximate the indicators of any prescribed finite collection of disjoint
nonempty subintervals of $(-c,c)$ by real polynomials in this norm.
The limiting Gram matrix for $L$ is negative definite, so the polynomial
Gram matrix is negative definite once the errors are sufficiently small.
This proves $\nu_n\to\infty$ by nestedness.

For an upper bound, use the classical hyperbolic secant transform
\[
  \int_{\mathbb R} w(t)e^{zt}\,dt=\frac1{\cos(z/2)},
  \qquad |\operatorname{Re}z|<\pi.
\]
For example, $v=e^{2\pi t}$ reduces the integral to Euler's beta integral.
Define real polynomials $P_n$ by their generating function at zero,
\[
  U(t,s)=\sum_{n\ge0}P_n(t)s^n
   =\frac{\exp(2t\arctan s)}{\sqrt{1+s^2}}.
\]
These are the classical Meixner--Pollaczek polynomials with parameters
$\lambda=1/2$ and $\phi=\pi/2$ \cite[Eq.~18.23.7]{DLMFMP}. Each $P_n$ has degree $n$. The transform gives
$\int w(t)U(t,s)U(t,z)\,dt=(1-sz)^{-1}$, so the $P_n$ are orthonormal.
Integration by parts twice gives
\[
  \int w''(t)U(t,s)U(t,z)\,dt
   =\frac{4(\arctan s+\arctan z)^2}{1-sz}.
\]
These identities hold near $(s,z)=(0,0)$, where exponential domination
justifies coefficient extraction and the boundary terms vanish.
In particular,
\[
 L(P_n^2)=8\sum_{\substack{1\le l\le n\\l\ \mathrm{odd}}}\frac1{l^2},
 \qquad \pi^2-L(P_n^2)=8\sum_{\substack{l>n\\l\ \mathrm{odd}}}\frac1{l^2}.
\]
Let $G_n$ be the Gram matrix in this basis. The inequality $w''\le\pi^2w$
gives $G_n\le\pi^2 I$. Each negative eigenvalue of $G_n$ contributes
at least $\pi^2$ to $\operatorname{tr}(\pi^2 I-G_n)$, whose remaining
contributions are nonnegative. The last display shows that this trace is
$O(\log(n+1))$, proving the bound.
\end{proof}

\begin{proof}[Proof of Theorem~\ref{brsign:main}]
Put $F(x)=\sum_{r\ge0}(-1)^r/(x+r+1/2)^3$ for $x>0$. The gamma integral
for the reciprocal cube, followed by the secant transform in the proof
of Lemma~\ref{brsign:inertia}, gives
\begin{equation}\label{brsign:cauchy}
 F(x)=\frac14\int_0^\infty u^2 e^{-xu}\operatorname{sech}(u/2)\,du
     =-\frac14\int_{\mathbb R}\frac{w''(t)}{x-it}\,dt.
\end{equation}
In the second equality, use $(x-it)^{-1}=\int_0^\infty e^{-xu+itu}\,du$
and integrate twice by parts in $t$. The integrals converge absolutely.
Since $|w''|$ has every moment, finite geometric division gives a full
asymptotic expansion as $x\to+\infty$, with coefficient of $x^{-k-1}$
equal to $-i^k L(t^k)/4$. Only finite truncations will be needed.

Write $R_j=p_j/q_j$ for the convergents of the original algorithm, with
$p_0=0$, $q_0=1$, $p_{-1}=1$, $q_{-1}=0$. Its defining recurrences are
$p_j=b_jp_{j-1}+a_jp_{j-2}$ and $q_j=b_jq_{j-1}+a_jq_{j-2}$, where
$b_j$ is monic. Write $d_j=\deg b_j$ and $r_j=\deg q_j$. The source's
Theorem~4.1 and Proposition~5.1, together with its degree-selection rule,
give
\[
 N_j=fq_j(x)q_j(x+1)\bigl(1/f(x)-R_j(x)-R_j(x+1)\bigr),
 \quad d_{j+1}=3-\deg N_j.
\]
Here $N_0=1$, each nonzero $N_j$ is even in $x+1/2$ and has degree at
most two, and $b_j$ is odd. Thus $d_j\in\{1,3\}$ and all $a_j$ are
nonzero whenever the algorithm reaches this step. It continues forever.
Indeed $N_j=0$ would give a rational function $R$ satisfying
$R(x)+R(x+1)=1/(x+1/2)^3$. Every nonempty finite set of poles of $R$ in
one orbit modulo integral translations supplies two distinct uncancelled
poles of the sum, one at each extreme (the leftmost one shifted by $-1$).
A polynomial $R$ cannot work either. Hence no such rational $R$ exists.
The same pole argument and degree rules hold for the scaled summand in
(A.3). More directly, replacing $f$ by $8f$, $a_1$ by $a_1/8$,
and retaining every $b_j$ and all $a_j$ for $j\ge2$ scales all $p_j$
for $j\ge1$ by $1/8$, leaves $q_j$ and $N_j$ fixed, and meets every
unique cancellation prescribed by the algorithm.

By the recurrence, $r_j=\sum_{h=1}^j d_h$, and $q_j$ is monic of parity
$r_j$ (with $r_j\equiv j\pmod2$). Put $T_j=F-R_j$. Since
$F(x)+F(x+1)=1/f(x)$, the defining identity for $N_j$ gives
\[
 T_j(x)+T_j(x+1)=\frac{N_j(x)}{f(x)q_j(x)q_j(x+1)}.
\]
In an inverse-power expansion, adding the unit shift doubles its first
nonzero coefficient. The rational right side therefore shows
that the first nonzero term of $T_j$ has order $x^{-(2r_j+d_{j+1})}$.
This also excludes an identically zero formal expansion for $T_j$.
For $r=r_{j-1}$ and $d=d_j$, the continuant determinant gives
\[
 R_j-R_{j-1}=\frac{C_j}{q_jq_{j-1}},
 \qquad C_j=(-1)^{j-1}\prod_{h=1}^j a_h.
\]
As $T_j$ has strictly higher inverse order than $R_j-R_{j-1}$, we obtain
\begin{equation}\label{brsign:pivotres}
 q_{j-1}F-p_{j-1}= C_j x^{-r-d}+O(x^{-r-d-1}).
\end{equation}

Let $Q(t)=i^{-r}q_{j-1}(it)$, a real monic polynomial. By
\eqref{brsign:cauchy}, the coefficient of $x^{-h-1}$ on the left of
\eqref{brsign:pivotres} is $-i^{r+h} L(Qt^h)/4$. Hence
\[
 L(Qt^h)=0\quad (0\le h<r+d-1),\qquad
 h_j:=L(Q^2t^{d-1})=L(Qt^{r+d-1})\ne0.
\]
The two expressions for $h_j$ agree by monicity and the preceding
orthogonality. Since $r\equiv j-1\pmod2$, its sign is
\begin{equation}\label{brsign:pivotsign}
 \operatorname{sgn}h_j=(-1)^{(d-1)/2} S_j,
 \qquad S_j=-\operatorname{sgn}\prod_{h=1}^j a_h.
\end{equation}

Adding the polynomials $Q,tQ,\ldots,t^{d-1}Q$ extends the old space of
degree less than $r$ to that of degree less than $r+d$. This block is
orthogonal to the old space. Its entries below total exponent $d-1$
vanish and its anti-diagonal is $h_j$. For $d=1$, it contributes one
negative direction exactly when $S_j=-1$. If $d=3$, we have
\[
 h_j=L\bigl(Q^2(t)(t^2-c^2)\bigr)>0,
\]
because $L(Q^2)=0$ and $(t^2-c^2)w''(t)>0$ away from $\pm c$.
Thus $d=3$ can only occur when $S_j=-1$. Its Gram matrix has the form
\[
 \begin{pmatrix}0&0&h_j\\0&h_j&u\\h_j&u&v\end{pmatrix}.
\]
Its determinant is negative and it has a positive diagonal entry, so
its negative index is one. In either case the increase in negative
index equals $\mathbf{1}_{\{S_j=-1\}}$. Consequently
\begin{equation}\label{brsign:exact}
 \nu_{r_k}=\#\{j\le k:S_j=-1\},\qquad k\le r_k\le3k.
\end{equation}
Lemma~\ref{brsign:inertia} now shows that the set on the right has an
unbounded counting function of order $O(\log(k+1))$. In particular
$S_j$ takes both signs infinitely often. Its sign therefore changes
infinitely often. Each sign change is adjacent to a negative $S_j$,
and each such $S_j$ accounts for at most two changes. Since
$S_j/S_{j-1}=\operatorname{sgn}a_j$ for $j\ge2$, the number of
negative $a_j$ through index $k$ tends to infinity and is at most
$2\nu_{r_k}+1=O(\log(k+1))$. This also proves the claim on positive signs.
\end{proof}


\paragraph{Where the first sign reversals occur.}
For Appendix~A.3 of the source, the negative partial numerators among
$a_1,\ldots,a_{598}$ have indices $2$, $450$, and $451$.
In particular, the positive run $a_3,\ldots,a_{449}$ does not persist.
To reproduce these finite indices, expand the bivariate Gram generating
function from Lemma~\ref{brsign:inertia} to degrees below $600$.
The first three degrees form the initial size-three block. Take its
Schur complement and perform symmetric elimination separately on
the even and odd degrees, which are orthogonal by parity. At 192-bit
precision, every remaining diagonal pivot interval excludes zero.
Thus all subsequent blocks in this computation have size one and the
ordered pivot signs are the $S_j$ in \eqref{brsign:pivotsign}. A change
in $S_j$ gives a negative $a_j$. The self-contained source includes
the interval-arithmetic verifier after its end-of-document marker.

\section{Verification and unresolved requirements}
\label{sec:status}

The stability argument is algebraic and applies to arbitrary finite
elementary towers. The graph argument is combinatorial. The other
arguments use moments, generating functions, or complex analysis. The explicit Wright witness in
Theorem~\ref{thm:wright-interior} uses rational error bounds rather than a
floating-point sign test. The negative third cumulant in
Theorem~\ref{thm:yyy-explicit} is certified by an exact trigonometric
inequality. The total-positivity argument applies to arbitrary finite
minors, and the monotonicity argument applies to every positive real
argument. Finite experiments are not substituted for any of these
quantifiers.

Separate internal AI derivations checked the principal algebraic identities,
parameter domains, and convergence arguments. This is useful error checking,
but it is neither external peer review nor formal proof-assistant
verification. In particular, the manuscript should be read as a research
draft whose arguments are exposed for scrutiny.

\subsection{Literature coverage}

For the Bell criterion, the original question was compared with
Zhu's 2024 results~\cite{Zhu2024}, the exponential-Riordan criteria of
Zhu~\cite{Zhu2021}, and the Stieltjes-moment results of Wang and
Zhu~\cite{WangZhu2016}. Zhu's 2024 Proposition~6.18 establishes the
pointwise Stieltjes property and the shifted-cumulant input, while its
Conjecture~6.16(ii) still asks for coefficientwise positivity in the
component variable. The distinction is essential. The proof here supplies
its own positive measure for the input and does not depend on a displayed
indexing formula in that article.

Dominici's author preprint for the Hankel-expansion
paper~\cite{Dominici2020} was also inspected. It treats discrete weights
$\omega(k)z^k$ and leading Hankel determinants, with Vandermonde and Schur
expansions; no generic shifted-cumulant criterion was found there. Differences
between that preprint and every detail of the final article were not checked.

Two relevant full texts could not be inspected in this research session:
Zhu's 2026 paper~\cite{Zhu2026} and the general criterion in
Ding--Mu--Zhu~\cite{DingMuZhu2026}. Their accessible metadata and previews are
insufficient to exclude an equivalent theorem. Thus the general criterion
and its application must not be announced as a verified world-first result
on the basis of this draft.

For the Wright question, the original full text and subsequent work
including Ferreira--Simon~\cite{FerreiraSimon2025} were checked. For the
negative-power question, the original remark, the related GB2 question of
Bosch--Simon~\cite{BoschSimon2015}, and the stated ranges in
Wang~\cite{Wang2024,Wang2026} were compared. The full original Yang
statement and its journal version were also checked. Kang's 28 October
2026 preprint \cite{Kang2026} already refutes the separate
Conjecture~1(ii), and its argument was inspected directly. That result
is excluded from this draft's contribution count. Exact-title,
conjecture, counterexample, and method searches located no identical
resolution of the four older targets in the literature inspected. This records
the extent of the search; it is not a proof that no earlier solution
exists.

For elementary stability, the original Chen--Kauers question,
Chen's 2022 stability paper, the Sahu--Seiss--Srinivasan generalized
Liouville theorem, and the 2026 complete-reduction paper were compared.
The 2026 paper of Sahu--Srinivasan~\cite{stab:SS2026} concerns elementary
first integrals of triangular integral systems; no identical all-depth
stability characterization was located there. Three independent internal
reviews checked the local valuations, finite coefficient extension,
logarithmic divisor argument, quotient-module descent, and finite-test
classification. No fatal gap was found. This is a proposed complete
proof for review, not an independently established novelty certificate.

For the all-order Yang problem, both the original preprint and final
journal statement were checked, including the variable $x$ and the
prescribed decay. Orthogonal-polynomial Pad\'e methods have substantial
precedents: Pr\'evost and Rivoal \cite{PrevostRivoal2021} give effective
error estimates for nearly diagonal approximants to digamma asymptotic
expansions. Only their abstract and first-page preview were accessible;
the full paper remains a material priority-check gap. Li and Chen
\cite{LiChen2017} treat Pad\'e approximants for normalized gamma
asymptotics; the complete-monotonicity lemma inspected there concerns
ordinary Stirling truncations. Standard Stieltjes remainder identities
give complete monotonicity in $s=x^2$, which does not imply the required
property in $x$.

Section~\ref{sec:yang-allorders} supplies the additional matrix-marginal
argument and its exact transform identity. The underlying Heine and
Hermite--Genocchi identities, spectral Jacobian, Schur complement,
determinant lemma, Bernstein representation, and gamma integral are
classical ingredients. Two separate internal audits checked the full
proof, including the gamma shape, rate normalization, bounded-transform
interchanges, uniqueness, and every-order quantifier. A bounded search
did not locate an identical all-order result or the exact intermediate
matrix-marginal theorem; it does not certify their absence from the
literature.

For Li--Chen Conjecture~4.1, the original published inequalities and
coefficient system were read in full. Section~\ref{sec:li-chen-pade}
proves every order, proves denominator positivity, and identifies the
quadrature expressions with those exact original approximants.
Conjecture~3.1 on coefficient symmetry follows in the same theorem;
McCabe's inspected primary preview \cite[pp.~67--68]{McCabe1989}
shows that this symmetry is already classical, explicitly for the scaled
gamma series. It is not counted as another new target. An independent internal
audit checked the complete analytic argument. Exact rational calculations
through order $16$ and finite high-precision sign checks were additional
diagnostics, not the justification of the all-order conclusion.
A bounded search located no earlier identical solution. However, Luke's
1970 paper \cite{Luke1970}, N\'emeth--Zim\'anyi
\cite{NemethZimanyi1982}, and the remainder of McCabe's chapter
\cite{McCabe1989} remain unread and potentially relevant to the
alternating inequality theorem.
Historical priority is therefore unresolved. The positive-real
representation, moment extraction, and quadrature sign mechanisms are
classical; no new foundational theorem is claimed for them.

\paragraph{Graffiti.pc 177.}
West's statement and the 2008 theorem of DeLaVi\~na--Waller were
compared directly. The latter authors explicitly distinguished their
$2\alpha+1$ bound from Conjecture~177. Corpus [066] treated claw-free
graphs. The rooted deletion lemma and the leaf bound in
Section~\ref{sec:wow177} apply to arbitrary connected graphs.
Ducoffe's work on graph extremities \cite{Ducoffe2022} treats the
related property that deleting a closed neighborhood leaves a connected
graph. The lemma here prescribes an anchor in a given maximal independent
set and keeps a single neighbor. An independent argument audit checked
every step and the case $K_2$. As a diagnostic, the deletion statement
was also tested on the 995 connected graphs on two to seven vertices in
the graph atlas. A bounded search of citing works and public conjecture
records located no earlier proof of the unrestricted inequality.

\paragraph{Garrappa--Gerhold--Popolizio--Simon 25.}
The published statement and its parameter boundaries were checked.
The source already uses sectorial asymptotics in other parameter
regions. Here the first reciprocal-gamma coefficient vanishes at
$\beta=\alpha-1$. A two-term expansion along oblique rays detects
failure inside the conjectured region. A separate audit checked the
Hankel contour and the uniform remainder. Two implementations enclosed
the displayed finite witness using interval arithmetic. The source's
Kummer question and its Conjecture~26 have different statements.
The literature examined did not contain this interior counterexample.

\paragraph{Bhatnagar--Rajkumar 6.1 (2026).}
The source's degree rule, parity theorem, and Appendix~A.3 were compared
against the proof. The argument derives its sign sequence from the
original Algorithm TCF and permits degree-three steps. An independent
audit checked the rotated moment functional, residual orders, pivot signs,
and negative index of each added block. Exact rational calculations
reproduced the displayed coefficients and the index identity through
30 steps. The primary evidence for infinite oscillation is the index
argument valid for all orders. A bounded literature search located no
earlier proof of the sign conclusion for the cubic alternating tail.
Indefinite moment problems and generalized Jacobi matrices have an
extensive earlier theory, including Derevyagin--Derkach \cite{IndefJacobi}.
The contribution proposed here is the index estimate for this signed
weight and its exact identification with the signs of the source algorithm.

\subsection{What remains open in this project}

The ten complete mathematical arguments have undergone internal checks.
Independent specialist review and historical comparison are pending.
In particular, the unread works identified above prevent a definitive
priority claim for the Bell and Pad\'e results. The tools used in
several sections have well-established precedents. The manuscript
makes no ranking against the entire corpus.

Erd\H{o}s Problem~461 still requires a uniform linear lower bound, and
the work here also leaves the unrestricted Erd\H{o}s and multivariate
forest questions open. Chen's moment-set Skolem conjecture remains
separate from the classification in Section~\ref{stab:section}.
Appendix~\ref{sec:br-coefficients} establishes stabilization for a
2023 algorithm whose conjectural wording omits an index. Its result
is ancillary. The 2026 sign conjecture in Section~\ref{sec:br-sign}
has an explicit example and is counted separately.

The accompanying source is one self-contained \texttt{.tex} file with an
embedded bibliography. It requires no external figures, data files,
\texttt{.bib} file, or subsidiary \texttt{.tex} files. Run
\texttt{pdflatex} until its internal-reference warnings disappear;
the delivered version was compiled with three passes.

\appendix
\section{A classical proof of Sokal's forest-triangle clause}
\label{app:forest-triangle}

This appendix concerns exactly Conjecture~6.11(a) of
Sokal~\cite{Sokal2022}. It does not resolve parts~(b),~(c), or the shifted
Conjecture~6.12. The proof is an elementary composition of known totally
positive matrices. Its inclusion records a complete mathematical argument,
not a claim of a new foundational method or certified historical priority.

Let $T(z)=ze^{T(z)}$ be the tree function and define the exponential
Riordan matrix
\[
 R[1,G]_{n,k}=\frac{n!}{k!}[z^n]G(z)^k.
\]
Write $F=R[1,T]$ for the rooted-forest triangle, which is totally positive
by Sokal's Theorem~1.1(a). The original conjecture asks whether
\[
 F^{\sharp}(w)=R\left[1,\frac{e^{wT(z)}-1}{w}\right]
\]
is coefficientwise totally positive in $w$. The expression at $w=0$ is
interpreted by its polynomial limit.

\begin{proposition}
The matrix $F^{\sharp}(w)$ is coefficientwise totally positive.
\end{proposition}
\begin{proof}
Let $S(n,k)$ denote a Stirling number of the second kind and put
$S_w(n,k)=S(n,k)w^{n-k}$ for $n\ge k$, and zero otherwise.
The exponential generating function for $S(n,k)$ gives
\[
 S_w=R\left[1,\frac{e^{wz}-1}{w}\right].
\]
The ordinary Stirling matrix is totally positive. One standard verification
uses $S(n,k)=h_{n-k}(1,2,\ldots,k)$: these are path weights in a planar
network with horizontal edges of weight $j$ at level $j$ and vertical
edges of weight one. Sources $(-n,1)$ and sinks $(-k,k)$ give the displayed
complete homogeneous polynomial; the nonintersecting-path determinant
identity gives nonnegative minors. The $k=0$ row and column form a separate
identity block.

For increasing index sets $I,J$ of the same size, every nonzero term in
the minor of $S_w$ has the same power of $w$, so
\[
 \det (S_w)_{I,J}
 =w^{\sum_{i\in I}i-\sum_{j\in J}j}\det S_{I,J}.
\]
If the exponent would be negative, every determinant term is zero by
lower triangularity. Hence $S_w$ is coefficientwise totally positive.
Riordan multiplication now gives the exact factorization
\[
 F^{\sharp}(w)=F S_w.
\]
Cauchy--Binet expresses every finite minor of this product as a finite
sum of products of minors of $F$ and $S_w$. Each summand has nonnegative
coefficients, proving the claim.
\end{proof}

No explicit later resolution of this exact clause was located in the
bounded literature check. Nevertheless, the method is classical, and
Zhu~\cite{Zhu2024} uses an analogous forest--Stirling factorization for
functional digraphs. The two unread 2026 sources discussed above remain
part of the priority check. This appendix must not be advertised as a
high-difficulty new-method solution, or as a solution of all of
Conjecture~6.11.

\section{Exact boundaries of the Bell-kernel method}
\label{app:boundaries}

The criterion in Section~\ref{sec:tp} concerns a fixed \emph{real}
cumulant sequence and coefficientwise positivity in one Bell parameter.
The following exact examples rule out three tempting extensions. They
were formulated while developing the present method; no pre-existing
named open conjecture is identified with these extensions. Consequently
none is counted toward the requested ten resolutions.

\subsection{Polynomial cumulants do not give a coefficientwise extension}

Let $q$ be an indeterminate, and define $f_n(q)$ by the finite Stieltjes
continued fraction
\begin{equation}\label{eq:boundary-sfrac}
 \sum_{n\ge0}f_n(q)z^n
 =\frac{1}{1-\dfrac{z}{1-\dfrac{qz}{1-q^8z}}}
 =\frac{1-(q+q^8)z}{1-(1+q+q^8)z+q^8z^2}.
\end{equation}
The nonnegative continued-fraction weights $1,q,q^8,0,\ldots$ imply
that $(f_n(q))$ is coefficientwise Hankel-totally positive. This follows
from the standard weighted-Dyck-path and nonintersecting-path
interpretation, so it is a formal polynomial assertion, not just
positivity at fixed $q\ge0$.

Set $\kappa_n(q)=f_{n-1}(q)$ for $n\ge1$, and define
\[
 \sum_{n\ge0}m_n(t,q)\frac{z^n}{n!}
 =\exp\left(t\sum_{n\ge1}\kappa_n(q)\frac{z^n}{n!}\right).
\]
Thus the shifted cumulants satisfy the polynomial analogue of the input
condition in Section~\ref{sec:tp}.

\begin{proposition}
The sequence $(m_n(t,q))$ is not coefficientwise Hankel-totally positive
jointly in $t,q$.
\end{proposition}
\begin{proof}
Equation~\eqref{eq:boundary-sfrac} gives $f_0=f_1=1$ and
\[
 f_n=(1+q+q^8)f_{n-1}-q^8f_{n-2}\qquad(n\ge2).
\]
The exponential formula gives the exact Bell recurrence
\[
 m_0=1,\qquad
 m_n=t\sum_{j=1}^{n}\binom{n-1}{j-1}f_{j-1}m_{n-j}.
\]
Applying these finite recurrences through $n=6$ to the minor
$D=m_4m_6-m_5^2$ gives
\begin{align*}
 [t^3]D={}&7q^{25}+8q^{18}-10q^{17}+9q^{11}+18q^{10}+9q^9\\
          &+8q^4+32q^3+48q^2+32q+8.
\end{align*}
In particular, $[t^3q^{17}]D=-10$. This is a negative coefficient of a
$2\times2$ Hankel minor, proving the assertion.
\end{proof}

For every \emph{fixed} $q\ge0$, the numerical hypothesis of
Section~\ref{sec:tp} still applies. There is no contradiction: positivity
of coefficients in $t$ at each fixed $q$ does not imply positivity of the
individual coefficients in $q$.

\subsection{Geometric mixing and factorial lifting can fail}

Consider
\[
 K(z)=z+\frac{z^2}{1-21z}
     =\sum_{n\ge1}\kappa_n\frac{z^n}{n!},\qquad
 m_n(x)=n![z^n]\frac1{1-xK(z)}.
\]
Here $\kappa_1=1$ and $\kappa_n=n!21^{n-2}$ for $n\ge2$.
Let $G$ be the gamma probability measure of shape $2$ and scale $21$.
The positive measure
\[
 \mu=\frac{20}{21}\delta_0+\frac1{21}G
\]
represents $(\kappa_{n+1})_{n\ge0}$: for $n\ge1$,
$\int u^n\dd\mu(u)=(n+1)!21^{n-1}=\kappa_{n+1}$, and the
zeroth moment is one.

\begin{proposition}
The shifted cumulants in this example form a numerical Stieltjes moment
sequence, and $(m_n(x)/n!)$ is coefficientwise Hankel-totally positive,
but $(m_n(x))$ is not coefficientwise Hankel-totally positive.
\end{proposition}
\begin{proof}
The moment representation was just verified. The normalized generating
function has the positive Stieltjes fraction
\[
 \frac1{1-xK(z)}
 =\frac1{1-\dfrac{xz}{1-\dfrac{z}{1-20z}}},
\]
which proves the coefficientwise assertion for $(m_n/n!)$.
For the unnormalized sequence, direct coefficient extraction gives
\begin{align*}
 m_1={}&x, & m_2={}&2x^2+2x,\\
 m_3={}&6x^3+12x^2+126x,\\
 m_4={}&24x^4+72x^3+1032x^2+10584x,\\
 m_5={}&120x^5+480x^4+7920x^3+110880x^2+1111320x.
\end{align*}
Expanding the shifted leading minor yields
\begin{align*}
 \det(m_{i+j+1})_{i,j=0}^{2}
 =24x^3\big(&x^6+6x^5+195x^4-860x^3\\
            &+86295x^2+79926x+1120581\big).
\end{align*}
The coefficient of $x^6$ is $-20640<0$.
\end{proof}

The same construction with $K_b(z)=z+z^2/(1-bz)$, $b\ge1$, gives
\[
 [x^6]\det(m_{i+j+1})_{i,j=0}^{2}=96(-b^2+11b-5).
\]
Thus the failure is systematic. It does not refute Sokal's
Conjecture~6.7, which fixes the specific tree function $K=T$.
The already known failure of general coefficientwise Hadamard closure
also cautions against factorial lifting; the example here supplies an
explicit certificate inside the proposed geometric construction.

\subsection{A nested exponential formula also needs extra hypotheses}

Keep the same numerical $K$, and define
\[
 H_w(z)=\frac{e^{wK(z)}-1}{w},\qquad
 \sum_{n\ge0}M_n(t,w)\frac{z^n}{n!}=e^{tH_w(z)}.
\]
Although the innermost shifted cumulants are numerical Stieltjes
moments, the nested formula need not give joint coefficientwise Hankel
positivity. Indeed, with
$D=\det(M_{i+j+1}(t,w))_{i,j=0}^{2}$, one obtains
\[
 [t^4]D=2w^4-82w^3-5144w^2+1035540w+959112.
\]
For completeness, this identity can be checked without any infinite
series manipulation. Start with
\[
 (\kappa_1,\ldots,\kappa_5)=(1,2,126,10584,1111320),
\]
form $b_0=1$ and
$b_n=w\sum_{j=1}^{n}\binom{n-1}{j-1}\kappa_jb_{n-j}$,
put $h_n=b_n/w$, and then form
$M_0=1$ and
$M_n=t\sum_{j=1}^{n}\binom{n-1}{j-1}h_jM_{n-j}$.
These recurrences through $n=5$ give the displayed determinant
coefficient. The negative $t^4w^2$ coefficient is $-5144$.
This rules out a general nested-Bell shortcut, but does not decide
Sokal's forest conjectures, which have additional structure.

\subsection{Scope of these examples}

The three obstructions identify limits of the Bell-kernel method.
The multivariate forest assertions remain open in this work.
The classical forest-triangle argument and the three method obstructions
are ancillary to the ten targets in the opening table.

\section{Ancillary result: stabilization of the Bhatnagar--Rajkumar algorithm}
\label{sec:br-coefficients}

\paragraph{The precise sequence in the original question.}
Bhatnagar and Rajkumar \cite{BhatnagarRajkumar2023} construct a triangular array of continued-fraction
coefficients from Legendre rational approximants. Their Section~4
observes experimentally that the coefficient in position $m$ stabilizes when
$k\ge m+1$, where $k$ is the Legendre order. Conjecture~6.2 asks
whether the sequence generated by their Theorem~5.2 agrees with the
classical Stieltjes continued fraction for Binet's function, without
specifying the Legendre index in that statement.
We prove the stabilized identity by identifying exactly what the
leading-degree cancellation algorithm matches.

The dependence on $k$ must be kept explicit: the fixed-$k$ rows do
not equal the Binet coefficients at every position, as the original
Table~2 already shows. The stabilized sequence is the one displayed
in their equation (6.1) and discussed immediately before their
Conjecture~6.2. This supports a stabilized interpretation of their
question, but that interpretation is not explicit in its wording.
Accordingly, this result is recorded as an ancillary theorem, not
counted as an unchanged-statement resolution of Conjecture~6.2.

Set
\[
 \begin{aligned}
 J(x)&=\log\Gamma(x+1)-(x+\tfrac12)\log x+x-\tfrac12\log(2\pi),\\
 \epsilon(x)&=J(x)-J(x+1)
       =(x+\tfrac12)\log\frac{x+1}{x}-1.
 \end{aligned}
\]
Let $P_k$ be the usual Legendre polynomial and $P_k^*$ its associated
numerator polynomial, with the normalization used in the original
paper. Put
\[
 E_k(x)=(2x+1)\frac{P_k^*(2x+1)}{P_k(2x+1)}-1
        =\frac{f_k^*(z)}{f_k(z)},\qquad z=x(x+1),
\]
where $d_k=\deg f_k\ge1$ and $\deg f_k^*=d_k-1$ for $k\ge2$.
For a list of coefficients $a_1,\ldots,a_m$, write
\[
 G_m(x)=\cfrac{a_1}{x+\cfrac{a_2}{x+
                 \cfrac{\ddots}{x+\cfrac{a_m}{x}}}}
        =\frac{N_m(x)}{D_m(x)},
\]
with the continuant normalization
$D_{-1}=0,D_0=1$ and $D_j=xD_{j-1}+a_jD_{j-2}$.
Thus $D_m$ is monic of degree $m$.
The polynomial whose leading coefficient is cancelled in the
original algorithm is
\[
 \Delta_{k,m}(z)=f_k^*(z)D_m(x)D_m(x+1)
  -f_k(z)\{N_m(x)D_m(x+1)-N_m(x+1)D_m(x)\}.
\]
Its dependence only on $z$ follows from the parity of the continuants,
or from the original paper's Lemma~2.3.

\begin{theorem}\label{br:main}
Let $b_j>0$ denote the coefficients of the classical continued fraction
\[
 J(x)=\cfrac{b_1}{x+\cfrac{b_2}{x+\cfrac{b_3}{x+\ddots}}},
 \qquad b_1=\frac1{12},\quad b_2=\frac1{30},\quad
 b_3=\frac{53}{210},\ \ldots.
\]
For every $m\ge1$ and $k\ge m+1$, the first $m$ successive
leading-coefficient cancellations in the Bhatnagar--Rajkumar
algorithm are well-defined and uniquely select
\[
 a_j(k)=b_j,\qquad 1\le j\le m.
\]
The resulting polynomial $\Delta_{k,m}$ has degree at most $d_k-1$.
If $k\ge m+2$, this degree is exactly $d_k-1$.
In particular, the stabilized sequence exists and agrees term by
term with the classical Binet sequence displayed in the source.
\end{theorem}

\begin{proof}
All order statements in this proof may be read in the formal Laurent
series field at infinity; convergence of the Stirling series is not
used. Write $\delta T(x)=T(x)-T(x+1)$. If
$T(x)=c x^{-r}+O(x^{-r-1})$ with $c\ne0$ and $r>0$, then
\begin{equation}\label{br:valuation}
 \delta T(x)=r c x^{-r-1}+O(x^{-r-2}).
\end{equation}

We first recall the exact approximation orders of the two classical
continued fractions appearing in the original paper. The Legendre
quotient $P_k^*(y)/P_k(y)$ is the $k$-node Gaussian approximation to
\[
 \frac12\log\frac{y+1}{y-1}
       =\frac12\int_{-1}^1\frac{dt}{y-t}.
\]
Its moments agree through degree $2k-1$, so finite geometric
division gives an error $O(y^{-2k-1})$. Consequently
\begin{equation}\label{br:legendre-order}
 E_k(x)-\epsilon(x)=O(x^{-2k}).
\end{equation}

Let $B_m$ be the $m$th convergent of the classical Binet fraction.
The standard Stieltjes continued-fraction construction gives
\begin{equation}\label{br:binet-order}
 J(x)-B_m(x)=c_m x^{-2m-1}+O(x^{-2m-3}),
 \qquad c_m=(-1)^m\prod_{j=1}^{m+1}b_j\ne0.
\end{equation}
For clarity, the positive moment measure underlying this standard
construction is supplied directly by Binet's integral:
\[
 \frac{J(\sqrt s)}{\sqrt s}
 =\int_0^\infty\frac{d\nu(u)}{s+u},\qquad
 d\nu(u)=\frac{-\log(1-e^{-2\pi\sqrt u})}{2\pi\sqrt u}\,du.
\]
This measure is positive on $(0,\infty)$ and has every moment;
the ordinary and shifted Hankel forms are strictly positive.
The classical Stieltjes algorithm therefore has positive,
nonzero coefficients at every step. Each added level changes the
convergent first at order $x^{-2m-1}$, and the continuant determinant
identity gives exactly the coefficient in
\eqref{br:binet-order}. These are also the coefficients $b_j$ and
the approximation property described in Section~6 of the original
paper, not a different continued fraction.

For $k\ge m+1$, equations \eqref{br:valuation}--\eqref{br:binet-order}
give
\[
 E_k(x)-\delta B_m(x)
  =E_k(x)-\epsilon(x)+\delta(J-B_m)(x)
  =O(x^{-2m-2}).
\]
Since $f_k(z)$ has degree $d_k$ and
$D_m(x)D_m(x+1)$ is a monic polynomial of degree $m$ in $z$, the
corresponding cross-multiplied residual satisfies
\[
 \Delta_{k,m}(z)=O(x^{2d_k-2})=O(z^{d_k-1}).
\]
Thus choosing $a_1=b_1,\ldots,a_m=b_m$ makes precisely the
leading coefficient required by the algorithm vanish.

We verify uniqueness and absence of a breakdown before reaching
position $m$. Suppose the previous coefficients are already fixed
at $b_1,\ldots,b_{j-1}$ and regard $a_j$ as variable.
The continuant identity is
\[
 G_j(x)-G_{j-1}(x)
   =\frac{(-1)^{j-1}a_1\cdots a_j}{D_j(x)D_{j-1}(x)}.
\]
Changing $a_j$ by an amount $v$ therefore changes $G_j$ first by
\[
 (-1)^{j-1}(b_1\cdots b_{j-1})v\,x^{-2j+1}.
\]
By \eqref{br:valuation}, it changes $\delta G_j$ first by
\[
 (2j-1)(-1)^{j-1}(b_1\cdots b_{j-1})v\,x^{-2j}.
\]
The factor multiplying $v$ is nonzero. Before this choice the lower
inverse-power coefficients have already been matched. Hence the
coefficient of $z^{d_k}$ in the cross-multiplied residual is affine
in $a_j$ with a nonzero slope. There is exactly one cancellation,
and the existence argument shows that it is $a_j=b_j$.
Induction proves the assertion for every $j\le m$.

One can also check directly that the degree at every preceding
stage is the degree needed for the next cancellation. For $j<m$,
we have $k\ge j+2$, so the Legendre error in
\eqref{br:legendre-order} occurs strictly after the first nonzero
term of $\delta(J-B_j)$. The latter is
$(2j+1)c_jx^{-2j-2}$ by \eqref{br:binet-order}. It cannot cancel,
and therefore $\deg\Delta_{k,j}=d_k-1$. The same argument gives
$\deg\Delta_{k,m}=d_k-1$ whenever $k\ge m+2$.
At the boundary $k=m+1$, degree at most $d_k-1$ suffices to define
the unique coefficient already selected; no later continuation of
that fixed row is required for stabilization.
\end{proof}

\paragraph{Scope and attribution.}
The result identifies the stabilized coefficients from the original
algorithm. It does not assert every fixed-$k$ row continues forever,
or prove every coefficient-sign observation about all residual
polynomials in Section~4 of the source. Those stronger assertions
are not proved here. The usual alternating Binet bounds follow
from the established coefficient identity and classical Stieltjes
theory; they are not counted as a second closure. In particular, the
reversed signs in the display accompanying Conjecture~6.1, which
conflict with its prose and the first convergent, are not treated as
a new counterexample.

\paragraph{Literature boundary.}
The original article was read in full. The same authors' September 2026
follow-up \cite{BhatnagarRajkumar2026} was also inspected: it proves
uniqueness for a general rational-kernel cancellation algorithm, but
does not give the Binet coefficient-matching argument above.
A bounded search found no identical stabilization theorem. This does
not establish historical priority; the formal-series, quadrature, and
continuant tools used here are classical. The original notation issue
is sufficient to exclude this appendix from the strict target count.


\begin{thebibliography}{99}
\setlength{\itemsep}{2pt plus .3pt minus .3pt}

\bibitem{Sokal2022}
A. D. Sokal,
\emph{Total positivity of some polynomial matrices that enumerate labeled
trees and forests I: forests of rooted labeled trees},
Monatshefte f\"ur Mathematik \textbf{200} (2023), 389--452;
published online 2022.
\href{https://doi.org/10.1007/s00605-022-01687-0}{doi:10.1007/s00605-022-01687-0}.
\href{https://arxiv.org/abs/2105.05583}{arXiv:2105.05583}.

\bibitem{Zhu2024}
B.-X. Zhu,
\emph{Coefficientwise Hankel-total positivity of the row-generating
polynomials for the output matrices of certain production matrices},
Advances in Mathematics \textbf{447} (2024), 109695.
\href{https://arxiv.org/abs/2202.03793v2}{arXiv:2202.03793v2}.

\bibitem{Zhu2021}
B.-X. Zhu,
\emph{Total positivity from the exponential Riordan arrays},
SIAM Journal on Discrete Mathematics \textbf{35} (2021), 2971--3003.
\href{https://arxiv.org/abs/2006.14485}{arXiv:2006.14485}.

\bibitem{WangZhu2016}
Y. Wang and B.-X. Zhu,
\emph{Log-convex and Stieltjes moment sequences},
Advances in Applied Mathematics \textbf{81} (2016), 115--127.
\href{https://arxiv.org/abs/1612.04114}{arXiv:1612.04114}.

\bibitem{Lambert2012}
G. A. Kalugin, D. J. Jeffrey, R. M. Corless and P. B. Borwein,
\emph{Stieltjes and other integral representations for functions of Lambert
$W$}, Integral Transforms and Special Functions \textbf{23} (2012), 581--593.
\href{https://doi.org/10.1080/10652469.2011.613830}{doi:10.1080/10652469.2011.613830}.

\bibitem{Schilling2012}
R. L. Schilling, R. Song and Z. Vondra\v{c}ek,
\emph{Bernstein Functions: Theory and Applications}, second edition,
De Gruyter Studies in Mathematics 37, De Gruyter, Berlin, 2012.

\bibitem{Macdonald1995}
I. G. Macdonald,
\emph{Symmetric Functions and Hall Polynomials}, second edition,
Oxford Mathematical Monographs, Oxford University Press, 1995.

\bibitem{Yang2014}
Z.-H. Yang,
\emph{Several completely monotone functions related to DeTemple's sequence},
preprint, 2014.
\href{https://arxiv.org/abs/1409.6411}{arXiv:1409.6411}.

\bibitem{YangChuZhang2015}
Z.-H. Yang, Y.-M. Chu and X.-H. Zhang,
\emph{Necessary and sufficient conditions for functions involving the psi
function to be completely monotonic},
Journal of Inequalities and Applications \textbf{2015} (2015), 157.
\href{https://doi.org/10.1186/s13660-015-0674-8}{doi:10.1186/s13660-015-0674-8}.

\bibitem{FerreiraSimon2024}
R. A. C. Ferreira and T. Simon,
\emph{On the log-concavity of the Wright function},
Constructive Approximation \textbf{60} (2024), 309--338.
\href{https://doi.org/10.1007/s00365-023-09666-w}{doi:10.1007/s00365-023-09666-w}.
\href{https://arxiv.org/abs/2212.07974v2}{arXiv:2212.07974v2}.

\bibitem{YanoYanoYor2009}
K. Yano, Y. Yano and M. Yor,
\emph{On the laws of first hitting times of points for one-dimensional
symmetric stable L\'evy processes},
S\'eminaire de Probabilit\'es XLII, Lecture Notes in Mathematics 1979,
Springer, 2009, pp. 187--227.
\href{https://doi.org/10.1007/978-3-642-01763-6_8}{doi:10.1007/978-3-642-01763-6\_8}.
\href{https://arxiv.org/abs/0811.2046}{arXiv:0811.2046}.

\bibitem{BoschSimon2015}
P. Bosch and T. Simon,
\emph{On the infinite divisibility of inverse Beta distributions},
preprint version consulted, 2014.
\href{https://arxiv.org/abs/1405.4176v2}{arXiv:1405.4176v2}.

\bibitem{Wang2024}
M. Wang,
\emph{Moments of Gamma type and three-parametric Mittag-Leffler function},
preprint, 2024.
\href{https://arxiv.org/abs/2410.19330}{arXiv:2410.19330}.

\bibitem{Wang2026}
M. Wang,
\emph{Infinite divisibility of $\alpha$-Cauchy distributions},
preprint, version 3, 15 April 2026.
\href{https://arxiv.org/abs/2512.23164v3}{arXiv:2512.23164v3}.

\bibitem{FerreiraSimon2025}
R. A. C. Ferreira and T. Simon,
\emph{Mittag-Leffler functions and convex ordering}, preprint, 2025.
\href{https://arxiv.org/abs/2512.04940}{arXiv:2512.04940}.

\bibitem{Zhu2026}
B.-X. Zhu,
\emph{Coefficientwise total positivity and Stieltjes moment properties from
Riordan arrays}, Journal of the London Mathematical Society
\textbf{113} (2026), e70450.
\href{https://doi.org/10.1112/jlms.70450}{doi:10.1112/jlms.70450}.
Full text not inspected; cited as a remaining priority-check item.

\bibitem{Dominici2020}
D. Dominici,
\emph{Power series expansion of a Hankel determinant},
Linear Algebra and its Applications \textbf{601} (2020), 17--54.
\href{https://doi.org/10.1016/j.laa.2020.04.014}{doi:10.1016/j.laa.2020.04.014}.
The \href{https://www.dk-compmath.jku.at/publications/dk-reports/2018-06-12/at_download/file}{author preprint, DK-Report 2018-08}, was inspected.

\bibitem{DingMuZhu2026}
M.-J. Ding, L. Mu and B.-X. Zhu,
\emph{Stieltjes moment property of the generalized Eulerian polynomials},
Advances in Applied Mathematics \textbf{177} (2026), 103077.
\href{https://doi.org/10.1016/j.aam.2026.103077}{doi:10.1016/j.aam.2026.103077}.
Publisher preview inspected; the full general Theorem~2.1 was unavailable
and remains a priority-check item.

% Insert these items inside the parent manuscript's thebibliography.
\bibitem{stab:CK}
S.~Chen and M.~Kauers,
Some open problems related to creative telescoping,
\emph{Journal of Systems Science and Complexity} \textbf{30} (2017),
154--172.
\href{https://doi.org/10.1007/s11424-017-6202-9}{doi:10.1007/s11424-017-6202-9}.
Preprint: \href{https://arxiv.org/abs/1609.03768}{arXiv:1609.03768} (2016).

\bibitem{stab:Chen}
S.~Chen,
Stability problems in symbolic integration,
\emph{Proceedings of ISSAC 2022} (2022), 517--524.
\href{https://doi.org/10.1145/3476446.3535502}{doi:10.1145/3476446.3535502}.
\href{https://arxiv.org/abs/2202.06305}{arXiv:2202.06305}.

\bibitem{stab:SSS}
C.~Sahu, M.~Seiss and V.~R.~Srinivasan,
Iterated and generalized iterated integrals,
\href{https://arxiv.org/abs/2504.04846}{arXiv:2504.04846} (2025).

\bibitem{stab:CDGHLL}
S.~Chen, H.~Du, Y.~Gao, H.~Huang, W.~Li and Z.~Li,
Complete reduction for derivatives in a transcendental Liouvillian extension,
\href{https://arxiv.org/abs/2602.03592}{arXiv:2602.03592v3},
4 August 2026.

\bibitem{stab:Rosenlicht}
M.~Rosenlicht,
Liouville's theorem on functions with elementary integrals,
\emph{Pacific Journal of Mathematics} \textbf{24} (1968), 153--161.
\href{https://doi.org/10.2140/pjm.1968.24.153}{doi:10.2140/pjm.1968.24.153}.

\bibitem{stab:Risch}
R.~H.~Risch,
The problem of integration in finite terms,
\emph{Transactions of the American Mathematical Society}
\textbf{139} (1969), 167--189.
\href{https://doi.org/10.1090/S0002-9947-1969-0237477-8}{doi:10.1090/S0002-9947-1969-0237477-8}.

\bibitem{stab:SS2026}
C.~Sahu and V.~R.~Srinivasan,
Elementary first integrals of integral differential systems,
\href{https://arxiv.org/abs/2608.09664}{arXiv:2608.09664},
10 August 2026.

% Add only these two new entries to the existing bibliography.
% The insertion reuses the existing Chen 2022 entry under key stab:Chen:
% S. Chen, Stability problems in symbolic integration, ISSAC 2022, 517--524,
% doi:10.1145/3476446.3535502; arXiv:2202.06305.  Do not duplicate it.

\bibitem{stab:Rosenlicht1969}
M.~Rosenlicht,
On the explicit solvability of certain transcendental equations,
\emph{Publications Math\'ematiques de l'IH\'ES} \textbf{36}
(1969), 15--22.
\href{https://www.numdam.org/item/PMIHES_1969__36__15_0/}{Primary source}.

\bibitem{stab:BBY}
J.~P.~Bell, S.~N.~Burris and K.~Yeats,
On the set of zero coefficients of a function satisfying a linear
differential equation,
\emph{Mathematical Proceedings of the Cambridge Philosophical Society}
\textbf{153} (2012), 235--247.
\href{https://doi.org/10.1017/S0305004112000114}{doi:10.1017/S0305004112000114}.
\href{https://arxiv.org/abs/1105.6078}{arXiv:1105.6078}.

\bibitem{Kang2026}
M. Kang,
\emph{A failure of the DeTemple complete-monotonicity threshold},
AI-assisted preprint, submitted 28 October 2026.
\href{https://mathdb.com/abs/2610.00788v1}{MathDB 2610.00788v1}.
The two-page mathematical argument was inspected; the accompanying
Lean certificate was not independently rebuilt in this project.


\bibitem{PrevostRivoal2021}
M. Pr\'evost and T. Rivoal,
\emph{Application of Pad\'e approximation to Euler's constant and
Stirling's formula}, Ramanujan Journal \textbf{54} (2021), 177--195.
\href{https://doi.org/10.1007/s11139-019-00201-9}{doi:10.1007/s11139-019-00201-9}.
Only the abstract and first-page preview were inspected.

\bibitem{LiChen2017}
X. Li and C.-P. Chen,
\emph{Pad\'e approximant related to asymptotics for the gamma function},
Journal of Inequalities and Applications \textbf{2017} (2017), 53.
\href{https://doi.org/10.1186/s13660-017-1315-1}{doi:10.1186/s13660-017-1315-1}.


\bibitem{Cauer1932}
W. Cauer,
\emph{The Poisson integral for functions with positive real part},
Bulletin of the American Mathematical Society \textbf{38} (1932), 713--717.
\href{https://doi.org/10.1090/S0002-9904-1932-05510-0}{doi:10.1090/S0002-9904-1932-05510-0}.
The representation is also stated in \cite[Eq.~(71)]{Kalugin2011}.

\bibitem{Kalugin2011}
G. A. Kalugin, D. J. Jeffrey and R. M. Corless,
\emph{Stieltjes, Poisson and other integral representations for functions
of Lambert W}, arXiv preprint (2011).
\href{https://arxiv.org/abs/1103.5640}{arXiv:1103.5640}.
Equation (71) refers to this preprint version, not the shortened final
journal version.

\bibitem{Luke1970}
Y. L. Luke,
\emph{Evaluation of the Gamma Function by Means of Pad\'e Approximations},
SIAM Journal on Mathematical Analysis \textbf{1} (1970), 266--281.
\href{https://doi.org/10.1137/0501024}{doi:10.1137/0501024}.
Full text not inspected; cited as an unresolved historical-overlap lead.

\bibitem{NemethZimanyi1982}
G. N\'emeth and M. Zim\'anyi,
\emph{Polynomial type Pad\'e approximants},
Mathematics of Computation \textbf{38} (1982), 553--565.
\href{https://doi.org/10.1090/S0025-5718-1982-0645671-7}{doi:10.1090/S0025-5718-1982-0645671-7}.
Full text not inspected; its treatment of scaled gamma approximants is
described in the inspected pages of \cite{McCabe1989}.

\bibitem{McCabe1989}
J. H. McCabe,
\emph{On continued fractions associated with polynomial type Pad\'e
approximants, with an application},
in L. Jacobsen (ed.), \emph{Analytic Theory of Continued Fractions III},
Lecture Notes in Mathematics \textbf{1406}, Springer (1989),
67--75.
\href{https://doi.org/10.1007/BFb0096166}{doi:10.1007/BFb0096166}.
Only the publisher's two-page preview, pp.~67--68, was inspected.


\bibitem{BhatnagarRajkumar2023}
G. Bhatnagar and K. Rajkumar,
\emph{Telescoping continued fractions for the error term in Stirling's
formula}, Journal of Approximation Theory \textbf{293} (2023), 105943.
\href{https://doi.org/10.1016/j.jat.2023.105943}{doi:10.1016/j.jat.2023.105943}.
\href{https://arxiv.org/abs/2204.00962}{arXiv:2204.00962v2}.

\bibitem{BhatnagarRajkumar2026}
G. Bhatnagar and K. Rajkumar,
\emph{The method of telescoping continued fractions},
preprint, 7 September 2026.
\href{https://arxiv.org/abs/2609.07092}{arXiv:2609.07092v1}.

\bibitem{wow:West}
D. B. West,
\emph{Some Conjectures of Graffiti.pc (2004--07)}, REGS 2009,
entry for Conjectures 174, 177, and 179 (2005).
\url{https://dwest.web.illinois.edu/regs/graffiti.html}.

\bibitem{wow:DW}
E. DeLaVi\~na and B. Waller,
\emph{Spanning trees with many leaves and average distance},
Electronic Journal of Combinatorics \textbf{15} (2008), R33.
\href{https://doi.org/10.37236/757}{doi:10.37236/757}.

\bibitem{Ducoffe2022}
G. Ducoffe,
\emph{Obstructions to Faster Diameter Computation: Asteroidal Sets},
17th International Symposium on Parameterized and Exact Computation
(IPEC 2022), LIPIcs 249, article 10.
\href{https://doi.org/10.4230/LIPIcs.IPEC.2022.10}{doi:10.4230/LIPIcs.IPEC.2022.10}.

\bibitem{GGPS2025}
R. Garrappa, S. Gerhold, M. Popolizio and T. Simon,
\emph{On some inequalities for the two-parameter Mittag--Leffler
function in the complex plane}, Journal of Mathematical Analysis and
Applications \textbf{551} (2025), 129588.
\href{https://doi.org/10.1016/j.jmaa.2025.129588}{doi:10.1016/j.jmaa.2025.129588}.
\href{https://arxiv.org/abs/2410.11852}{arXiv:2410.11852v4}.

\bibitem{DLMFMP}
NIST Digital Library of Mathematical Functions,
Meixner--Pollaczek generating function, Eq.~18.23.7.
\url{https://dlmf.nist.gov/18.23.E7}.

\bibitem{IndefJacobi}
M. Derevyagin and V. Derkach,
\emph{Spectral problems for generalized Jacobi matrices},
Linear Algebra and its Applications \textbf{382} (2004), 1--24.
\href{https://doi.org/10.1016/j.laa.2003.11.003}{doi:10.1016/j.laa.2003.11.003}.

\end{thebibliography}
\end{document}

% The sections below are optional standalone Python-FLINT verifiers.

%% BEGIN mittag_witness.py
"""Finite witness for Garrappa--Gerhold--Popolizio--Simon Conjecture 25.

Requires python-flint (Arb/Acb); every enclosure comparison is strict.
The analytic proof in the companion TeX file is independent of this calculation.
"""
from flint import acb, arb, ctx

ctx.prec = 320
alpha = arb(3) / 2
beta = arb(501) / 1000
cutoff = 300


def enclose_modulus(z):
    power = acb(1)
    partial = acb(0)
    for k in range(cutoff):
        partial += power * (alpha * k + beta).rgamma()
        power *= z
    r = abs(z)
    g = alpha * cutoff + beta
    two_step_ratio = r**2 / (g * (g + 1) * (g + 2))
    assert two_step_ratio < arb(1) / 500
    first_omitted = r**cutoff * g.rgamma()
    second_omitted = r**(cutoff + 1) * (g + alpha).rgamma()
    remainder = (first_omitted + second_omitted) * 500 / 499
    assert remainder < arb(1) / 10**210
    return partial, remainder


real_value, real_error = enclose_modulus(acb(-400))
complex_value, complex_error = enclose_modulus(acb(-400, 100))
assert real_value.imag.is_zero()
assert real_value.real - real_error > arb(410213) / 10**11
assert real_value.real + real_error < arb(410214) / 10**11
assert abs(complex_value) - complex_error > arb(390048) / 10**11
assert abs(complex_value) + complex_error < arb(390049) / 10**11
print('E(-400) strictly between 410213/10^11 and 410214/10^11')
print('|E(-400+100i)| strictly between 390048/10^11 and 390049/10^11')
print('All strict interval checks passed at', ctx.prec, 'bits.')
%% END mittag_witness.py

%% BEGIN continued_fraction_signs.py
"""Certify the signs through index 598 by Arb ball LDL, independent of sign values.

Run with python-flint. Positive/negative decisions require an interval strictly
on the selected side of zero. If any interval contains zero the run stops.
The first three basis polynomials form the rigorous block from the proof.
"""
import json
from flint import arb, ctx

ctx.prec=192
order=600
h=[arb(0)]; h2=[arb(0)]
for k in range(order//2+1):
    reciprocal=arb(1)/(2*k+1)
    h.append(h[-1]+reciprocal)
    h2.append(h2[-1]+reciprocal**2)

def gram(i,j):
    """Coefficient of s^i z^j in 4*(atan(s)+atan(z))^2/(1-s*z)."""
    if i>j:return gram(j,i)
    if (i-j)%2:return arb(0)
    if i==j:return 8*h2[(i+1)//2]
    d=(j-i)//2
    return 4*(-1)**(d+1)*(h[(j+1)//2]-h[(i+1)//2])/d

def schur(i,j):
    """Complement of leading block [[0,0,4],[0,8,0],[4,0,8]]."""
    return (gram(i,j)+gram(i,0)*gram(0,j)/2
            -(gram(i,0)*gram(2,j)+gram(i,2)*gram(0,j))/4
            -gram(i,1)*gram(1,j)/8)

sign_at_degree={}
for parity in (0,1):
    degrees=[i for i in range(3,order) if i%2==parity]
    a=[[schur(i,j) for j in degrees[:k+1]] for k,i in enumerate(degrees)]
    for k,degree in enumerate(degrees):
        pivot=a[k][k]
        if pivot>0:sign_at_degree[degree]=1
        elif pivot<0:sign_at_degree[degree]=-1
        else:raise ArithmeticError(f'precision insufficient at degree {degree}: {pivot}')
        # Lower triangular rank-one update of the remaining Schur complement.
        for i in range(k+1,len(degrees)):
            ratio=a[i][k]/pivot
            for j in range(k+1,i+1):a[i][j]-=ratio*a[j][k]

# The size-three first block has S_1=-1. Each remaining nonzero Schur pivot
# is a unit normal-index increment; the proof identifies its sign with S_j.
signs=[-1]+[sign_at_degree[i] for i in range(3,order)]
negative=[j for j in range(2,len(signs)+1) if signs[j-1]!=signs[j-2]]
report={'precision_bits':ctx.prec,'gram_order':order,'certified_through_coefficient':len(signs),
        'negative_coefficients':negative,'always_nonzero_pivots_after_initial_degree_three':True,
        'initial_coefficient_positive':True,'method':'outward-rounded Arb triangular elimination on parity blocks'}
print(json.dumps(report,indent=2))
%% END continued_fraction_signs.py
